\documentclass[12pt, a4paper, openany, oneside]{book}
\usepackage[utf8]{inputenc}
\usepackage[T1]{fontenc}
\usepackage[brazilian]{babel}
\usepackage{libertine}

\usepackage{geometry}
\geometry{
    a4paper,            % Redundante com a classe, mas boa prática manter
    left=3cm,           % Margem esquerda
    right=3cm,          % Margem direita (Garante a simetria com a esquerda)
    top=3cm,            % Margem superior
    bottom=3cm,         % Margem inferior
    marginparwidth=2cm, % Define exatamente 2cm para as notas do comando \fonte
    marginparsep=4mm    % Define um respiro de 4mm entre o texto principal e a nota
}

\usepackage{mathtools}
\usepackage{amsthm}
\usepackage{amsmath}
\usepackage{amssymb}
\usepackage{amsfonts}
\usepackage{etoolbox}

\usepackage{csquotes}
\usepackage{microtype}
\usepackage{enumitem}
\usepackage{esvect}
\usepackage{verbatim}
%\usepackage{cancel}
\usepackage{xcolor}
\usepackage{marginnote}

%%%%%%%%%%%%%%%%%%%%%%%%%%%%%%

\setlength{\parskip}{8pt}
\setlength{\parindent}{0pt}
\setlist[itemize]{itemsep=0pt, topsep=0pt}
\setlist[description]{itemsep=0pt, topsep=0pt}
\setlist[enumerate]{itemsep=0pt, topsep=0pt}

\newlist{iffe}{description}{1}
\setlist[iffe]{
  font=\normalfont,
  labelwidth=2.6em,
  labelsep=0.5em,
  leftmargin=3.1em,
  align=right,
  nosep
}

\allowdisplaybreaks
\raggedbottom
%\setlength\emergencystretch{0.5em}
%\tolerance=300

\setcounter{tocdepth}{3}
\setcounter{secnumdepth}{3}

%%%%%%%%%%%%%%%%%%%%%%%%%%%%%%

\usepackage[backend=biber]{biblatex}
\addbibresource{bibliografia.bib}

%%%%%%%%%%%%%%%%%%%%%%%%%%%%%%

% Inverte a posição padrão das notas de margem (joga para a esquerda)
\reversemarginpar

% Cria a chave (toggle) para controlar a exibição das notas internas
\newtoggle{mostrarfontes}

% Escolha UM dos dois abaixo comentando o outro:
\toggletrue{mostrarfontes}  % MODO DE REVISÃO: Mostra as fontes na margem
%\togglefalse{mostrarfontes} % MODO FINAL: Oculta todas as fontes, PDF limpo

% Comando para inserir a fonte na margem quando o modo de revisão estiver ativo
\newcommand{\fonte}[2][]{%
  \iftoggle{mostrarfontes}{%
    \marginnote{\tiny\sffamily\color{blue}Fonte: \cite[#1]{#2}}%
  }{}%
}

%%%%%%%%%%%%%%%%%%%%%%%%%%%%%%

\newcommand{\F}{\mathbb{F}}
\newcommand{\C}{\mathbb{C}}
\newcommand{\R}{\mathbb{R}}
\newcommand{\Q}{\mathbb{Q}}
\newcommand{\Z}{\mathbb{Z}}
\newcommand{\N}{\mathbb{N}}
\newcommand{\Pb}{\mathbb{P}}
\newcommand{\ds}{\displaystyle}
\newcommand{\op}{\operatorname}
\newcommand{\id}{\operatorname{Id}}
\renewcommand{\Im}{\operatorname{Im}}
\DeclareMathOperator{\ind}{Ind}
\DeclareMathOperator{\Img}{Im}
\DeclareMathOperator{\sen}{sen}
\DeclareMathOperator{\Dom}{Dom}
\DeclareMathOperator{\Int}{int} %interior

\newcommand{\itemproof}{\hfill$\blacksquare$\par}
\newcommand{\blackproof}{\hfill$\blacksquare$\par}
\newcommand{\whiteproof}{\hfill$\square$\par}

\let\oldrestriction\restriction
\DeclareRobustCommand{\restriction}{\mathord{\oldrestriction}}

\makeatletter
\NewDocumentCommand{\lowerint}{}{\mathop{}\mathpalette\lowerint@\relax\!\int}
\NewDocumentCommand{\upperint}{}{\mathop{}\mathpalette\upperint@\relax\!\int}
\newcommand{\lowerint@}[2]{%
  \begingroup
  \sbox\z@{$\m@th#1\int$}%
  \lowup@{l}{\underline}{#1}%
  \endgroup
}
\newcommand{\upperint@}[2]{%
  \begingroup
  \sbox\z@{$\m@th#1\int$}%
  \lowup@{r}{\overline}{#1}%
  \endgroup
}
\newcommand{\lowup@}[3]{%
  % #1 = l (lower) or r (upper)
  % #2 = \underline (lower) or \overline (upper)
  % #3 = math style
  \rlap{%
    \hspace{0.05\wd\z@}%
    \makebox[0.9\wd\z@][#1]{%
      $\m@th#2{%
        \hspace{0.4\wd\z@}%
        \ifx#3\displaystyle\else\hspace{0.1\wd\z@}\fi
        \vphantom{\copy\z@}%
      }$%
    }%
    \hspace{0.05\wd\z@}%
  }%
}
\makeatother

\makeatletter
\newcommand{\repeatitem}{\addtocounter{\@enumctr}{-1}\item}
\makeatother

%%%%%%%%%%%%%%%%%%%%%%%%%%%%%%

\usepackage[
    colorlinks=true,
    linkcolor=blue,
    citecolor=blue,
    urlcolor=blue,
    breaklinks=true
]{hyperref}
\usepackage[noabbrev]{cleveref}

\newtheoremstyle{meuestilo}
  {8pt}                     % Espaço acima
  {0pt}                     % Espaço abaixo
  {\normalfont}             % Fonte do corpo do teorema
  {}                        % Indentação (vazio = nenhuma)
  {\bfseries}               % Fonte do cabeçalho do teorema
  {.}                       % Pontuação após o cabeçalho
  { }                       % Espaço após o cabeçalho
  {}                        % Especificação do cabeçalho (vazio = `Teorema 1.1`)

\theoremstyle{meuestilo}

\newtheorem{definition}{Definição}[section]
\newtheorem{theorem}{Teorema}[section]
\newtheorem{proposition}{Proposição}[section]
\newtheorem{lemma}{Lema}[section]
\newtheorem{corollary}{Corolário}[section]
\newtheorem{remark}{Observação}[section]
\newtheorem{example}{Exemplo}[section]
\newtheorem{axiom}{Axioma}[section]
\newtheorem{notation}{Notação}[section]
\newtheorem{claim}{Afirmação}

\crefname{definition}{definição}{definições}
\crefname{theorem}{teorema}{teoremas}
\crefname{proposition}{proposição}{proposições}
\crefname{lemma}{lema}{lemas}
\crefname{corollary}{corolário}{corolários}
\crefname{remark}{observação}{observações}
\crefname{example}{exemplo}{exemplos}
\crefname{axiom}{axioma}{axiomas}
\crefname{notation}{notação}{notações}
\crefname{claim}{afirmação}{afirmações}

\crefname{equation}{equação}{equações}
\crefname{chapter}{capítulo}{capítulos}
\crefname{section}{seção}{seções}
\crefname{subsection}{subseção}{subseções}
\crefname{appendix}{apêndice}{apêndices}
\crefname{figure}{figura}{figuras}
\crefname{table}{tabela}{tabelas}

\newcommand{\crefpairconjunction}{ e }
\newcommand{\crefmiddleconjunction}{, }
\newcommand{\creflastconjunction}{ e }

\AtBeginEnvironment{proof}{\setcounter{claim}{0}} % 2. Forçamos o contador a voltar para 0 sempre que um 'proof' começar
\renewenvironment{proof}[1][Prova.]{\par\noindent\textbf{#1} }{}
\newenvironment{midproof}[1][Prova.]{\par\noindent\emph{#1} }{}

%%%%%%%%%%%%%%%%%%%%%%%%%%%%%%

\title{Fundamentos}
\author{Alexander Kahleul}
\date{\today}

%%%%%%%%%%%%%%%%%%%%%%%%%%%%%%

\begin{document}
\maketitle
\setcounter{page}{2}
\tableofcontents

%%%%%%%%%%%%%%%%%%%%%%%%%%%%%%

\part{Teoria dos Conjuntos}
%!TEX root = main.tex

\chapter{Os Axiomas}

\section{Primeiros Axiomas} \label{sec1}

\subsection{O Axioma da Extensão}

\begin{axiom}[da Extensão] \label{ax:1}
    Se dois conjuntos têm os mesmos elementos, então eles são iguais.
        \[
            \boxed{
                \forall x \forall y (\forall z((z \in x) \leftrightarrow (z \in y)) \rightarrow (x=y))
            }
        \]
\end{axiom}

\begin{remark}
    Com a definição de $\subseteq$, o axioma da extensão pode ser enunciado como
        \[
            \forall x \forall y (((x \subseteq y) \land (y \subseteq x)) \rightarrow (x=y)).
        \] 
\end{remark}

\begin{corollary}
    \leavevmode
        \begin{enumerate}[leftmargin=*, align=left, label=\textbf{(\alph*)}]
            \item $\forall x (x \subseteq x)$.
            \item $\forall x \forall y ((x \subseteq y) \land (y \subseteq x) \rightarrow (x=y))$.
            \item $\forall x \forall y \forall z (((x \subseteq y) \land (y \subseteq z)) \rightarrow (x \subseteq z))$.
        \end{enumerate}
\end{corollary}

\begin{proof}
    \blackproof
\end{proof}

\subsection{O Axioma do Vazio}

\begin{definition}
    Um conjunto $x$ é \emph{vazio} se $\forall y (y \notin x)$.
\end{definition}

\begin{axiom}[do Vazio] \label{ax:2}
    Existe um conjunto vazio.
        \[
            \boxed{
                \exists x \forall y (y \notin x)
            }
        \]
\end{axiom}

\begin{proposition} \label{prop.set:extensaounico1}
    Quaisquer dois conjuntos vazios são iguais.
        \[
            \forall x_1 \forall x_2 ( (\forall y (y \notin x_1) \land \forall y (y \notin x_2)) \rightarrow (x_1 = x_2) ).
        \]
\end{proposition}

\begin{proof}
Se $x_1 \neq x_2$, então ou existe $z \in x_1$ tal que $z \notin x_2$, ou existe $z \in x_2$ tal que $z \notin x_1$. No primeiro caso, $x_1$ não é vazio, e, no segundo caso, $x_2$ não é vazio. Em ambos os casos, temos uma contradição. Logo $x_1 = x_2$. \blackproof
\end{proof}

\begin{remark}
    O axioma do vazio, junto com a \cref{prop.set:extensaounico1}, nos permite estabelecer que existe um único conjunto vazio:
        \[
            \exists x (\forall y  (y \notin x) \land \forall z(\forall y ( y \notin z ) \rightarrow (z=x) )).
        \]
    Podemos então falar \emph{do} conjunto vazio (em vez de \emph{de um}). Ele é denotado por $\emptyset$.
\end{remark}

\begin{proposition}
    O conjunto vazio está contido em qualquer conjunto:
        \[
          \forall x (\emptyset \subseteq x).
        \]
\end{proposition}

\begin{proof}
    Se existisse $x$ tal que $\emptyset \not\subseteq x$, então existiria $y \in \emptyset$ tal que $y \notin x$, uma contradição pois $\forall y (y \notin \emptyset)$. Logo, $\forall x (\emptyset \subseteq x)$. \blackproof
\end{proof}

\subsection{O Axioma do Par}

\begin{axiom}[do Par] \label{ax:3}
    Para quaisquer conjuntos $x$ e $y$, existe um conjunto cujos elementos são $x$ e $y$.
        \[
            \boxed{
                \forall x \forall y \exists z \forall w ((w \in z) \leftrightarrow ((w = x) \lor (w = y)))
            }
        \]
\end{axiom}

\begin{proposition} \label{prop.set:extensaounico2}
    O conjunto $z$ do axioma do par é único. Ele é denotado por $\{x,y\}$.
\end{proposition}

\begin{proof}
    Se
        $
            \forall w ( (w \in z') \leftrightarrow ((w=x) \lor (w=y)) ),
        $
    então, como
        $
            \forall w ((w \in z) \leftrightarrow ((w = x) \lor (w = y))),
        $
    temos $\forall w ((w \in z') \leftrightarrow (w \in z) )$, de modo que, pelo axioma da extensão, $z' = z$. \blackproof
\end{proof}

\begin{proposition}
    $\forall x \forall y((x \in y) \leftrightarrow \{x\} \subseteq y)$.
\end{proposition}

\begin{proof}
    Por um lado ($\Rightarrow$), suponha que $x \in y$. Se $z \in \{x\}$, então $z = x$. Como, por hipótese, $x \in y$, de $z=x$ vem $z \in y$. Logo $\{x\} \subseteq y$. Por outro lado ($\Leftarrow$), suponha que $\{x\} \subseteq y$. Temos $x \in \{x\}$, e como $\{x\} \subseteq y$, temos $x \in y$. \blackproof
\end{proof}

\subsection{O Axioma da União}

\begin{axiom}[da União] \label{ax:4}
    Para todo conjunto $x$ existe o conjunto de todos os conjuntos que pertencem a algum elemento de $x$.
        \[
            \boxed{
                \forall x \exists y \forall z ((z \in y) \leftrightarrow \exists w ( (z \in w) \land (w \in x)))
            }
        \]
\end{axiom}

\begin{proposition} \label{prop.set:extensaounico3}
    O conjunto $y$ do axioma da união é único. Ele é denotado por $\bigcup x $.
\end{proposition}

\begin{proof}
    Se
        $
            \forall z ((z \in y') \leftrightarrow \exists w ( (z \in w) \land (w \in x))),
        $
    então, como
        $
            \forall z ((z \in y) \leftrightarrow \exists w ( (z \in w) \land (w \in x))),
        $
    temos $\forall z ((z \in y') \leftrightarrow (z \in y)) $, de modo que, pelo axioma da extensão, $y' = y$. \blackproof
\end{proof}

\begin{theorem}
    Para quaisquer conjuntos $x$ e $y$, existe o conjunto dos conjuntos que pertencem a $x$ ou a $y$.
        \[
            \forall x \forall y \exists z \forall w ((w \in z) \leftrightarrow ((w \in x) \lor (w \in y)))
        \]
    Ademais, esse conjunto é único, sendo denotado por $x \cup y$.
\end{theorem}

\begin{proof}
    Provemos que $z := \bigcup \{x,y\}$, que existe pelos axiomas do par e da união, funciona. De fato, para todo $w$, temos $w \in z$ se, e somente se, existe $u \in \{x,y \}$ tal que $w \in u$. Como $u \in \{ x,y \}$ se, e somente se, $u = x$ ou $u=y$, temos $w \in x$ ou $w \in y$. A unicidade de $z$ segue do axioma da extensão, de modo que podemos denotar $x \cup y := z$. \blackproof
\end{proof}

\subsection{O Axioma das Partes}

\begin{axiom}[das Partes] \label{ax:5}
    Para todo conjunto $x$, existe o conjunto dos subconjuntos de $x$.
        \[
            \boxed{
                \forall x \exists y \forall z ( (z \in y) \leftrightarrow (z \subseteq x) )
            }
        \]
\end{axiom}

\begin{proposition} \label{prop.set:extensaounico4}
    O conjunto $y$ do axioma das partes é único. Ele é denotado por $\mathcal{P}(x)$.
\end{proposition}

\begin{proof}
    Pelo axioma das partes, temos $\forall z ( (z \in y) \leftrightarrow (z \subseteq x) )$. Seja $y'$ tal que $\forall z ( (z \in y') \leftrightarrow (z \subseteq x) )$. Logo, $\forall z ((z \in y' )\leftrightarrow(z \in y))$, de modo que, pelo axioma da extensão, $y' = y$. \blackproof
\end{proof}

\subsection{O Esquema de Axiomas da Separação}

\begin{axiom}[da Separação] \label{ax:6}
    Para toda fórmula $P$ em que $z$ não ocorre livre, a fórmula
        \[
            \boxed{
                \forall y \exists z \forall x ( (x \in z) \leftrightarrow ((x \in y) \land P) )
            }    
        \]
    é um axioma.
\end{axiom}

\begin{remark}
    O axioma da separação também é chamado de axioma da compreensão ou axioma da especificação. O conjunto $y$ é chamado de ``universo'' da discussão.
\end{remark}

\begin{proposition}
    O conjunto $z$ do axioma da separação é único. Ele é denotado por $\{x \in y : P(x) \}$.
\end{proposition}

\begin{proof}
    Segue do axioma da extensão. \blackproof
\end{proof}

\begin{theorem}[Paradoxo de Russell] Não existe o conjunto de todos os conjuntos.
        \[
            \forall x \exists y (y \notin x)
        \]
\end{theorem}

\begin{proof}
    Suponha, por absurdo, que $\exists x \forall y (y \in x)$. Pelo axioma da separação em $x$ com a fórmula $y \notin y$, existe $z$ tal que $\forall y ( (y \in z) \leftrightarrow ((y \in x) \land (y \notin y)))$ (note que $z$ não ocorre livre em $y \notin y$). Como $\forall y (y \in x)$, temos $\forall y ((y \in z) \leftrightarrow (y \notin y))$. Particularmente para $y = z$, temos $((z \in z) \leftrightarrow (z \notin z))$, uma contradição. Logo $\forall x \exists y (y \notin x)$. \blackproof
\end{proof}

\begin{theorem} \label{teo.set:intfamilia}
    Para todo conjunto $x \neq \emptyset$ existe o conjunto de todos os conjuntos que pertencem simultaneamente a todos os elementos de $x$.
        \[
            \forall x ((x \neq \emptyset) \rightarrow \exists y \forall z ((z \in y) \leftrightarrow \forall w ( (w \in x) \rightarrow (z \in w))))
        \]
    Ademais, esse conjunto é único, sendo denotado por $\bigcap x$.
\end{theorem}

\begin{proof}
    Precisamos provar que existe $y$ tal que
        \[
           \forall z ( (z \in y) \leftrightarrow \forall w ( (w \in x) \rightarrow (z \in w)) ). 
        \]
    Como $x \neq \emptyset$, tome $v \in x$. Pelo axioma da separação com universo $v$ e a fórmula $\forall w ( (w \in x) \rightarrow (z \in w))$, onde $y$ não ocorre livre, existe $y$ tal que
        \[
            \forall z ( (z \in y) \leftrightarrow ( (z \in v) \land  (\forall w ( (w \in x) \rightarrow (z \in w))) ) ),
        \]
    isto é, existe 
        \[
            y := \{ z \in v : \forall w ( (w \in x) \rightarrow (z \in w)) \}.
        \]
    Afirmamos que esse $y$ funciona. De fato, por um lado ($\Rightarrow$), se $z \in y$, então, trivialmente, $\forall w ((w \in x) \rightarrow (z \in w))$. Por outro lado ($\Leftarrow$), se $z$ é tal que $\forall w ((w \in x) \rightarrow (z \in w))$, então, particularmente para $w = v$, temos $v \in x \rightarrow z \in v$, e como $v \in x$, temos $z \in v$. Como $z \in v$ e $\forall w ((w \in x) \rightarrow (z \in w))$, temos que $z \in y$. Isso completa a prova da existência. A unicidade de $y$ segue do axioma da extensão, de modo que podemos denotar $\bigcap x := y$. \blackproof
\end{proof}

\begin{definition} \label{defi.set:capminus}
    Sejam $x$ e $y$ conjuntos.
        \begin{enumerate}[leftmargin=*, align=left, label=\textbf{(\alph*)}]
            \item A \emph{interseção} entre $x$ e $y$ é definida como
                \[
                    x \cap y :=  \{ z \in x : z \in y\}.
                \]
            Dizemos que $x$ e $y$ são \emph{disjuntos} se $x \cap y = \emptyset$. 
            \item A \emph{diferença} entre $x$ e $y$ é definida como
                \[
                    x \setminus y := \{z \in x : z \notin y \}.
                \]
            Dizemos que $x \setminus y$ é o \emph{complementar} de $y$ relativo a $x$ se $y \subseteq x$. Isso é denotado por $y^{C} := x \setminus y$.
            \item A \emph{diferença simétrica} entre $x$ e $y$ é definida como
                \[
                    x \Delta y := \{ z \in x \cup y :z \notin x \cap y \}.
                \]
        \end{enumerate}
\end{definition}

\begin{remark}
    A boa definição da \cref{defi.set:capminus} segue do \cref{ax:6} da separação. Vejamos como isso é feito, por exemplo, na definição de $x \cap y$. Sendo $x$ o universo, o axioma da separação é a fórmula $\forall x \exists w \forall z ((z \in w) \leftrightarrow ((z \in x) \land P) )$, onde $P$ é uma fórmula em que $w$ não ocorre livre. Se $P$ é a fórmula $z \in y$, então existe $w$ que cumpre $\forall z ( (z \in w) \leftrightarrow (z \in x) \land (z \in y) )$, isto é, $w = \{z \in x : z \in y \}$. Denotamos esse $w$ por $x \cap y$.
\end{remark}

\subsection{Propriedades Algébricas}

Apresentamos a seguir algumas propriedades algébricas de algumas operações com conjuntos, como união e interseção. Por serem simples e de caráter operacional, omitimos as suas demonstrações.

\begin{proposition}[Propriedades da União]
    \leavevmode
        \begin{enumerate}[leftmargin=*, align=left, label=\textbf{(\alph*)}]
            \item $\forall x (x \cup x = x)$.
            \item $\forall x (x \cup \emptyset = x)$.
            \item $\forall x \forall y (x \cup y = y \cup x)$.
            \item $\forall x \forall y \forall z (x \cup (y \cup z) = (x \cup y) \cup z)$.
            \item $\forall x \forall y (x \cup y = y \leftrightarrow x \subseteq y)$.
            \item $\forall x \forall y ( (x \subseteq x \cup y) \land (y \subseteq x \cup y))$.
            \item $\forall x \forall y \forall z (x \subseteq y \rightarrow x \cup z \subseteq y \cup z )$.
            \item $\forall x \forall y (x \subseteq y \rightarrow \bigcup x \subseteq \bigcup y)$.
        \end{enumerate}
\end{proposition}

\begin{proposition}[Propriedades da Interseção]
    \leavevmode
        \begin{enumerate}[leftmargin=*, align=left, label=\textbf{(\alph*)}]
            \item $\forall x (x \cap x = x)$.
            \item $\forall x (x \cap \emptyset = \emptyset)$.
            \item $\forall x \forall y (x \cap y = y \cap x)$.
            \item $\forall x \forall y \forall z (x \cap (y \cap z) = (x \cap y) \cap z)$.
            \item $\forall x \forall y (x \cap y = x \leftrightarrow x \subseteq y)$.
            \item $\forall x \forall y ( (x \cap y \subseteq x) \land (x \cap y \subseteq y))$.
            \item $\forall x \forall y \forall z (x \subseteq y \rightarrow x \cap z \subseteq y \cap z)$.
            \item $\forall x \forall y (x \subseteq y \land x \neq \emptyset \rightarrow \bigcap y \subseteq \bigcap x)$.
        \end{enumerate}
\end{proposition}

\begin{proposition}[Distributividade]
    \leavevmode
        \begin{enumerate}[leftmargin=*, align=left, label=\textbf{(\alph*)}]
            \item $\forall x \forall y \forall z (x \cap (y \cup z) = (x \cap y) \cup (x \cap z))$.
            \item $\forall x \forall y \forall z (x \cup (y \cap z) = (x \cup y) \cap (x \cup z))$.
        \end{enumerate}
\end{proposition}

\begin{proposition}[Propriedades da Diferença]
    \leavevmode
        \begin{enumerate}[leftmargin=*, align=left, label=\textbf{(\alph*)}]
            \item (Imediatas).
                \begin{enumerate}[label=\roman*.]
                    \item $\forall x (x \setminus \emptyset = x)$.
                    \item $\forall x (x \setminus x = \emptyset)$.
                    \item $\forall x (\emptyset \setminus x = \emptyset)$.
                \end{enumerate}
            \item \leavevmode
                \begin{enumerate}[label=\roman*.]
                    \item $\forall x \forall y (x \setminus y = x \leftrightarrow x \cap y = \emptyset)$.
                    \item $\forall x \forall y (x \setminus y = \emptyset \leftrightarrow x \subseteq y)$.
                \end{enumerate}
            \item (Leis de De Morgan).
                \begin{enumerate}[label=\roman*.]
                    \item $\forall x \forall y \forall z ( x \setminus (y \cup z) = (x \setminus y) \cap (x \setminus z) )$.
                    \item $\forall x \forall y \forall z ( x \setminus (y \cap z) = (x \setminus y) \cup (x \setminus z) )$.
                \end{enumerate}
            \item $\forall x \forall y \forall z ( x \setminus (y \setminus z) = (x \setminus y) \cup (x \cap z) )$.
            \item (Diferenças entre interseções).
                \begin{enumerate}[label=\roman*.]
                    \item $\forall x \forall y \forall z ( x \cap (y \setminus z) = (x \cap y) \setminus (x \cap z) )$.
                    \item $\forall x \forall y \forall z ( (x \setminus y) \cap z = (x \cap z) \setminus y )$.
                \end{enumerate}
            \item (Monotonocidade da diferença).
                \begin{enumerate}[label=\roman*.]
                    \item $\forall x \forall y \forall z (x \subseteq y \rightarrow x \setminus z \subseteq y \setminus z)$.
                    \item $\forall x \forall y \forall z (y \subseteq z \rightarrow x \setminus z \subseteq x \setminus y)$.
                \end{enumerate}
        \end{enumerate}
\end{proposition}

\begin{proposition}[Propriedades do Complemento Relativo]
    \leavevmode
        \begin{enumerate}[leftmargin=*, align=left, label=\textbf{(\alph*)}]
            \item $\forall x \forall y (y \subseteq x \rightarrow (y^C)^C = y)$.
            \item $\forall x \forall y (y \subseteq x \rightarrow (y \cup y^C = x) \land (y \cap y^C = \emptyset))$.
            \item $\forall x \forall y \forall z (z \subseteq y \subseteq x \rightarrow y^C \subseteq z^C)$.
            \item $\forall x \forall y \forall z (y \subseteq x \land z \subseteq x \rightarrow y \setminus z = y \cap z^C)$.
            \item (Leis de De Morgan para complementos).
                \begin{enumerate}[label=\roman*.]
                    \item $\forall x \forall y \forall z (y \subseteq x \land z \subseteq x \rightarrow (y \cup z)^C = y^C \cap z^C)$.
                    \item $\forall x \forall y \forall z (y \subseteq x \land z \subseteq x \rightarrow (y \cap z)^C = y^C \cup z^C)$.
                \end{enumerate}
        \end{enumerate}
\end{proposition}

\begin{proposition}[Variadas]
    \leavevmode
        \begin{enumerate}[leftmargin=*, align=left, label=\textbf{(\alph*)}]
            \item $\forall x \forall y \forall z ( ((x \in y) \land (y \in z)) \rightarrow ( ( x \in \bigcup z) \land (y \subseteq \bigcup z)))$.
            \item $\forall A (\bigcup \mathcal{P}{(A)} = A)$.
        \end{enumerate}
\end{proposition}

\section{Relações e Funções}

\subsection{Produto Cartesiano}

\begin{definition}[Par ordenado]
    Sejam $a$ e $b$ conjuntos. O \emph{par ordenado} $(a,b)$ é definido como o conjunto $\{ \{ a \}, \{a,b \} \}$, isto é,
        \[
            (a,b) := \{ \{ a \}, \{a,b \} \}.
        \]
\end{definition}

\begin{proposition}
    Dois pares ordenados $(a,b)$ e $(c,d)$ são iguais se, e somente se, $a = c$ e $b=d$.
        \[
            \forall a \forall b \forall c \forall d ( ((a,b) = (c,d)) \leftrightarrow ((a=c) \land (b=d)))
        \]
\end{proposition}

\begin{proof}
    Ver \cite{fajardo2024conjuntos}, teorema 4.2, página 95. Ver \cite{jechhrbacek1999}, teorema 1.2, página 18. \blackproof
\end{proof}

\begin{theorem} \label{teo.set:prodcart}
    Para quaisquer conjuntos $A$ e $B$, existe um conjunto com todos os pares ordenados $(a,b)$ tais que $a \in A$ e $b \in B$.
        \[
            \forall A \forall B \exists C \forall x (x \in C \leftrightarrow \exists a \exists b (a \in A \land b \in B \land x = (a,b))) 
        \]
    Ademais, esse conjunto é único, sendo denotado por $A \times B$.
\end{theorem}

\begin{proof}
    Defina
        \[
            C := \{x \in \mathcal{P}{(\mathcal{P}{(A \cup B)})} : \exists a \exists b (a \in A \land b \in B \land x = (a,b)) \}.
        \]
    Afirmamos que $C$ funciona. \blackproof
\end{proof}

\begin{definition}
    O \emph{produto cartesiano} dos conjuntos $A$ e $B$ é definido como $A \times B$.
\end{definition}

\begin{proposition}
    \leavevmode
        \begin{enumerate}[leftmargin=*, align=left, label=\textbf{(\alph*)}]
            \item Para quaisquer $A$ e $B$, temos $A \times B = \emptyset$ se, e somente se, $A = \emptyset$ ou $B = \emptyset$.
            \item Para quaisquer $A$, $B$ e $C$,
                \begin{enumerate}[label=\roman*.]
                    \item $(A \cup B) \times C = A \times C \cup B \times C$.
                    \item $A \times (B \cup C) = A \times B \cup A \times C$.
                \end{enumerate}
            \item Para quaisquer $A$, $B$ e $C$,
                \begin{enumerate}[label=\roman*.]
                    \item $(A \cap B) \times C = A \times C \cap B \times C$.
                    \item $A \times (B \cap C) = A \times B \cap A \times C$.
                \end{enumerate}
            \item Para quaisquer $A$, $B$ e $C$,
                \begin{enumerate}[label=\roman*.]
                    \item $(A \setminus B) \times C = A \times C \setminus B \times C$.
                    \item $A \times (B \setminus C) = A \times B \setminus A \times C$.
                \end{enumerate}
        \end{enumerate}
\end{proposition}

\begin{proof}
    \blackproof
\end{proof}

\subsection{Relações}

\begin{definition}
    Uma \emph{relação binária}, ou, simplesmente, uma \emph{relação}, é um conjunto de pares ordenados.
        \begin{enumerate}[leftmargin=*, align=left, label=\textbf{(\alph*)}]
            \item Dizemos que $R$ é uma relação de $A$ em $B$ se $R \subseteq A \times B$. Dizemos que $R$ é uma relação em $A$ se $R \subseteq A \times A$. Denotamos $(a,b) \in R$ por $aRb$.
        \end{enumerate}
\end{definition}

\begin{theorem}
    $R$ é uma relação se, e somente se, existem $A$ e $B$ tais que $R \subseteq A \times B$.
\end{theorem}

\begin{proof}
    \leavevmode
        \begin{itemize}[align=left]
            \item[($\Rightarrow$)] Seja $R$ uma relação. Defina
                \begin{align*}
                    A &:= \left\{ a \in \bigcup \bigcup R : \exists b \left(b \in \bigcup \bigcup R \land aRb \right) \right\} \\
                    B &:= \left\{ b \in \bigcup \bigcup R : \exists a \left(a \in \bigcup \bigcup R \land aRb \right) \right\}
                \end{align*}
            Seja $x \in R$. Então existem $a$ e $b$ tais que $x = (a,b)$. Se $\{ \{a \}, \{a,b \} \} \in R$, então $\{ \{a \}, \{a,b \} \} \subseteq \bigcup R$, donde $\{a,b \} \in \bigcup R$, donde $\{ a,b \} \subseteq \bigcup \bigcup R$, donde $a,b \in \bigcup \bigcup R$. Como $aRb$, pelas definições de $A$ e $B$, temos $a \in A$ e $b \in B$. Como $x = (a,b)$ e $a \in A$ e $b \in B$, pelo \cref{teo.set:prodcart} temos $x \in A \times B$, donde, por fim, segue que $R \subseteq A \times B$.
            \item[($\Leftarrow$)] Os elementos de $A \times B$ são pares ordenados; logo, qualquer subconjunto de $A \times B$ terá pares ordenados como elementos.
        \end{itemize}
    \blackproof
\end{proof}

\begin{definition}
    Seja $R$ uma relação.
        \begin{enumerate}[leftmargin=*, align=left, label=\textbf{(\alph*)}]
            \item O \emph{domínio} de $R$ é definido como
                \[
                    \Dom R := \left\{ a \in \bigcup \bigcup R : \exists b ((a,b) \in R) \right\}.
                \]
            \item A \emph{imagem} de $R$ é definida como
                \[
                    \Img R := \left\{ b \in \bigcup \bigcup R : \exists a ((a,b) \in R) \right\}.
                \]
            \item A \emph{relação inversa} de $R$ é definida como
                \[
                    R^{-1} := \{ (b,a) \in \Img R \times \Dom R : (a,b) \in R \}.
                \]
            \item A \emph{imagem de um conjunto $X$ por $R$} é definida como
                \[
                    R[X] := \{ b \in \bigcup \bigcup R : \exists a (a \in X \land (a,b) \in R) \}.
                \]
            \item A \emph{imagem inversa de um conjunto $Y$ por $R$} é definida como
                \[
                    R^{-1}[Y] := \left\{ a \in \bigcup \bigcup R : \exists b (b \in Y \land (a,b) \in R) \right\}
                \]
            Equivalentemente, $R^{-1}[Y]$ é a imagem de $Y$ pela relação $R^{-1}$. 
            \item A \emph{restrição de $R$ a $X$} é definida como
                \[
                    R \restriction_X := \{ (a,b) \in R : a \in X \}.
                \]
            \item A \emph{composição de $R$ e $S$} é definida como 
                \[
                    S \circ R := \{ (a,c) \in \Dom R \times \Img S : \exists b (b \in \Img R \cap \Dom S \land (aRb \land bSc)) \}.
                \]
        \end{enumerate}
\end{definition}

\begin{proposition} \label{prop.set:domeimdeRAB}
    Sejam $R$ e $S$ relações e $A$, $B$, $C$ e $D$ tais que $R \subseteq A \times B$ e $S \subseteq C \times D$.
        \begin{enumerate}[leftmargin=*, align=left, label=\textbf{(\alph*)}]
            \item $\Dom R = R^{-1}[B] = \{x \in A : \exists y (y \in B \land xRy)\}$.
            \item $\Img R = R[A] = \{ y \in B : \exists x (x \in A \land xRy)\}$.
            \item $S \circ R \subseteq A \times D$.
        \end{enumerate}
\end{proposition}

\begin{proof}
    \leavevmode
        \begin{enumerate}[leftmargin=*, align=left, label=\textbf{(\alph*)}]
            \item Se $x \in \Dom R$, então existe $y \in \Img R$ tal que $(x,y) \in R$. Como $R \subseteq A \times B$, temos $(x,y) \in A \times B$, donde $x \in A$ e $y \in B$. Assim, $x \in A$ e existe $y \in B$ tal que $xRy$, o que prova a inclusão $\subseteq$. Reciprocamente, se $x \in A$ e existe $y \in B$ tal que $(x,y) \in R$, então $x \in \Dom R$, o que prova a inclusão $\supseteq$. \blackproof
            \item Se $y \in \Img R$, então existe $x \in \Dom R$ tal que $(x,y) \in R$. Como $R \subseteq A \times B$, temos $(x,y) \in A \times B$, donde $x \in A$ e $y \in B$. Assim, $y \in B$ e existe $x \in A$ tal que $xRy$, o que prova a inclusão $\subseteq$. Reciprocamente, se $y \in B$ e existe $x \in A$ tal que $(x,y) \in R$, então $y \in \Img R$, o que prova a inclusão $\supseteq$. \blackproof
            \item Se $(x,z) \in S \circ R$, então $x \in \Dom R$ e $z \in \Img S$. Como $\Dom R \subseteq A$ (primeiro item) e $\Img S \subseteq D$ (segundo item), temos $x \in A$ e $z \in D$, de modo que $(x,z) \in A \times D$. Com isso, $S \circ R \subseteq A \times D$. \blackproof
        \end{enumerate}
\end{proof}

\begin{proposition} \label{prop.set:relacoes}
    Sejam $R$, $S$ e $T$ relações.
        \begin{enumerate}[leftmargin=*, align=left, label=\textbf{(\alph*)}]
            \item 
                \begin{enumerate}[label=\roman*.]
                    \item $(R^{-1})^{-1} = R$.
                    \item $\Dom R^{-1} = \Img R$.
                    \item $\Img R^{-1} = \Dom R$.
                    \item $\id_{\Dom R} \subseteq R^{-1} \circ R$.
                    \item $\id_{\Img R} \subseteq R \circ R^{-1}$.
                \end{enumerate}
            \item 
                \begin{enumerate}[label=\roman*.]
                    \item Se $a \notin \Dom R$, então $R[\{a\}] = \emptyset$.
                    \item Se $b \notin \Img R$, então $R^{-1}[\{b\}] = \emptyset$.
                \end{enumerate}
            \item Para quaisquer $X$ e $Y$,
                \begin{enumerate}[label=\roman*.]
                    \item $R[X \cup Y] = R[X] \cup R[Y]$ e $R^{-1}[X \cup Y] = R^{-1}[X] \cup R^{-1}[Y]$.
                    \item $R[X \cap Y] \subseteq R[X] \cap R[Y]$ e $R^{-1}[X \cap Y] \subseteq R^{-1}[X] \cap R^{-1}[Y]$.
                    \item $R[X \setminus Y] \supseteq R[X] \setminus R[Y]$ e $R^{-1}[X \setminus Y] \supseteq R^{-1}[X] \setminus R^{-1}[Y]$.
                \end{enumerate}
            \item $T \circ (S \circ R) = (T \circ S) \circ R$.
            \item $(S \circ R)^{-1} = R^{-1} \circ S^{-1}$.
            \item Se $\Img R \subseteq \Dom S$, então $\Dom \, (S \circ R) = \Dom R$.
            \item Se $\Dom S \subseteq \Img R$, então $\Img \, (S \circ R) = \Img S$.
        \end{enumerate}
\end{proposition}

\begin{proof}
    \leavevmode
        \begin{enumerate}[leftmargin=*, align=left, label=\textbf{(\alph*)}]
            \item 
                \begin{enumerate}[label=\roman*.]
                    \item Temos $(a,b) \in R$ se, e somente se, $(b,a) \in R^{-1}$, o que é equivalente a $(a,b) \in (R^{-1})^{-1}$. Logo, $R = (R^{-1})^{-1}$.
                    \item 
                    \item 
                    \item 
                    \item 
                \end{enumerate}
            \item 
                \begin{enumerate}[label=\roman*.]
                    \item 
                    \item 
                \end{enumerate}
            \item 
                \begin{enumerate}[label=\roman*.]
                    \item 
                    \item 
                    \item 
                \end{enumerate}
            \item Por um lado ($\subseteq$), se $(a,d) \in T \circ (S \circ R)$, então $a \in \Dom \, (S \circ R)$, $d \in \Img T$ e existe $c \in \Img \, (S \circ R) \cap \Dom T$ tal que $(a,c) \in S \circ R$ e $(c,d) \in T$. De $(a,c) \in S \circ R$ segue que $a \in \Dom R$, $c \in \Img S$ e existe $b \in \Img R \cap \Dom S$ tal que $(a,b) \in R$ e $(b,c) \in S$. Como $b \in \Dom S$, $d \in \Img T$ e existe $c \in \Dom T \cap \Img S$ tal que $(b,c) \in S$ e $(c,d) \in T$, temos que $(b,d) \in T \circ S$. Com isso, $b \in \Dom \, (T \circ S)$ e $d \in \Img \, (T \circ S)$. Como $a \in \Dom R$, $d \in \Img \, (T \circ S)$ e existe $b \in \Dom \, (T \circ S) \cap \Img R$ tal que $(a,b) \in R$ e $(b,d) \in T \circ S$, temos que $(a,d) \in (T \circ S) \circ R$. Com isso, $T \circ (S \circ R) \subseteq (T \circ S) \circ R$. A prova de que $(T \circ S) \circ R \subseteq T \circ (S \circ R)$ é completamente análoga. Logo, $T \circ (S \circ R) = (T \circ S) \circ R$. \blackproof
            \item Por um lado ($\subseteq$), se $(c,a) \in (S \circ R)^{-1}$, então $(a,c) \in S \circ R$, $a \in \Dom R$, $c \in \Img S$ e existe $b \in \Img R \cap \Dom S$ tal que $(a,b) \in R$ e $(b,c) \in S$ isto é, $(c,b) \in S^{-1}$, $(b,a) \in R^{-1}$, com $c \in \Dom S^{-1}$, $a \in \Img R^{-1}$ e $b \in \Img S^{-1} \cap \Dom R^{-1}$. Com isso, $(c,a) \in R^{-1} \circ S^{-1}$, de modo que $(S \circ R)^{-1} \subseteq R^{-1} \circ S^{-1}$. Por outro lado ($\supseteq$), a prova de que $R^{-1} \circ S^{-1} \subseteq (S \circ R)^{-1}$ é completamente análoga. Logo, $(S \circ R)^{-1} = R^{-1} \circ S^{-1}$. \blackproof
            \item Se $a \in \Dom R$, então existe $b \in \Img R$ tal que $(a,b) \in R$. Se $\Img R \subseteq \Dom S$, então $b \in \Dom S$, de modo que existe $c \in \Img S$ tal que $(b,c) \in S$. Assim, $(a,c) \in S \circ R$, de modo que $a \in \Dom \, (S \circ R)$. Com isso, $\Dom R \subseteq \Dom \, (S \circ R)$. Agora, se $a \in \Dom \, (S \circ R)$, então existe $c \in \Img \, (S \circ R)$ tal que $(a,c) \in S \circ R$, donde $a \in \Dom R$ e $\Dom \, (S \circ R) \subseteq \Dom R$. Com isso, $\Dom R = \Dom \, (S \circ R)$. \blackproof
            \item Se $c \in \Img S$, então existe $b \in \Dom S$ tal que $(b,c) \in S$. Se $\Dom S \subseteq \Img R$, então $b \in \Img R$, de modo que existe $a \in \Dom R$ tal que $(a,b) \in R$. Assim, $(a,c) \in S \circ R$, de modo que $c \in \Img \, (S \circ R)$. Com isso, $\Img S \subseteq \Img \, (S \circ R)$. Agora, se $c \in \Img \, (S \circ R)$, então existe $a \in \Dom \, (S \circ R)$ tal que $(a,c) \in S \circ R$, donde $c \in \Img S$ e $\Img \, (S \circ R) \subseteq \Img S$. Com isso, $\Img S = \Img \, (S \circ R)$. \blackproof
        \end{enumerate}
\end{proof}

%\begin{proposition}
%    Seja $R$ uma relação. A imagem inversa de $X$ por $R$ é igual à imagem de $X$ por $R^{-1}$.
%\end{proposition}

\begin{definition}
    \leavevmode
        \begin{enumerate}[leftmargin=*, align=left, label=\textbf{(\alph*)}]
            \item A \emph{relação de identidade} em um conjunto $A$ é definida como
                \[
                    \id_{A} := \{(a,b) \in A \times A : a = b\}.
                \]
            \item A \emph{relação de pertencimento} em um conjunto $A$ é definida como
                \[
                    \in_{A} := \{(a,b) \in A \times A : a \in b\}.
                \]
            \item A \emph{relação de subconjunto} em um conjunto $A$ é definida como
                \[
                    \subseteq_{A} := \{(a,b) \in A \times A : a \subseteq b\}.
                \]
            \item A \emph{relação de subconjunto próprio} em um conjunto $A$ é definida como
                \[
                    \subsetneq_{A} := \{(a,b) \in A \times A : a \subsetneq b\}.
                \]
        \end{enumerate}
\end{definition}

\subsection{Funções}

\begin{definition}[Função] \label{defi.set:função}
    \leavevmode
        \begin{enumerate}[leftmargin=*, align=left, label=\textbf{(\alph*)}]
            \item Uma relação $f$ é uma \emph{função} se
                \[
                    \forall a \forall b \forall c ((a,b) \in f \land (a,c) \in f \rightarrow b=c).
                \]
            A \emph{imagem} de $a \in \Dom f$ por $f$ é denotada por $f(a)$.
            \item Uma \emph{função parcial de $A$ em $B$} é uma função $f$ tal que $\Dom f \subseteq A$ e $\Img f \subseteq B$.
            \item Uma \emph{função total de $A$ em $B$} é uma função $f$ tal que $\Dom f = A$ e $\Img f \subseteq B$. Isso é denotado por $f : A \to B$.
            \item O conjunto de todas as funções totais de $A$ em $B$ é denotado por $B^{A}$, isto é,
                \begin{align*}
                    B^{A} := \{f \in \mathcal{P}{(A \times B)} : & \forall a \forall b \forall c ((a,b) \in f \land (a,c) \in f \rightarrow b=c) \\ 
                    & \land \forall x(x \in A \rightarrow \exists y((x,y) \in f)) \}.
                \end{align*}
        \end{enumerate}
\end{definition}

\begin{definition}
    A \emph{função identidade} de um conjunto $A$ é a função $\id_A : A \to A$ definida por $\id_A{(x)} = x$ para todo $x \in A$.
\end{definition}

\begin{proposition}
    Seja $f: A \to B$. Então $\id_B \circ f = B$ e $f \circ \id_A = A$.
\end{proposition}

\begin{proof}
    \blackproof
\end{proof}

\begin{proposition}
    Sejam $f$ e $g$ funções.
        \begin{enumerate}[leftmargin=*, align=left, label=\textbf{(\alph*)}]
            \item $f \restriction_X$ é uma função, com $\Dom f \restriction_X =  \Dom f \cap X$.
            \item $g \circ f$ é uma função, com $\Dom \, (g \circ f) = \Dom f \cap f^{-1}[\Dom g]$.
        \end{enumerate}
\end{proposition}

\begin{proof}
    \leavevmode
        \begin{enumerate}[leftmargin=*, align=left, label=\textbf{(\alph*)}]
            \item Por definição, $f\restriction_X = \{ (a,b) \in f : a \in X \}$, isto é, $f$ é um conjunto de pares ordenados e, portanto, uma relação. Se $(a,b) \in f\restriction_X$ e $(a,c) \in f\restriction_X$, então, pela definição de $f\restriction_X$, $(a,b) \in f$, $(a,c) \in f$ e $a \in X$; como $f$ é função, $b=c$, de modo que $f\restriction_X$ é também uma função. Provemos, por fim, que $\Dom f\restriction_X =  \Dom f \cap X$.
                \begin{itemize}
                    \item[($\subseteq$)] Se $a \in \Dom f\restriction_X$, então existe $b \in \Img f\restriction_X$ tal que $(a,b) \in f\restriction_X$; pela definição de $f\restriction_X$, vem $(a,b) \in f$ e $a \in X$, e de $(a,b) \in f$ vem $a \in \Dom f$. Com isso, $a \in \Dom f \cap X$, de modo que $\Dom f\restriction_X \subseteq \Dom f \cap X$.
                    \item[($\supseteq$)] Se $a \in \Dom f \cap X$, então de $a \in \Dom f$ segue que existe $b \in \Img f$ tal que $(a,b) \in f$, e como $a \in X$, vem $(a,b) \in f\restriction_X$, de modo que $\Dom f \cap X \subseteq \Dom f\restriction_X$. 
                \end{itemize}
            Logo, $\Dom f\restriction_X = \Dom f \cap X$. Em particular, temos $f\restriction_X = f \restriction_{\Dom f \cap X}$. \blackproof
            \item Se $(a,x) \in g \circ f$ e $(a,y) \in g \circ f$, então existe $b \in \Img f \cap \Dom g$ tal que $(a,b) \in f$ e $(b,x) \in g$ e existe $c \in \Img f \cap \Dom g$ tal que $(a, c) \in f$ e $(c,y) \in g$. Como $f$ é função, vem $b=c$; daí, vem $(b,x) \in g$ e $(b,y) \in g$, e como $g$ é função, vem $x=y$.
                \begin{itemize}
                    \item[($\subseteq$)] Se $a \in \Dom \, (g \circ f)$, então existe $c \in \Img \, (g \circ f)$ tal que $(a,c) \in g \circ f$, e então existe $b \in \Img f \cap \Dom g$ tal que $(a,b) \in f$ e $(b,c) \in g$. De $(a,b) \in f$ segue que $a \in \Dom f$, e de $b \in \Dom g$ com $(a,b) \in f$ segue que $a \in f^{-1}[\Dom g]$. Com isso, $a \in \Dom f \cap f^{-1}[\Dom g]$, de modo que $\Dom \, (g \circ f) \subseteq \Dom f \cap f^{-1}[\Dom g]$.
                    \item[($\supseteq$)] Se $a \in \Dom f \cap f^{-1}[\Dom g]$, então de $a \in \Dom f$ segue que existe $b \in \Img f$ tal que $(a,b) \in f$, e de $a \in f^{-1}[\Dom g]$ segue que existe $b' \in \Dom g$ tal que $(a,b') \in f$. Como $f$ é função, temos $b = b'$, de modo que $b \in \Dom g$ e existe $c \in \Img g$ tal que $(b,c) \in g$, de modo que $(a,c) \in g \circ f$. Com isso, $a \in \Dom \, (g \circ f)$, de modo que $\Dom f \cap f^{-1}[\Dom g] \subseteq \Dom \, (g \circ f)$.  
                \end{itemize}
            Logo, $\Dom \, (g \circ f) = \Dom f \cap f^{-1}[\Dom g]$. \blackproof
        \end{enumerate}
\end{proof}

\begin{definition}
    \leavevmode
        \begin{enumerate}[leftmargin=*, align=left, label=\textbf{(\alph*)}]
            \item Duas funções $f$ e $g$ são \emph{compatíveis} se $f(x) = g(x)$ para todo $x \in \Dom f \cap \Dom g$.
            \item Um \emph{sistema de funções compatíveis} é um conjunto $F$ de funções duas a duas compatíveis.
        \end{enumerate}
\end{definition}

\begin{proposition} \fonte[26]{jechhrbacek1999}
    \leavevmode
        \begin{enumerate}[leftmargin=*, align=left, label=\textbf{(\alph*)}]
            \item Sejam $f$ e $g$ funções. As seguintes afirmações são equivalentes.
                \begin{enumerate}[label=\roman*.]
                    \item $f$ e $g$ são compatíveis.
                    \item $f \cup g$ é uma função.
                    \item $f \restriction_{\Dom f {\cap} \Dom g} = g \restriction_{\Dom f \cap \Dom g}$. 
                \end{enumerate}
            \item Se $F$ é um sistema de funções compatíveis, então $\bigcup F$ é uma função e
                \[
                    \Dom  \left( \bigcup F \right) = \bigcup \{\Dom f : f \in F \}.
                \]
        \end{enumerate}
\end{proposition}

\begin{proof}
    \leavevmode
        \begin{enumerate}[leftmargin=*, align=left, label=\textbf{(\alph*)}]
            \item 
            \item Observe inicialmente que
                \[
                    \{\Dom f : f \in F \} = \left\{z \in \mathcal{P}\left( \bigcup \bigcup \bigcup F \right) : \exists f (f \in F \land z = \Dom f) \right\}.
                \]
        \end{enumerate}
\end{proof}

\subsection{Funções Injetivas}

\begin{definition}
    Uma função $f$ é \emph{injetiva} se
        \[
            \forall x \forall y (x,y \in \Dom f \rightarrow (x \neq y \rightarrow f(x) \neq f(y))).
        \]
\end{definition}

\begin{proposition} \label{prop.set:feginj}
    Sejam $f$ e $g$ funções.
        \begin{enumerate}[leftmargin=*, align=left, label=\textbf{(\alph*)}]
            \item Se $f$ e $g$ são injetivas, então $g \circ f$ é injetiva.
            \item Se $g \circ f$ é injetiva e $\Img f \subseteq \Dom g$, então $f$ é injetiva.
        \end{enumerate}
\end{proposition}

\begin{proof}
    \leavevmode
        \begin{enumerate}[leftmargin=*, align=left, label=\textbf{(\alph*)}]
            \item Sejam $x,y \in \Dom \, (g \circ f)$. Se $g(f(x)) = g(f(y))$, então, pela injetividade de $g$, temos $f(x) = f(y)$. Se $f(x) = f(y)$, então, pela injetividade de $f$, temos $x = y$. Logo, $g \circ f$ é injetiva. \blackproof
            \item Sejam $x,y \in \Dom f$. Se $f(x) = f(y)$, então, como $\Img f \subseteq \Dom g$, temos $g (f(x)) = g(f(y))$. Daí, como $\Dom f = \Dom \, (g \circ f)$ (\cref{prop.set:relacoes}) e $g \circ f$ é injetiva, temos $x=y$. Logo, $f$ é injetiva. \blackproof
        \end{enumerate}
\end{proof}

\begin{theorem} \label{teo.set:inj}
    Seja $f$ uma função.
        \begin{enumerate}[leftmargin=*, align=left, label=\textbf{(\alph*)}]
            \item Se a relação $f^{-1}$ é uma função, então $f^{-1}$ é injetiva.
            \item A relação $f^{-1}$ é uma função se, e somente se,
                \begin{enumerate}[label=\roman*.]
                    \item $f$ é injetiva.
                    \item $f^{-1} \circ f = \id_{\Dom f}$.
                \end{enumerate}
        \end{enumerate}
\end{theorem}

\begin{proof}
    \leavevmode
        \begin{enumerate}[leftmargin=*, align=left, label=\textbf{(\alph*)}]
            \item Se $(y,x) \in f^{-1}$ e $(z,x) \in f^{-1}$, então $(x,y) \in f$ e $(x,z) \in f$, e como $f$ é função, temos $y=z$. Logo, $f^{-1}$ é uma função injetiva. \blackproof
            \item A equivalência mais importante é com $f$ ser injetiva.
                \begin{enumerate}[label=\roman*.]
                    \item Por um lado ($\Rightarrow$), sejam $(y,x) \in f$ e $(z,x) \in f$. Então, $(x,y) \in f^{-1}$ e $(x,z) \in f^{-1}$. Como $f^{-1}$ é uma função, temos $y=z$. Logo, $f$ é injetiva. Por outro lado ($\Leftarrow$), sejam $(x,y) \in f^{-1}$ e $(x,z) \in f^{-1}$. Então, $(y,x) \in f$ e $(z,x) \in f$. Como $f$ é injetiva, temos $y=z$. Logo, $f^{-1}$ é uma função.
                \end{enumerate}
            Provemos que $f$ é injetiva se, e somente se, $f^{-1} \circ f = \id_{\Dom f}$.
                \begin{enumerate}[label=\roman*., resume]
                    \item Por um lado ($\Rightarrow$), se $(x,z) \in f^{-1} \circ f$, então existe $y \in \Img f \cap \Dom f^{-1}$ tal que $(x,y) \in f$ e $(y,z) \in f^{-1}$. Com isso, $(z,y) \in f$, e como $f$ é injetiva, temos $z=x$, de modo que $(x,x) \in \id_{\Dom f}$, isto é, $f^{-1} \circ f \subseteq \id_{\Dom f}$. Agora, se $(x,x) \in \id_{\Dom f}$, então existe $y \in \Img f$ tal que $(x, y) \in f$, isto é, $(y,x) \in f^{-1}$. Com isso, $(x,x) \in f^{-1} \circ f$, e então $\id_{\Dom f} \subseteq f^{-1} \circ f$. Logo, $f^{-1} \circ f = \id_{\Dom f}$. Por outro lado ($\Leftarrow$), sendo $(x,y) \in f$ e $(z,y) \in f$, temos $(y,x) \in f^{-1}$, de modo que $(z,x) \in f^{-1} \circ f$, e como $f^{-1} \circ f = \id_{\Dom f}$, temos $x=z$. Logo, $f$ é injetiva.
                \end{enumerate}
            Com isso, todas as equivalências foram provadas. \blackproof
        \end{enumerate}
\end{proof}

\subsection{Funções Sobrejetivas}

\begin{definition}
    Uma função $f : A \to B$ é \emph{sobrejetiva em $B$} se $\Img f = B$.
\end{definition}

\begin{proposition}
    Uma função $f: A \to B$ é sobrejetiva em $B$ se, e somente se, para todo $y \in B$ existe $x \in A$ tal que $(x,y) \in f$.
\end{proposition}

\begin{proof}
    Segue da \cref{prop.set:domeimdeRAB}. \blackproof    
\end{proof}

\begin{lemma} \label{lem.set:fegsob}
    Sejam $f: A \to B$ e $g: C \to D$ funções.
        \begin{enumerate}[leftmargin=*, align=left, label=\textbf{(\alph*)}]
            \item $\Dom \, (g \circ f) = \{ x \in A : f(x) \in C\}$.
            \item $\Dom \, (g \circ f) = A$ se, e somente se, $f(A) \subseteq C$.
        \end{enumerate}
\end{lemma}

\begin{proof}
    \leavevmode
        \begin{enumerate}[leftmargin=*, align=left, label=\textbf{(\alph*)}]
            \item Pela \cref{prop.set:domeimdeRAB}, temos
                \[
                    \Dom \, (g \circ f) = \{ x \in A : \exists y(y \in D \land (x,y \in g \circ f))\}.
                \]
                \begin{itemize}
                    \item[($\subseteq$)] Se $x \in \Dom \, (g \circ f)$, então existe $y \in D$ tal que $(x,y) \in g \circ f$. Logo, existe $z \in \Img f \cap \Dom g \subseteq B \cap C$ tal que $(x,z) \in f$ e $(z,y) \in g$, isto é, $z = f(x)$ e $y = g(z)$. Com isso, temos $x \in A$ e $f(x) \in C$, de modo que $x \in \{ x \in A : f(x) \in C\}$, isto é, $\Dom \, (g \circ f) \subseteq \{ x \in A : f(x) \in C\}$.
                    \item[($\supseteq$)] Se $x \in \{ x \in A : f(x) \in C\}$, então $x \in A$ e $f(x) \in C$, isto é, existe (um único) $z \in C$ tal que $(x,z) \in f$. Como $z \in C$, existe $y \in \Img g \subseteq D$ tal que $(z,y) \in g$. Com isso, $x \in A$ e existe $y \in D$ tal que $(x,y) \in g \circ f$, de modo que $x \in \Dom \, (g \circ f)$, isto é, $\{ x \in A : f(x) \in C\} \subseteq \Dom{(g \circ f)}$.  
                \end{itemize}
            Logo, $\Dom \, (g \circ f) = \{ x \in A : f(x) \in C\}$. \blackproof
            \item A volta ($\Leftarrow$) já foi provada (\cref{prop.set:relacoes}). Por outro lado ($\Rightarrow$), se $y \in f(A)$, então existe $x \in A$ tal que $y = f(x)$, e como $A = \Dom \, (g \circ f)$, temos $f(x) \in C$, isto é, $y \in C$. Logo, $f(A) \subseteq C$. \blackproof
        \end{enumerate}
\end{proof}

\begin{proposition} \label{prop.set:fegsob}
    Sejam $f : A \to B$ e $g : C \to D$ funções.
        \begin{enumerate}[leftmargin=*, align=left, label=\textbf{(\alph*)}]
            \item Se $f$ e $g$ são sobrejetivas e $B=C$, então $g \circ f$ é sobrejetiva.
            \item Se $g \circ f$ é sobrejetiva e $f(A) \subseteq C$, então $g$ é sobrejetiva.
        \end{enumerate} 
\end{proposition}

\begin{proof}
    \leavevmode
        \begin{enumerate}[leftmargin=*, align=left, label=\textbf{(\alph*)}]
            \item Se $B=C$, então $\Dom \, (g \circ f) = A$ pelo \cref{lem.set:fegsob}. Se $g$ é sobrejetiva, então para todo $z \in D$ existe $y \in C = B$ tal que $(y,z) \in g$. Se $f$ é sobrejetiva, então para esse $y \in B$ existe $x \in A$ tal que $(x,y) \in f$. Com isso, para todo $z \in D$ existe $x \in A$ tal que $(x,z) \in g \circ f$, o que prova a sobrejetividade de $g \circ f$. \blackproof
            \item Se $f(A) \subseteq C$, então $\Dom \, (g \circ f) = A$ pelo \cref{lem.set:fegsob}. Se $g \circ f$ é sobrejetiva, então para todo $z \in D$ existe $x \in A$ tal que $(x,z) \in g \circ f$. Com isso, existe $y \in f(A) \cap C = C$ tal que $(x,y) \in f$ e $(y,z) \in g$. o que prova a sobrejetividade de $g$. \blackproof
        \end{enumerate}
\end{proof}

\begin{lemma} \label{lema.set:fbolafmenos1}
    Para qualquer função $f$, tem-se $f \circ f^{-1} = \id_{\Img f}$.
\end{lemma}

\begin{proof}
    Por um lado ($\subseteq$), $(y,z) \in f \circ f^{-1}$, então existe $x \in \Dom f \cap \Img f^{-1}$ tal que $(y,x) \in f^{-1}$ e $(x,z) \in f$. Daí, $(x,y) \in f$, e como $f$ é uma função, temos $y=z$. Logo $f \circ f^{-1} \subseteq \id_{\Img f}$. Por outro lado ($\supseteq$), se $(y,y) \in \id_{\Img f}$, então existe $x \in \Dom f$ tal que $(x,y) \in f$. Logo $(y,x) \in f^{-1}$, de modo que $(y, y) \in f \circ f^{-1}$, isto é, $f \circ f^{-1} \supseteq \id_{\Img f}$. Logo, vem $f \circ f^{-1} = \id_{\Img f}$. \blackproof
\end{proof}

\begin{theorem} \label{teo.set:sob}
    Uma função $f : A \to B$ é sobrejetiva em $B$ se, e somente se, $f \circ f^{-1} = \id_B$.
\end{theorem}

\begin{proof}
    \leavevmode
        \begin{itemize}[align=left]
            \item[($\Rightarrow$)] Se $f$ é sobrejetiva em $B$, então $\Img f = B$, de modo que, pelo \cref{lema.set:fbolafmenos1}, $f \circ f^{-1} = \id_B$.
            \item[($\Leftarrow$)] Se $f \circ f^{-1} = \id_B$, então $f$ é sobrejetiva em $B$ porque, como também $f \circ f^{-1} = \id_{\Img f}$ (\cref{lema.set:fbolafmenos1}), temos $\Img f = B$.
        \end{itemize}
     \blackproof
\end{proof}

\subsection{Funções Bijetivas}

\begin{definition}
    Uma função $f: A \to B$ é \emph{bijetiva em relação a $B$} se é injetiva e sobrejetiva em $B$.
\end{definition}

\begin{proposition} \label{prop.set:defbij}
    Uma função $f: A \to B$ é bijetiva em $B$ se, e somente se, para todo $y \in B$ existe um único $x \in A$ tal que $(x,y) \in f$.
\end{proposition}

\begin{proof}
    Segue imediatamente da definição. \blackproof
\end{proof}

\begin{proposition}
    Se $f : A \to B$ e $g : B \to C$ são funções bijetivas, então a função $g \circ f : A \to C$ é uma função bijetiva.
\end{proposition}

\begin{proof}
    Segue como corolário imediato das \cref{prop.set:feginj,prop.set:fegsob}. \blackproof
\end{proof}

\begin{theorem} \label{teo.set:bij1}
    Seja $f : A \to B$ uma função.
        \begin{enumerate}[leftmargin=*, align=left, label=\textbf{(\alph*)}]
            \item Se a relação $f^{-1}$ é uma função de $B$ em $A$, então $f^{-1}$ é bijetiva em $A$.
            \item A relação $f^{-1}$ é uma função de $B$ em $A$ se, e somente se,
                \begin{enumerate}[label=\roman*.]
                    \item $f$ é bijetiva em $B$.
                    \item $f^{-1} \circ f = \id_{A}$ e $f \circ f^{-1} = \id_{B}$.
                \end{enumerate}
        \end{enumerate}
\end{theorem}

\begin{proof}
    \leavevmode
        \begin{enumerate}[leftmargin=*, align=left, label=\textbf{(\alph*)}]
            \item Como $f$ é função, para todo $x \in A$ existe um único $y \in B$ tal que $(x,y) \in f$, isto é, $(y,x) \in f^{-1}$. Logo, pela \cref{prop.set:defbij}, $f^{-1}$ é bijetiva em $A$. \blackproof
            \item A equivalência que mais importa é com $f$ ser bijetiva em $B$.
                \begin{enumerate}[label=\roman*.]
                    \item Por um lado ($\Rightarrow$), se $f^{-1} \subseteq B \times A$ é uma função com $\Dom f^{-1} = B$, então para todo $y \in B$ existe um único $x \in A$ tal que $(y,x) \in f^{-1}$, isto é, $(x,y) \in f$. Logo, pela \cref{prop.set:defbij}, $f$ é bijetiva. Por outro lado ($\Leftarrow$), se $f$ é bijetiva, então, pela \cref{prop.set:defbij}, para todo $y \in B$ existe um único $x \in A$ tal que $(x,y) \in f$, isto é, $(y,x) \in f^{-1}$. Daí, pela definição de função, $f^{-1}$ é uma função de $B$ em $A$. \blackproof
                \end{enumerate}
            Provemos que $f$ é bijetiva em $B$ se, e somente se, $f^{-1} \circ f = \id_A$ e $f \circ f^{-1} = \id_{B}$.
                \begin{enumerate}[label=\roman*., resume]
                    \item Por um lado ($\Rightarrow$), se $f$ é bijetiva em $B$, então $f$ é injetiva e sobrejetiva em $B$, de modo que, pelos \cref{teo.set:inj,teo.set:sob}, temos $f^{-1} \circ f = \id_A$ e $f \circ f^{-1} = \id_B$, respectivamente. Agora, se $f^{-1} \circ f = \id_A$ e $f \circ f^{-1} = \id_B$, então pelos mesmos teoremas $f$ é injetiva e sobrejetiva em $B$, isto é, $f$ é bijetiva em $B$. \blackproof
                \end{enumerate}
            Com isso, todas as equivalências foram provadas. \blackproof
        \end{enumerate}
\end{proof}

\subsection{Funções Inversas}

\begin{definition}
    \leavevmode
        \begin{enumerate}[leftmargin=*, align=left, label=\textbf{(\alph*)}]
            \item Uma função $f$ é \emph{invertível à esquerda} se existe uma função $g$ tal que $g \circ f = \id_{\Dom f}$. Dizemos, nesse caso, que $g$ é uma \emph{inversa à esquerda} de $f$.
            \item Uma função $f : A \to B$ é \emph{invertível à esquerda} se existe uma função $g : B \to A$ tal que $g \circ f = \id_A$. Dizemos, nesse caso, que $g$ é uma \emph{inversa à esquerda} de $f$.
        \end{enumerate}
\end{definition}

\begin{theorem} \label{teo.set:invesquerda}
    \leavevmode
        \begin{enumerate}[leftmargin=*, align=left, label=\textbf{(\alph*)}]
            \item Uma função $f$ é invertível à esquerda se, e somente se, $f$ é injetiva.
            \item Uma função $f : A \to B$ é invertível à esquerda se, e somente se, $f$ é injetiva.
        \end{enumerate}
\end{theorem}

\begin{proof}
    \leavevmode
        \begin{enumerate}[leftmargin=*, align=left, label=\textbf{(\alph*)}]
            \item Por um lado ($\Rightarrow$), se $f$ é injetiva, então, pelo \cref{teo.set:inj}, $f^{-1}$ é função e $f^{-1} \circ f = \id_{\Dom f}$. Logo, $f$ é invertível à esquerda. Por outro lado ($\Leftarrow$), sejam $(x,y) \in f$ e $(z,y) \in f$. Como $f$ é invertível à esquerda, existe uma função $g$ tal que $g \circ f = \id_{\Dom f}$. Como $(x,x) \in g \circ f$, existe $w \in \Img f \cap \Dom g$ tal que $(x,w) \in f$ e $(w,x) \in g$. Como $f$ é função, temos $w=y$, de modo que $(y,x) \in g$. Analogamente, temos $(y,z) \in g$, e como $g$ é função, temos $x=z$. Logo, $f$ é injetiva. \blackproof
            \item 
        \end{enumerate}
\end{proof}

\begin{definition}
    Uma função $f:A \to B$ é \emph{invertível à direita} se existe uma função $g : B \to A$ tal que $f \circ g = \id_{B}$. Dizemos, nesse caso, que $g$ é uma \emph{inversa à direita} de $f$.
\end{definition}

\begin{theorem} \label{teo.set:invdireita}
    Uma função $f:A \to B$ é invertível à direita se, e somente se, $f$ é sobrejetiva em $B$.
\end{theorem}

\begin{remark}
    A prova da volta ($\Leftarrow$) deste teorema depende do axioma da escolha. Mais precisamente, de um enunciado equivalente ao axioma da escolha: para toda relação $R$ existe uma função $f \subseteq R$ tal que $\Dom f = \Dom R$. Ainda assim, enunciamos este resultado aqui por uma questão de organização didática.  
\end{remark}

\begin{proof}
    \blackproof
\end{proof}

\begin{definition}
    Uma função $f: A \to B$ é \emph{invertível} se existe uma função $g : B \to A$ tal que $g \circ f = \id_A$ e $f \circ g = \id_B$. Dizemos, nesse caso, que $g$ é a \emph{inversa} de $f$.
\end{definition}

\begin{lemma} \label{lema.set:invesqdiriguais}
    As inversas laterais de uma função invertível à esquerda e à direita coincidem.
\end{lemma}

\begin{proof}
    Seja $f : A \to B$ invertível à esquerda e à direita. Sejam $g_1 : B \to A$ a inversa à esquerda de $f$ e $g_2 : B \to A$ a inversa à direta de $f$. Temos
        \[
            g_1 = g_1 \circ \id_B = g_1 \circ (f \circ g_2) = (g_1 \circ f) \circ g_2 = \id_A \circ g_2 = g_2.    
        \]
    Logo, as inversas à esquerda e à direita de $f$ são iguais. \blackproof
\end{proof}

\begin{proposition}
    \leavevmode
        \begin{enumerate}[leftmargin=*, align=left, label=\textbf{(\alph*)}]
            \item Uma função invertível tem uma única inversa.\footnote{Note que é esta proposição que nos permite dizer ``a função inversa'' em vez de ``uma função inversa''. }
            \item Uma função é invertível se, e somente se, é invertível à esquerda e à direita.
        \end{enumerate}
\end{proposition}

\begin{proof}
    \leavevmode
        \begin{enumerate}[leftmargin=*, align=left, label=\textbf{(\alph*)}]
            \item Seja $f : A \to B$ invertível. Sejam $g_1, g_2 : B \to A$ inversas de $f$. Provemos que $g_1 = g_2$. De fato,
                \[
                    g_1 = g_1 \circ \id_B = g_1 \circ (f \circ g_2) = (g_1 \circ f) \circ g_2 = \id_A \circ g_2 = g_2.    
                \]
            Logo, a função inversa da função $f : A \to B$, quando existe, é única. \blackproof
            \item A ida ($\Rightarrow$) segue imediatamente das definições. A volta ($\Leftarrow$) também, considerando o \cref{lema.set:invesqdiriguais}. \blackproof 
        \end{enumerate}
\end{proof}

\begin{theorem} \label{teo.set:bij2}
    Seja $f: A \to B$.
        \begin{enumerate}[leftmargin=*, align=left, label=\textbf{(\alph*)}]
            \item $f$ é invertível se, e somente se, $f$ é bijetiva.
            \item Se $f$ é invertível, então a função inversa de $f$ é a relação inversa $f^{-1}$.
        \end{enumerate}
\end{theorem}

\begin{proof}
    \leavevmode
        \begin{enumerate}[leftmargin=*, align=left, label=\textbf{(\alph*)}]
            \item Por um lado ($\Rightarrow$), se $f$ é invertível, então existe $g: B \to A$ tal que $g \circ f = \id_A$ e $f \circ g = \id_B$. Logo, $f$ é invertível à esquerda e à direita, isto é, $f$ é injetiva (\cref{teo.set:invesquerda}) e sobrejetiva em $B$ (\cref{teo.set:invdireita}), isto é, $f$ é bijetiva em $B$. Por outro lado ($\Leftarrow$), se $f$ é bijetiva em $B$, então, pelo \cref{teo.set:bij1}, sua relação inversa $f^{-1} \subseteq B \times A$ é uma função de $B$ em $A$ e, mais ainda, satisfaz $f^{-1} \circ f = \id_{A}$ e $f \circ f^{-1} = \id_{B}$. Logo, $f$ é invertível. \blackproof   
            \item Segue imediatamente do primeiro item. \blackproof
        \end{enumerate}
\end{proof}

\subsection{Relações de Ordem}

\begin{definition}
    Seja $R$ uma relação em $X$.
        \begin{enumerate}[leftmargin=*, align=left, label=\textbf{(\alph*)}]
            \item Dizemos que $R$ é \emph{reflexiva} se
                \[
                    \forall x (x \in X \rightarrow (x,x) \in R).
                \]
            \item Dizemos que $R$ é \emph{irreflexiva} se
                \[
                    \forall x (x \in X \rightarrow (x,x) \notin R).
                \]
            \item Dizemos que $R$ é \emph{simétrica} se
                \[
                    \forall x \forall y (x,y \in X \rightarrow ((x,y) \in R \rightarrow (y,x) \in R)).
                \]
            \item Dizemos que $R$ é \emph{antissimétrica} se
                \[
                    \forall x \forall y (x,y \in X \rightarrow ((x,y) \in R \land (y,x) \in R \rightarrow x=y)).
                \]
            \item Dizemos que $R$ é \emph{assimétrica} se
                \[
                    \forall x \forall y (x,y \in X \rightarrow ( (x,y) \in R \rightarrow (y,x) \notin R)).
                \]
            \item Dizemos que $R$ é \emph{transitiva} se
                \[
                    \forall x \forall y \forall z (x,y,z \in X \rightarrow ((x,y) \in R \land (y,z) \in R \rightarrow (x,z) \in R)).
                \]
        \end{enumerate}
\end{definition}

\begin{definition}[Ordem Parcial] \label{defi.set:ordemparcial}
    \leavevmode
        \begin{enumerate}[leftmargin=*, align=left, label=\textbf{(\alph*)}]
            \item Uma \emph{relação de ordem parcial} em $X$ é uma relação ${\leq} \subseteq X \times X$ que tem as seguintes propriedades.
                \begin{enumerate}[label=\roman*.]
                    \item Reflexividade: $\forall x(x \in X \rightarrow (x,x) \in \leq)$.
                    \item Antissimetria: $\forall x \forall y (x,y \in X \rightarrow (x \leq y \land y \leq x \rightarrow x=y))$.
                    \item Transitividade: $\forall x \forall y \forall z (x,y,z \in X \rightarrow (x \leq y \land y \leq z \rightarrow x \leq z))$. 
                \end{enumerate}
            Dizemos que $X$ é o \emph{domínio} de $\leq$.
            \item Um \emph{conjunto parcialmente ordenado} é um par $(X, \leq)$, onde ${\leq} \subseteq X \times X$ é uma relação de ordem parcial.
        \end{enumerate}
\end{definition}

\begin{example}
    Para todo $X$, a relação $\subseteq_X$ é uma ordem parcial em $X$.
\end{example}

\begin{definition}[Ordem Estrita] \label{defi.set:ordemestrita}
    \leavevmode
        \begin{enumerate}[leftmargin=*, align=left, label=\textbf{(\alph*)}]
            \item Uma \emph{relação de ordem estrita} em $X$ é uma relação ${<} \subseteq X \times X$ que tem as seguintes propriedades.
                \begin{enumerate}[label=\roman*.]
                    \item Assimetria: $\forall x \forall y (x,y \in X \rightarrow ((x,y) \in < \rightarrow (y,x) \notin <))$.
                    \item Transitividade: $\forall x \forall y \forall z (x,y,z \in X \rightarrow ( (x, y) \in < \land (y, z) \in < \rightarrow (x,z) \in <))$. 
                \end{enumerate}
            Dizemos que $X$ é o \emph{domínio} de $<$.
            \item Um \emph{conjunto estritamente ordenado} é um par $(X, <)$, onde ${<} \subseteq X \times X$ é uma relação de ordem estrita.
        \end{enumerate}
\end{definition}

\begin{example}
    Para todo $X$, a relação $\subsetneq_X$ é uma ordem estrita em $X$.
\end{example}

\begin{proposition} \label{prop.set:parcial_e_estrita}
    \leavevmode
        \begin{enumerate}[leftmargin=*, align=left, label=\textbf{(\alph*)}]
            \item Seja $\leq$ uma ordem parcial em $X$. A relação
                \[
                    {<} := \{(x,y) \in X \times X : x \leq y \land x \neq y \}
                \]
            é uma ordem estrita em $X$.
            \item Seja $<$ uma ordem estrita em $X$. A relação
                \[
                    {\leq} := \{(x,y) \in X \times X : x<y \lor x=y \}
                \]
            é uma ordem parcial em $X$.
        \end{enumerate}
\end{proposition}

\begin{proof}
    \leavevmode
        \begin{enumerate}[leftmargin=*, align=left, label=\textbf{(\alph*)}]
            \item 
        \end{enumerate}
\end{proof}

\begin{definition}
    Sejam $(X, \leq)$ parcialmente ordenado e $S \in \mathcal{P}(X)_{\neq \emptyset}$.
        \begin{enumerate}[leftmargin=*, align=left, label=\textbf{(\alph*)}]
            \item Dizemos que $M \in X$ é
                \begin{enumerate}[label=\roman*.]
                    \item uma \emph{cota superior} de $S$ se
                        \[
                            \forall x (x \in S \rightarrow x \leq M).
                        \]
                    Dizemos, nesse caso, que $S$ é \emph{limitado superiormente} em $X$.
                    \item o \emph{elemento máximo} de $S$ se
                        \[
                            \forall x (x \in S \rightarrow x \leq M) \land M \in S,
                        \]
                    isto é, se $M$ é uma cota superior de $S$ e $M \in S$. Isso é denotado por $\max{S} := M$.
                    \item o \emph{supremo} de $S$ se
                        \[
                            \forall x (x \in S \rightarrow x \leq M) \land \forall y (y \in X \land \forall x (x \in S \rightarrow x \leq y) \rightarrow M \leq y),
                        \]
                    isto é, se $M$ é a menor cota superior de $S$. Isso é denotado por $\sup{S} := M$.
                \end{enumerate}
            \item Dizemos que $m \in X$ é
                \begin{enumerate}[label=\roman*.]
                    \item uma \emph{cota inferior} de $S$ se
                        \[
                            \forall x (x \in S \rightarrow m \leq x).
                        \]
                    Dizemos, nesse caso, que $S$ é \emph{limitado inferiormente} em $X$.
                    \item o \emph{elemento mínimo} de $S$ se
                        \[
                            \forall x (x \in S \rightarrow m \leq x) \land m \in S,
                        \]
                    isto é, se $m$ é uma cota inferior de $S$ e $m \in S$. Isso é denotado por $\min{S} := m$.
                    \item o \emph{ínfimo} de $S$ se
                        \[
                            \forall x (x \in S \rightarrow m \leq x) \land \forall y (y \in X \land \forall x (x \in S \rightarrow y \leq x) \rightarrow y \leq m),
                        \]
                    isto é, se $m$ é a maior cota inferior de $S$. Isso é denotado por $\inf{S} := m$.
                \end{enumerate}
        \end{enumerate}
\end{definition}

\begin{remark}
    É fácil ver que o máximo de $S$, quando existe, é único. De fato, se $x, y \in S$ são máximos de $S$, então $x \leq y$ e $y \leq x$, de modo que $x = y$. Essa unicidade também vale para o supremo, o mínimo e o ínfimo de $S$, quando existem. Isso justifica o uso do artigo ``o'' (em \emph{o} máximo, em vez de \emph{um} máximo, por exemplo) e nos permite denotar esses elementos por $\max S$, $\sup S$, $\min S$ e $\inf S$, respectivamente.
\end{remark}

\begin{definition}
    Seja $(X, \leq)$ parcialmente ordenado.
        \begin{enumerate}[leftmargin=*, align=left, label=\textbf{(\alph*)}]
            \item Dizemos que $x,y \in X$ são \emph{comparáveis} se $x \leq y$ ou $y \leq x$.
            \item Dizemos que $\leq$ é uma relação de ordem \emph{total}\footnote{O termo \emph{ordem linear} também costuma ser usado. Dizemos, nesse caso, que o par $(X,\leq)$ é um \emph{conjunto linearmente ordenado}.}, e que o par $(X,\leq)$ é um \emph{conjunto totalmente ordenado}, se quaisquer dois elementos de $X$ são comparáveis, isto é, se
                \[
                    \forall x \forall y (x,y \in X \rightarrow (x \leq y \lor y \leq x)).
                \]
        \end{enumerate}
\end{definition}

\begin{definition}
    Seja $(X, \leq)$ parcialmente ordenado.
        \begin{enumerate}[leftmargin=*, align=left, label=\textbf{(\alph*)}]
            \item Dizemos que $\leq$ é uma \emph{boa ordem} em $X$, e que o par $(X,\leq)$ é um \emph{conjunto bem-ordenado}, se todo subconjunto não vazio de $X$ possui um elemento mínimo.
            \item Dizemos que $\leq$ é um \emph{reticulado} se para quaisquer $x,y \in X$, o conjunto $\{ x,y \}$ possui supremo e ínfimo.
            \item Dizemos que $\leq$ é uma \emph{árvore} se, para todo $x \in X$, o conjunto $S = \{ y \in X : y \leq x\}$ é tal que $(S, \leq \cap \ S^2)$ é um conjunto bem-ordenado.
        \end{enumerate}
\end{definition}

\begin{proposition}
    \leavevmode
        \begin{enumerate}[leftmargin=*, align=left, label=\textbf{(\alph*)}]
            \item Toda boa ordem é uma ordem total.
            \item Toda boa ordem é uma árvore.
            \item Toda ordem total é um reticulado.
        \end{enumerate}
\end{proposition}

\begin{proof}
    \leavevmode
        \begin{enumerate}[leftmargin=*, align=left, label=\textbf{(\alph*)}]
            \item Sejam $x, y \in X$. Como $\leq$ é uma boa ordem, todo subconjunto não vazio de $X$ possui elemento mínimo. Assim, o conjunto $\{x,y\} \subseteq X$ possui um elemento mínimo. Se $\min\{x,y\} = x$, então $x \leq y$. Se $\min\{x,y\} = y$, então $y \leq x$. Em qualquer caso, vale $x \leq y$ ou $y \leq x$, de modo que a ordem $\leq$ é total. \blackproof
            \item Pela definição de árvore, precisamos mostrar que, para todo $x \in X$, o conjunto $S = \{ y \in X : y \leq x\}$ é bem-ordenado pela ordem induzida. Seja $A \subseteq S$ um subconjunto não vazio. Como $S \subseteq X$, temos que $A \subseteq X$. Como $X$ é bem-ordenado, $A$ possui um elemento mínimo. Com isso, todo subconjunto não vazio de $S$ possui um elemento mínimo, isto é, $S$ é bem-ordenado. \blackproof
            \item Sejam $x, y \in X$. Como a ordem $\leq$ é total, temos $x \leq y$ ou $y \leq x$. Suponha, sem perda de generalidade, que $x \leq y$. Como $x \leq x$ e $x \leq y$, temos que $x$ é uma cota inferior de $\{x,y\}$. Se $z$ uma cota inferior qualquer de $\{x,y\}$, então, em particular, $z \leq x$. Logo, $x$ é a maior das cotas inferiores, isto é, $x = \inf\{x,y\}$. Como $y \geq x$ e $y \geq y$, temos que $y$ é uma cota superior de $\{x,y\}$. Se $w$ uma cota superior qualquer de $\{x,y \}$, então, em particular, $w \geq y$. Logo, $y$ é a menor das cotas superiores, isto é, $y = \sup\{x,y\}$. Portanto, $\{x,y\}$ possui supremo e ínfimo. \blackproof
        \end{enumerate}
\end{proof}

\begin{proposition}
    Se $(X,\leq)$ é parcialmente ordenado e $Y \subseteq X$, então $(Y, \leq \cap \ Y^2)$ é parcialmente ordenado.
\end{proposition}

\begin{proof}
    Defina $\leq_Y \ := \ \leq \cap \ Y^2$.
        \begin{enumerate}[label=\roman*.]
            \item (Reflexividade) Dado $y \in Y$, como $Y \subseteq X$, temos $y \in X$, donde $y \leq y$. Como $y \in Y$, vem $(y,y) \in \leq_Y$.
            \item (Antissimetria) Sejam $x, y \in Y$ tais que $(x,y) \in \leq_Y$ e $(y,x) \in \leq_Y$. Então $x \leq y$ e $y \leq x$. Pela antissimetria de $\leq$, vem $x=y$.
            \item (Transitividade) Sejam $x, y, z \in Y$ tais que $(x,y) \in \leq_Y$ e $(y,z) \in \leq_Y$. Então $x \leq y$ e $y \leq z$. Pela transitividade de $\leq$, vem $x \leq z$. Como $x, z \in Y$, vem $(x,z) \in \leq_Y$.
        \end{enumerate}
    Com isso, $\leq_Y \ := \ \leq \cap \ Y^2$ é uma ordem. \blackproof
\end{proof}

\begin{definition}
    Sejam $(X,\leq)$ parcialmente ordenado e $Y \subseteq X$. 
        \begin{enumerate}[leftmargin=*, align=left, label=\textbf{(\alph*)}]
            \item Dizemos que $\leq_Y \ := \ \leq \cap \ Y^2$ é uma \emph{subordem} de $\leq$.
            \item Dizemos que $(Y, \leq_Y)$ é um \emph{subconjunto parcialmente ordenado} de $(X, \leq)$. Isso é denotado por $(Y, \leq)$.
        \end{enumerate}
\end{definition}

\begin{proposition}
    Sejam $(X,\leq)$ parcialmente ordenado e $Y \subseteq X$.
        \begin{enumerate}[leftmargin=*, align=left, label=\textbf{(\alph*)}]
            \item Se $\leq$ é uma ordem total, então $\leq_Y$ é uma ordem total.
            \item Se $\leq$ é uma boa ordem, então $\leq_Y$ é uma boa ordem.
            \item Se $\leq$ é uma árvore, então $\leq_Y$ é uma árvore.
        \end{enumerate}
\end{proposition}

\begin{proof}
    \leavevmode
        \begin{enumerate}[leftmargin=*, align=left, label=\textbf{(\alph*)}]
            \item Sejam $x, y \in Y$. Como $Y \subseteq X$, temos $x, y \in X$, e como $\leq$ é total em $X$, temos $x \leq y$ ou $y \leq x$, de modo que $x \leq_Y y$ ou $y \leq_Y x$. \blackproof
            \item Seja $A \in \mathcal{P}(Y)_{\neq \emptyset}$. Como $Y \subseteq X$, temos $A \subseteq X$, e como $\leq$ é uma boa ordem em $X$, existe $\min{A} \in A$. Como para todo $z \in A \subseteq Y$ temos $\min{A} \leq z$, segue que $\min{A}$ é o mínimo de $A$ com relação a $\leq_Y$. \blackproof
            \item Sendo $x \in Y$ e $S := \{y \in Y : y \leq_Y x \}$, provemos que $(S, \leq_Y \cap \ S^2)$ é bem ordenado. Pondo $T := \{y \in X : y \leq x \}$, temos $S = T \cap Y$, e como $x \in X$ e $\leq$ é uma árvore, temos que $(T, \leq \cap \ T^2)$ é bem ordenado. Como $S \subseteq T$, temos que $(S, (\leq \cap \ T^2) \cap S^2 )$ é também uma boa ordem. Observando que
                \begin{align*}
                    \leq_Y \cap \ S^2 &= (\leq \cap \ Y^2) \cap (T \cap Y)^2 \\
                    &= \leq \cap \ Y^2 \cap T^2 \\
                    &= (\leq \cap \ T^2) \cap (T^2 \cap Y^2) \\
                    &= (\leq \cap \ T^2) \cap S^2, 
                \end{align*}
            segue que $(S, \leq_Y \cap \ S^2)$ é bem ordenado. \blackproof
        \end{enumerate}
\end{proof}

\begin{definition}
    Dois conjuntos parcialmente ordenados $(X_1, \leq_1)$ e $(X_2, \leq_2)$ são \emph{ordem-isomorfos} se existe uma bijeção $f:X_1 \to X_2$ tal que 
        \[
          \forall x \forall y (x,y \in X_1 \rightarrow (x \leq_1 y \leftrightarrow f(x) \leq_2 f(y))).  
        \]
    Dizemos, nesse caso, que $f$ é um \emph{isomorfismo de ordens parciais}.
\end{definition}

\begin{theorem}
    Se $(X, \leq)$ é um conjunto ordenado, então existe um conjunto ordenado $(Y, \preceq)$ ordem-isomorfo a $(X,\leq)$ tal que 
        \[
            \preceq \, \, = \{ (x,y) \in Y \times Y : x \subseteq y \}.
        \]
\end{theorem}

\begin{proof}
    Definindo $f : X \to \mathcal{P}{(X)}$ por $f(x) := \{y \in X : y \leq x \}$, temos que $f$ é bijetiva em relação a $Y := \Img{(f)}$. De fato, se $f(x) = f(y)$, então $x \in f(y)$ e $y \in f(x)$ já que $x \in f(x)$ e $y \in f(y)$; daí, pela definição de $f$, vem $x \leq y$ e $y \leq x$, de modo que $x = y$ e $f$ é injetiva. Provemos então que $x \leq y$ se, e somente se, $f(x) \subseteq f(y)$, para quaisquer $x,y \in X$. Se $z \in f(x)$, então $z \leq x$, e se $x \leq y$, então $z \leq y$, de modo que $z \in f(y)$ e $f(x) \subseteq f(y)$. Agora, se $x \in f(x) \subseteq f(y)$, então $x \in f(y)$, donde $x \leq y$. Com isso, $(X, \leq)$ é ordem-isomorfo a $(Y, \subseteq)$, como havíamos afirmado. \blackproof
\end{proof}

\section{O Axioma do Infinito e os Números Naturais}

\subsection{Sistemas de Peano}

\begin{definition}
    Um \emph{sistema de Peano} é uma tripla $(N, s, a)$ com as seguintes propriedades.
        \begin{enumerate}[label=\roman*.]
            \item $s: N \to N$ é uma função injetiva.
            \item $a \in N \setminus s[N]$.
            \item Para todo $A \subseteq N$, se $a \in A$ e $s(x) \in A$ sempre que $x \in A$, então $A = N$.
        \end{enumerate}
\end{definition}

\begin{example}
    Seja $(N,s,a)$ um sistema de Peano.
        \begin{enumerate}[leftmargin=*, align=left, label=\textbf{(\alph*)}]
            \item Para todo $n \in N$, $s(n) \neq n$.
            \item Para todo $n \in N_{\neq a}$ existe um único $k \in N$ tal que $n = s(k)$.
        \end{enumerate}
\end{example}

\begin{proof}
    \leavevmode
        \begin{enumerate}[leftmargin=*, align=left, label=\textbf{(\alph*)}]
            \item Defina $A := \{n \in N : s(n) \neq n\}$. Como $a \notin s[N]$ e $s(a) \in s[N]$, temos $s(a) \neq a$, de modo que $a \in A$. Agora, supondo que $n \in A$, provemos que $s(n) \in A$. Se $s(s(n)) = s(n)$, então, como $s$ é injetiva, teríamos $s(n)=n$, uma contradição, pois $n \in A$. Logo, $s(s(n)) \neq s(n)$, isto é, $s(n) \in A$. Logo, pelo princípio de indução, temos $A = N$. Logo, $s(n) \neq n$, para todo $n \in N$. \blackproof
            \item Defina
                \[
                    A := \{n \in N : n=a \text{ ou existe } k \in N \text{ tal que } n=s(k)\}.
                \]
            Como $a=a$, temos $a \in A$. Agora, supondo que $n \in A$, provemos que $s(n) \in A$. De fato, tomando $k=n$, temos $s(n)=s(k)$, de modo que $s(n) \in A$. Logo, pelo princípio de indução, temos $A=N$. Assim, se $n \in N_{\neq a}$, então $n \neq a$, e como $n \in A$, existe $k \in N$ tal que $n=s(k)$. Ademais, se $k,l \in N$ são tais que $n=s(k)$ e $n=s(l)$, então $s(k)=s(l)$; como $s$ é injetiva, temos $k=l$. Logo, para todo $n \in N_{\neq a}$ existe um único $k \in N$ tal que $n=s(k)$. \blackproof
        \end{enumerate}
\end{proof}

\begin{theorem}[Recursão] \label{teo.set:recursao_peano}
    Seja $(N,s,a)$ um sistema de Peano. Sejam $X \neq \emptyset$, $x_0 \in X$ e $f:X \to X$. Existe uma única $\varphi : N \to X$ tal que
        \begin{itemize}
            \item $\varphi(a) = x_0$;
            \item $\varphi(s(n)) = f(\varphi(n))$, para todo $n \in N$.
        \end{itemize}
\end{theorem}

\begin{proof} \fonte[113]{fajardo2024conjuntos}
    A ideia é considerar todas as relações de $N$ em $X$ que têm as propriedades desejadas e provar que a interseção de todas elas resulta em uma única função de domínio $N$. Defina
        \[
            \mathcal{C} := \left\{ R \subseteq N \times X : (a,x_0) \in R \land \forall n \forall x ((n,x) \in R \rightarrow (s(n),f(x)) \in R) \right\}.
        \]
    Temos $N \times X \in \mathcal{C}$; logo, $\mathcal{C} \neq \emptyset$, e então existe $\varphi := \bigcap \mathcal{C}$. Essa é a função procurada. Provemos isso.
        \begin{claim}
           $\varphi \in \mathcal{C}$.
        \end{claim}
        \begin{midproof}
            Como $(a,x_0) \in R$ para toda $R \in \mathcal{C}$, temos $(a,x_0) \in \varphi$. Analogamente, se $(n,x) \in \varphi$, então $(n,x) \in R$ para toda $R \in \mathcal{C}$, de modo que $(s(n),f(x)) \in R$ para toda $R \in \mathcal{C}$, isto é, $(s(n),f(x)) \in \varphi$. Logo, $\varphi \in \mathcal{C}$. \whiteproof
        \end{midproof}
        \begin{claim}
            $\Dom \varphi = N$.
        \end{claim}
        \begin{midproof}
            Por indução. Temos $\Dom \varphi \subseteq N$. Como $(a,x_0) \in \varphi$, temos $a \in \Dom \varphi$. Agora, se $n \in \Dom\varphi$, então existe $x \in X$ tal que $(n,x) \in \varphi$. Com isso, vem $(s(n),f(x)) \in \varphi$, donde $s(n) \in \Dom \varphi$. Logo, pelo princípio de indução, temos $\Dom \varphi = N$. \whiteproof
        \end{midproof}
        \begin{claim}
            A relação $\varphi \subseteq N \times X$ é uma função.
        \end{claim}
        \begin{midproof}
            Por indução. Seja
                \[
                    A := \{n \in N : \forall x \forall y ((n,x) \in \varphi \land (n,y) \in \varphi \rightarrow x=y) \}.
                \]
            Temos $(a, x_0) \in \varphi$. Suponha, por absurdo, que existe $y \in X \setminus \{x_0 \}$ tal que $(a, y) \in \varphi$. Então $(a, x_0) \in \varphi \setminus \{(a,y)\}$, e se $(m,z) \in \varphi \setminus \{(a,y)\}$, então $(s(m), f(z)) \in \varphi$, e como $s(m) \neq a$ para todo $m \in N$ pois $a \in N \setminus s[N]$, temos $(s(m), f(z)) \in \varphi \setminus \{(a,y)\}$, de modo que $\varphi \setminus \{(a,y)\} \in \mathcal{C}$, uma contradição. De fato, teríamos $\varphi \subseteq \varphi \setminus \{(a,y)\}$, e como $\varphi \setminus \{(a,y)\} \subseteq \varphi$, teríamos $\varphi \setminus \{(a,y)\} = \varphi$, donde $(a,y) \notin \varphi$, o que contraria $(a,y) \in \varphi$. Logo, se $(a,y) \in \varphi$, então $y=x_0$, isto é, $a \in A$. Agora, supondo que $n \in A$, provemos que $s(n) \in A$. Se $n \in \Dom \varphi$, então existe $x \in X$ tal que $(n,x) \in \varphi$. De $(n,x) \in \varphi$, vem $(s(n),f(x)) \in \varphi$. Pela hipótese de indução, se $(n,y) \in \varphi$, então $y=x$; provemos que, se $(s(n),y) \in \varphi$, então $y=f(x)$. Suponha, por absurdo, que existe $y \in X \setminus \{f(x) \}$ tal que $(s(n),y) \in \varphi$. Como $s(n) \neq a$, temos $(a,x_0) \in \varphi \setminus \{(s(n),y)\}$. Se $(m,z) \in \varphi \setminus \{(s(n),y)\}$, então $(s(m),f(z)) \in \varphi$. Se $s(m) \neq s(n)$, então $(s(m), f(z)) \in \varphi \setminus \{(s(n),y)\}$. Se $s(m) = s(n)$, então, pela injetividade de $s$, temos $m=n$, e pela hipótese de indução, temos $z=x$, de modo que $f(z) = f(x) \neq y$, o que implica $(s(m), f(z)) \in \varphi \setminus \{(s(n),y)\}$. Em qualquer caso, temos $\varphi \setminus \{(s(n),y)\} \in \mathcal{C}$, uma contradição (pelo mesmo motivo visto na base de indução). Logo, se $(s(n),y) \in \varphi$, então $y=f(x)$, isto é, $s(n) \in A$. Logo, pelo princípio de indução, temos $A = N$, de modo que $\varphi$ é uma função. \whiteproof
        \end{midproof}
    Como $\Dom \varphi = N$ e $\varphi$ é uma função, existe $\varphi: N \to X$ com as propriedades desejadas. Provemos, por fim, sua unicidade.
        \begin{claim}
            A função $\varphi : N \to X$ é única.
        \end{claim}
        \begin{midproof}
            Por indução. Seja $\Phi : N \to X$ tal que $\Phi(a) = x_0$ e $\Phi(s(n)) = f(\Phi(n))$ para todo $n \in N$, e seja $A := \{ n \in N : \Phi(n) = \varphi(n)\}$. Como $\Phi(a) = \varphi(a)$, temos $a \in A$. Agora, supondo que $n \in A$, provemos que $s(n) \in A$. Se $n \in A$, então, pela hipótese de indução, temos $\Phi(n) = \varphi(n)$, de modo que
                \[
                    \Phi(s(n)) = f(\Phi(n)) = f(\varphi(n)) = \varphi(s(n)),
                \]
            isto é, $s(n) \in A$. Logo, pelo princípio de indução, temos $A = N$. \whiteproof
        \end{midproof}
    Logo, existe uma única função $\varphi: N \to X$ tal que $\varphi(a) = x_0$ e $\varphi(s(n)) = f(\varphi(n))$ para todo $n \in N$. \blackproof
\end{proof}

\begin{theorem} \label{teo.set:isomorfismo_sistemas_peano} \fonte[52]{moschovakis}
    Sejam $(N, s, a)$ e $(M, t, b)$ sistemas de Peano. Existe uma única bijeção $\varphi : N \to M$ tal que
        \begin{itemize}
            \item $\varphi(a) = b$.
            \item $\varphi(s(n)) = t(\varphi(n))$ para todo $n \in N$.
        \end{itemize}
\end{theorem}

\begin{proof} \fonte[52]{moschovakis}
    Pelo teorema da recursão, existe uma única função $\varphi : N \to M$ tal que $\varphi(a) = b$ e $\varphi(s(n)) = t(\varphi(n))$ para todo $n \in N$. Resta provar que $\varphi$ é uma bijeção.
        \begin{claim}
            $\Img \varphi = M$.
        \end{claim}
        \begin{midproof}
            Por indução. Temos $\Img \varphi \subseteq M$. Como $\varphi(a) = b$, temos $b \in \Img \varphi$. Agora, supondo que $m \in \Img \varphi$, provemos que $t(m) \in \Img \varphi$. Se $m \in \Img \varphi$, então existe $n \in N$ tal que $\varphi(n) = m$. Logo, $\varphi(s(n)) = t(\varphi(n)) = t(m)$, de modo que $t(m) \in \Img \varphi$. Pelo princípio de indução, temos $\Img \varphi = M$. \whiteproof
        \end{midproof}
        \begin{claim}
            $\varphi$ é injetiva.
        \end{claim}
        \begin{midproof}
            Por indução. Seja
                \[
                    A := \{n \in N : \forall m (m \in N \rightarrow (\varphi(m) = \varphi(n) \rightarrow m=n)) \}.
                \]
            Seja $m \in N$ tal que $\varphi(m) = \varphi(a) = b$ e suponha, por absurdo, que $m \neq a$. Então existe $k \in N$ tal que $m = s(k)$, de modo que
                \[
                    b = \varphi(m) = \varphi(s(k)) = t(\varphi(k)) \in t(M),
                \]
            o que contradiz $b \notin t[M]$. Logo, $m=a$, isto é, $a \in A$. Agora, supondo que $n \in A$, provemos que $s(n) \in A$. Se $n \in A$, então, pela hipótese de indução, $\varphi(m) = \varphi(n)$ implica $m=n$; provemos que $\varphi(m) = \varphi(s(n))$ implica $m = s(n)$. Se $m=a$, então
                \[
                    b = \varphi(a) = \varphi(s(n)) = t(\varphi(n)) \in t(M),
                \]
            o que contradiz $b \notin t[M]$. Logo, $m \neq a$, e então existe $k \in N$ tal que $m = s(k)$. Com isso,
                \[
                    t(\varphi(k)) = \varphi(s(k)) = \varphi(m) = \varphi(s(n)) = t(\varphi(n)),
                \]
            e então $\varphi(k) = \varphi(n)$ pela injetividade de $t$. Pela hipótese de indução, temos $k=n$, de modo que $m = s(k) = s(n)$, isto é, $s(n) \in A$. Pelo princípio de indução, temos $A = N$. Logo, $\varphi$ é injetiva. \whiteproof
        \end{midproof}
    Logo, $\varphi : N \to M$ é uma bijeção. \blackproof
\end{proof}

\subsection{Aritmética em Sistemas de Peano}

\begin{theorem}[Recursão com parâmetro] \label{teo.set:recpar1} \fonte[47]{jechhrbacek1999}
    Seja $(N,s,a)$ um sistema de Peano. Sejam $X \neq \emptyset$, $x_0 \in X$ e $f : X \times N \to X$. Existe uma única função $\varphi : N \to X$ com as seguintes propriedades.
        \begin{itemize}
            \item $\varphi(a) = x_0$.
            \item $\varphi(s(n)) = f(\varphi(n), n)$, para todo $n \in N$.
        \end{itemize}
\end{theorem}

\begin{proof} \fonte[115]{fajardo2024conjuntos}
    Defina $g : X \times N \to X \times N$ por $g(y,n) = (f(y,n),s(n))$. Pelo \cref{teo.set:recursao_peano}, existe uma única função $\psi : N \to X \times N$ tal que $\psi(a) = (x_0,a)$ e $\psi(s(n)) = g(\psi(n))$ para todo $n \in N$.
        \begin{claim}
            $\pi_2(\psi(n)) = n$ para todo $n \in N$.
        \end{claim}
        \begin{midproof}
            Por indução. Seja $A := \{n \in N: \pi_2(\psi(n)) = n\}$. Como $\pi_2(\psi(a)) = \pi_2 (x_0,a) = a$, temos $a \in A$. Agora, supondo que $n \in A$, provemos que $s(n) \in A$. Se $n \in A$, então, pela hipótese de indução, temos $\pi_2(\psi(n)) = n$, de modo que
                \begin{align*}
                    \pi_2(\psi(s(n))) &= \pi_2(g(\psi(n))) \\
                    &= \pi_2(g(\pi_1(\psi(n)), \pi_2(\psi(n)))) \\
                    &= \pi_2 (g(\pi_1(\psi(n)), n)) \\
                    &= \pi_2 (f(\pi_1(\psi(n)),n),s(n)) \\
                    &= s(n),
                \end{align*}
            isto é, $s(n) \in A$. Logo, $\pi_2(\psi(n)) = n$ para todo $n \in N$. \whiteproof
        \end{midproof}
    Agora, defina $\varphi := \pi_1 \circ \psi : N \to X$, isto é, $\varphi(n) = \pi_1(\psi(n))$ para todo $n \in N$. Essa é a função procurada. Provemos isso.
        \begin{claim}
            $\varphi(a) = x_0$ e, para todo $n \in N$, $\varphi(s(n)) = f(\varphi(n), n)$.
        \end{claim}
        \begin{midproof}
            Como $\psi(a) = (x_0,a)$, temos $\varphi(a) = \pi_1(x_0, a) = x_0$. Como $\psi(n) = (\varphi(n),n)$, para todo $n \in N$, temos
                \[
                    \psi(s(n)) = g(\psi(n)) = g(\varphi(n),n) = (f(\varphi(n),n),s(n)),
                \]
            para todo $n \in N$. Logo,
                \[
                    \varphi(s(n)) = \pi_1(\psi(s(n))) = \pi_1(f(\varphi(n),n),s(n)) = f(\varphi(n),n)
                \]
            para todo $n \in N$. \whiteproof
        \end{midproof}
    Logo, existe $\varphi : N \to X$ com as propriedades desejadas. Provemos, por fim, sua unicidade.
        \begin{claim}
            A função $\varphi: N \to X$ é única.
        \end{claim}
        \begin{midproof}
            Por indução. Seja $\Phi : N \to X$ tal que $\Phi(a) = x_0$ e $\Phi(s(n)) = f(\Phi(n), n)$ para todo $n \in N$, e seja $A := \{ n \in N : \Phi(n) = \varphi(n)\}$. Como $\Phi(a) = \varphi(a)$, temos $a \in A$. Agora, supondo que $n \in A$, provemos que $s(n) \in A$. Se $n \in A$, então, pela hipótese de indução, temos $\Phi(n) = \varphi(n)$, de modo que
                \[
                    \Phi(s(n)) = f(\Phi(n), n) = f(\varphi(n), n) = \varphi(s(n)),
                \]
            isto é, $s(n) \in A$. Logo, pelo princípio de indução, temos $A = N$. \whiteproof
        \end{midproof}
    Logo, existe uma única função $\varphi: N \to X$ tal que $\varphi(a) = x_0$ e $\varphi(s(n)) = f(\varphi(n), n)$ para todo $n \in N$. \blackproof
\end{proof}

\begin{theorem} \label{lema.set:recpar2} \fonte[51]{jechhrbacek1999}
    Seja $(N,s,a)$ um sistema de Peano. Sejam $f : X \to Y$ e $g : X \times Y \times N \to Y$. Existe uma única função $\varphi : X \times N \to Y$ com as seguintes propriedades.
        \begin{itemize}
            \item $\varphi(x,a) = f(x)$ para todo $x \in X$.
            \item $\varphi(x, s(n)) = g(x,\varphi(x,n),n)$ para quaisquer $x \in X$ e $n \in N$.
        \end{itemize}
\end{theorem}

\begin{proof} \fonte[51]{jechhrbacek1999}
    Defina $G : Y^X \times N \to Y^X$ de modo que $G(h,n)(x) := g(x,h(x),n)$ para quaisquer $h \in Y^X$, $n \in N$ e $x \in X$.
        \begin{claim}
            A função $G$ está bem definida.
        \end{claim}
        \begin{midproof}
            Defina
                \[
                    G := \{((h,n),F) \in (Y^X \times N) \times Y^X : F = \{(x,y)\in X\times Y : y=g(x,h(x),n)\}\}.
                \]
            Sejam $h \in Y^X$ e $n \in N$. Para todo $x \in X$, como $h: X \to Y$, temos $h(x) \in Y$, de modo que $(x,h(x),n) \in X \times Y \times N$. Com isso, existe um único $y \in Y$ tal que $y = g(x,h(x),n)$, donde
                \[
                    \{(x,y)\in X\times Y : y=g(x,h(x),n)\} \in Y^X.
                \]
            Logo, $G$ é uma função de $Y^X \times N$ em $Y^X$. \whiteproof
        \end{midproof}
    Pelo \cref{teo.set:recpar1} da recursão, defina $\psi : N \to Y^X$ por $\psi(a) = f$ e $\psi(s(n)) = G(\psi(n),n)$ para todo $n \in N$. Por fim, defina $\varphi : X \times N \to Y$ por $\varphi(x,n) := \psi(n)(x)$, para quaisquer $x \in X$ e $n \in N$. Para todo $x \in X$, temos $\varphi(x,a) = \psi(a)(x) = f(x)$. Para todo $x \in X$ e todo $n \in N$, temos
        \begin{align*}
            \varphi(x,s(n)) &= \psi(s(n))(x) \\
            &= G(\psi(n),n)(x) \\
            &= g(x, \psi(n)(x), n) \\
            &= g(x, \varphi(x,n), n).
        \end{align*}
    Logo, existe $\varphi : X \times N \to Y$ com as propriedades desejadas. Provemos, por fim, sua unicidade.
        \begin{claim}
            A função $\varphi : X \times N \to Y$ é única.
        \end{claim}
        \begin{midproof}
            Por indução. Seja $\Phi : X \times N \to Y$ tal que $\Phi(x,a) = f(x)$ para todo $x \in X$ e $\Phi(x,s(n)) = g(x,\Phi(x,n), n)$ para quaisquer $x \in X$ e $n \in N$, e seja $A := \{ n \in N : \forall x (x \in X \rightarrow \Phi(x,n) = \varphi(x,n))\}$. Como $\Phi(x,a) = \varphi(x,a)$ para todo $x \in X$, temos $a \in A$. Agora, supondo que $n \in A$, provemos que $s(n) \in A$. Se $n \in A$, então, pela hipótese de indução, temos $\Phi(x,n) = \varphi(x,n)$ para todo $x \in X$, de modo que
                \[
                    \Phi(x,s(n)) = g(x,\Phi(x,n), n) = g(x,\varphi(x,n), n) = \varphi(x,s(n)),
                \]
            para todo $x \in X$, isto é, $s(n) \in A$. Logo, pelo princípio de indução, temos $A = N$. \whiteproof
        \end{midproof}
    Logo, existe uma única função $\varphi: X \times N \to Y$ tal que $\varphi(x,a) = f(x)$ para todo $x \in X$ e $\varphi(x,s(n)) = g(x,\varphi(x,n), n)$ para quaisquer $x \in X$ e $n \in N$. \blackproof
\end{proof}

\begin{comment}
\begin{proof}
    A ideia é considerar todas as relações de $X \times N$ em $Y$ que têm as propriedades desejadas e provar que a interseção de todas elas resulta em uma única função de domínio $X \times N$. Defina
        \[
            \mathcal{C} := \left\{ R \subseteq (X \times N) \times Y : ((x,a),f(x)) \in R \land \forall x \forall n \forall y ( ((x,n),y) \in R \rightarrow ((x,s(n)), g(x,y,n)) \in R) \right\}.
        \]
\end{proof}
\end{comment}

\subsubsection{Adição}

\begin{corollary}
    Seja $(N,s,a)$ um sistema de Peano. Existe uma única função $+_N : N \times N \to N$ com as seguintes propriedades.
        \begin{itemize}
            \item $+_N (m,a) = m$ para todo $m \in N$.
            \item $+_N (m,s(n)) =s(+_N (m,n))$ para quaisquer $m,n \in N$.
        \end{itemize}
\end{corollary}

\begin{proof} \fonte[52]{jechhrbacek1999}
    Defina $f : N \to N$ por $f(m) = m$ para todo $m \in N$, e $g : N \times N \times N \to N$ por $g(m,x,n) = s(x)$ para quaisquer $m,x,n \in N$. Pelo \cref{lema.set:recpar2}, existe uma única função $+_N : N \times N \to N$ tal que $+_N(m,a) = f(m) = m$ para todo $m \in N$, e $+_N(m,s(n)) = g(m,+_N(m,n),n) = s(+_N(m,n))$ para quaisquer $m,n \in N$. \blackproof
\end{proof}

\begin{lemma} \label{lema.set:comutatividade_soma_naturais}
    Seja $(N,s,a)$ um sistema de Peano.
        \begin{enumerate}[leftmargin=*, align=left, label=\textbf{(\alph*)}]
            \item $+_N(a,n)=n$, para todo $n \in N$.
            \item $+_N(s(m),n) = +_N(m,s(n))$, para quaisquer $m,n \in N$.
        \end{enumerate}
\end{lemma}

\begin{proof}
    \leavevmode
        \begin{enumerate}[leftmargin=*, align=left, label=\textbf{(\alph*)}]
            \item Defina $A := \{n \in N : +_N(a,n)=n\}$. Como $+_N (a,a) = a$, temos $a \in A$. Agora, supondo que $n \in A$, provemos que $s(n) \in A$. Se $n \in A$, então, pela hipótese de indução, temos $+_N (a, n) = n$, de modo que
                \[
                    +_N (a, s(n)) = s(+_N (a, n)) = s(n),
                \]
            isto é, $s(n) \in A$. Logo, pelo princípio de indução, temos $A = N$. Logo, $a +_N n = n$, para todo $n \in N$. \blackproof
            \item Fixe $m \in N$ e defina
                \[
                    A := \{n \in N : +_N(s(m),n) = +_N(m,s(n)) \}.    
                \]
            Como $+_N (s(m),a) = s(m)$ e $+_N (m,s(a)) = s(+_N (m,a)) = s(m)$, temos $+_N (s(m),a) = +_N (m,s(a))$, de modo que $a \in A$. Agora, supondo que $n \in A$, provemos que $s(n) \in A$. Se $n \in A$, então, pela hipótese de indução, temos $+_N (s(m),n) = +_N (m,s(n))$, de modo que
                \begin{align*}
                    +_N (s(m), s(n)) &= s(+_N (s(m), n)) \\
                    &= s(+_N (m, s(n))) \\
                    &= +_N (m, s(s(n))),
                \end{align*}
            isto é, $s(n) \in A$. Logo, pelo princípio de indução, temos $A = N$. Logo, $s(m) +_N n = m +_N s(n)$, para quaisquer $m,n \in N$. \blackproof
        \end{enumerate}
\end{proof}

\begin{theorem}
    Seja $(N,s,a)$ um sistema de Peano.
        \begin{enumerate}[leftmargin=*, align=left, label=\textbf{(\alph*)}]
            \item $m +_N n = n +_N m$, para quaisquer $m,n \in N$.
            \item $m +_N (n +_N p)=(m +_N n) +_N p$, para quaisquer $m,n,p \in N$.
        \end{enumerate}
\end{theorem}

\begin{proof}
    \leavevmode
        \begin{enumerate}[leftmargin=*, align=left, label=\textbf{(\alph*)}]
            \item Fixe $m \in N$ e defina
                \[
                    A := \{n \in N : m +_N n = n +_N m\}.
                \]
            Pelo \cref{lema.set:comutatividade_soma_naturais}, temos $a +_N m = m$, e como $m +_N a = m$, temos $m +_N a = a +_N m$, de modo que $a \in A$. Agora, supondo que $n \in A$, provemos que $s(n) \in A$. Se $n \in A$, então, pela hipótese de indução, temos $m +_N n = n +_N m$, de modo que
                \[
                    m +_N s(n) = s(m +_N n) = s(n +_N m) = n +_N s(m);
                \]
            como, pelo \cref{lema.set:comutatividade_soma_naturais}, $n +_N s(m) = s(n) +_N m$, temos $m +_N s(n) = s(n) +_N m$, de modo que $s(n) \in A$. Logo, pelo princípio de indução, temos $A = N$. Logo, a adição $+_N$ é comutativa. \blackproof
            \item Fixe $m,n \in N$ e defina
                \[
                    A := \{p \in N : m +_N (n +_N p) = (m +_N n) +_N p\}.    
                \]
            Como $n +_N a = n$, temos $m +_N (n +_N a) = m +_N n$, e como $(m +_N n) +_N a = m +_N n$, temos $m +_N (n +_N a) = (m +_N n) +_N a$, de modo que $a \in A$. Agora, supondo que $p \in A$, provemos que $s(p) \in A$. Se $p \in A$, então, pela hipótese de indução, temos $m +_N (n +_N p) = (m +_N n) +_N p$, de modo que
                \begin{align*}
                    m +_N (n +_N s(p)) &= m +_N s(n +_N p) \\
                    &= s(m +_N (n +_N p)) \\
                    &= s((m +_N n) +_N p) \\
                    &= (m +_N n) +_N s(p),
                \end{align*}
            isto é, $s(p) \in A$. Logo, pelo princípio de indução, temos $A = N$. Logo, a adição $+_N$ é associativa. \blackproof
        \end{enumerate}
\end{proof}

\begin{definition}
    Seja $(N,s,a)$ um sistema de Peano. Definimos $1_N := s(a)$.
\end{definition}

\begin{proposition}
    Seja $(N,s,a)$ um sistema de Peano. Para todo $n \in N$, temos
        \[
            s(n) = n +_N 1_N.
        \]
\end{proposition}

\begin{proof}
    Para todo $n \in N$, temos
        \[
            n +_N 1_N = n +_N s(a) = s(n +_N a) = s(n),
        \]
    isto é, $n +_N 1_N = s(n)$. \blackproof
\end{proof}

\begin{proposition}
    Sejam $(N,s,a)$ um sistema de Peano e $m,n,p \in N$.
        \begin{enumerate}[leftmargin=*, align=left, label=\textbf{(\alph*)}]
            \item Se $m +_N n = a$, então $m = n = a$.
            \item Se $m +_N p = n +_N p$, então $m=n$.
        \end{enumerate}
\end{proposition}

\begin{proof}
    \leavevmode
        \begin{enumerate}[leftmargin=*, align=left, label=\textbf{(\alph*)}]
            \item Sejam $m,n \in N$ tais que $m +_N n = a$. Se $n \neq a$, então existiria $k \in N$ tal que $n=s(k)$, donde
                \[
                    m +_N n = m +_N s(k) = s(m +_N k) = a,
                \]
            uma contradição pois $a \notin s[N]$. Logo, $n=a$. Com isso, como $m +_N n = a$ e $m +_N a = m$, temos $m = a$. \blackproof
            \item Fixe $m,n \in N$ e defina
                \[
                    A := \{p \in N : m +_N p = n +_N p \Rightarrow m=n\}.    
                \]
            Como $m +_N a = m$ e $n +_N a = n$, se $m +_N a = n +_N a$, temos $m = n$, de modo que $a \in A$. Agora, supondo que $p \in A$, provemos que $s(p) \in A$. Se $m +_N s(p) = n +_N s(p)$, então $s(m +_N p) = s(n +_N p)$, e como $s$ é injetiva, temos $m +_N p = n +_N p$; com isso, pela hipótese de indução, temos $m = n$, de modo que $s(p) \in A$. Logo, pelo princípio de indução, temos $A = N$. Logo, se $m +_N p = n +_N p$, então $m=n$. \blackproof
        \end{enumerate}
\end{proof}

\subsubsection{Multiplicação}

\begin{corollary}
    Seja $(N,s,a)$ um sistema de Peano. Existe uma única função $\cdot_N : N \times N \to N$ com as seguintes propriedades.
        \begin{itemize}
            \item $\cdot_N (m,a) = a$ para todo $m \in N$.
            \item $\cdot_N (m,s(n)) = +_N (\cdot_N (m,n),m)$ para quaisquer $m,n \in N$.
        \end{itemize}
\end{corollary}

\begin{proof} \fonte[52]{jechhrbacek1999}
    Defina $f : N \to N$ por $f(m) = a$ para todo $m \in N$, e $g : N \times N \times N \to N$ por $g(m,x,n) = +_N(x,m)$ para quaisquer $m,x,n \in N$. Pelo \cref{lema.set:recpar2}, existe uma única função $\cdot_N : N \times N \to N$ tal que $\cdot_N(m,a) = f(m) = a$ para todo $m \in N$, e $\cdot_N(m,s(n)) = g(m,\cdot_N(m,n),n) = +_N(\cdot_N(m,n),m)$ para quaisquer $m,n \in N$. \blackproof
\end{proof}

\begin{lemma} \label{lema.set:comutatividade_produto_naturais}
    Seja $(N,s,a)$ um sistema de Peano.
        \begin{enumerate}[leftmargin=*, align=left, label=\textbf{(\alph*)}]
            \item $\cdot_N (a, n) = a$, para todo $n \in N$.
            \item $\cdot_N (s(m),n) = +_N (\cdot_N (m,n), n)$, para quaisquer $m,n \in N$.
        \end{enumerate}
\end{lemma}

\begin{proof}
    \leavevmode
        \begin{enumerate}[leftmargin=*, align=left, label=\textbf{(\alph*)}]
            \item Defina $A := \{n \in N : \cdot_N (a,n) = a\}$. Como $\cdot_N (a,a) = a$, temos $a \in A$. Agora, supondo que $n \in A$, provemos que $s(n) \in A$. Se $n \in A$, então, pela hipótese de indução, temos $\cdot_N (a,n) = a$, de modo que
                \[
                    \cdot_N (a, s(n)) = +_N (\cdot_N (a, n), a) = +_N (a, a) = a,
                \]
            isto é, $s(n) \in A$. Logo, pelo princípio de indução, temos $A = N$. Logo, $a \cdot_N n = a$, para todo $n \in N$. \blackproof
            \item Fixe $m \in N$ e defina
                \[
                    A := \{n \in N : \cdot_N (s(m),n) = +_N (\cdot_N (m,n), n)\}.
                \]
            Como $\cdot_N (s(m),a) = a$ e $+_N (\cdot_N (m,a), a) = +_N (a, a) = a$, temos $\cdot_N (s(m),a) = +_N (\cdot_N (m,a), a)$, de modo que $a \in A$. Agora, supondo que $n \in A$, provemos que $s(n) \in A$. Se $n \in A$, então, pela hipótese de indução, temos $\cdot_N (s(m),n) = +_N (\cdot_N (m,n), n)$, de modo que
                \begin{align*}
                    \cdot_N (s(m), s(n)) &= +_N (\cdot_N (s(m), n), s(m)) \\
                    &= +_N (+_N (\cdot_N (m, n), n), s(m)) \\
                    &= +_N (\cdot_N (m, n), +_N (n, s(m))) \\
                    &= +_N (\cdot_N (m, n), +_N (s(m), n)) \\
                    &= +_N (\cdot_N (m, n), +_N (m, s(n))) \\
                    &= +_N (+_N (\cdot_N (m, n), m), s(n)) \\
                    &= +_N (\cdot_N (m, s(n)), s(n)),
                \end{align*}
            isto é, $s(n) \in A$. Logo, pelo princípio de indução, temos $A = N$. Logo, $s(m) \cdot_N n = (m \cdot_N n) +_N n$, para quaisquer $m,n \in N$. \blackproof
        \end{enumerate}
\end{proof}

\begin{proposition}
    Seja $(N,s,a)$ um sistema de Peano. Para quaisquer $m,n,p \in N$, temos
        \[
            m \cdot_N (n +_N p) = m \cdot_N n +_N m \cdot_N p.
        \] 
\end{proposition}

\begin{proof}
    Fixe $m,n \in N$ e defina
        \[
            A := \{p \in N : m \cdot_N (n +_N p) = (m \cdot_N n) +_N (m \cdot_N p)\}.
        \]
    Como $n +_N a = n$, temos $m \cdot_N (n +_N a) = m \cdot_N n$, e como $m \cdot_N a = a$, temos
        \[
            (m \cdot_N n) +_N (m \cdot_N a) = (m \cdot_N n) +_N a = m \cdot_N n,
        \]
    de modo que, $m \cdot_N (n +_N a) = (m \cdot_N n) +_N (m \cdot_N a)$, isto é, $a \in A$. Agora, supondo que $p \in A$, provemos que $s(p) \in A$. Se $p \in A$, então, pela hipótese de indução, temos $m \cdot_N (n +_N p) = (m \cdot_N n) +_N (m \cdot_N p)$, de modo que
        \begin{align*}
            m \cdot_N (n +_N s(p)) &= m \cdot_N s(n +_N p) \\
            &= (m \cdot_N (n +_N p)) +_N m \\
            &= ((m \cdot_N n) +_N (m \cdot_N p)) +_N m \\
            &= (m \cdot_N n) +_N ((m \cdot_N p) +_N m) \\
            &= (m \cdot_N n) +_N (m \cdot_N s(p)),
        \end{align*}
    isto é, $s(p) \in A$. Logo, pelo princípio de indução, temos $A = N$. Logo, $m \cdot_N (n +_N p) = (m \cdot_N n) +_N (m \cdot_N p)$, para quaisquer $m,n,p \in N$. \blackproof
\end{proof}

\begin{theorem}
    Seja $(N,s,a)$ um sistema de Peano.
        \begin{enumerate}[leftmargin=*, align=left, label=\textbf{(\alph*)}]
            \item $m \cdot_N n = n \cdot_N m$, para quaisquer $m,n \in N$.
            \item $m \cdot_N (n \cdot_N p)= (m \cdot_N n) \cdot_N p$, para quaisquer $m,n,p \in N$.
        \end{enumerate}
\end{theorem}

\begin{proof}
    \leavevmode
        \begin{enumerate}[leftmargin=*, align=left, label=\textbf{(\alph*)}]
            \item Fixe $m \in N$ e defina
                \[
                    A := \{n \in N : m \cdot_N n = n \cdot_N m \}.
                \]
            Pelo \cref{lema.set:comutatividade_produto_naturais}, temos $a \cdot_N m = a$, e como $m \cdot_N a = a$, temos $m \cdot_N a = a \cdot_N m$, de modo que $a \in A$. Agora, supondo que $n \in A$, provemos que $s(n) \in A$. Se $n \in A$, então, pela hipótese de indução, temos $m \cdot_N n = n \cdot_N m$, de modo que
                \[
                    m \cdot_N s(n) = (m \cdot_N n) +_N m = (n \cdot_N m) +_N m;
                \]
            como, pelo \cref{lema.set:comutatividade_produto_naturais}, $(n \cdot_N m) +_N m = s(n) \cdot_N m$, temos $m \cdot_N s(n) = s(n) \cdot_N m$, de modo que $s(n) \in A$. Logo, pelo princípio de indução, temos $A = N$. Logo, a multiplicação $\cdot_N$ é comutativa. \blackproof
            \item Fixe $m,n \in N$ e defina
                \[
                    A := \{p \in N : m \cdot_N (n \cdot_N p) = (m \cdot_N n) \cdot_N p\}.
                \]
            Como $n \cdot_N a = a$, temos $m \cdot_N (n \cdot_N a) = m \cdot_N a = a$, e como $(m \cdot_N n) \cdot_N a = a$, temos $m \cdot_N (n \cdot_N a) = (m \cdot_N n) \cdot_N a$, de modo que $a \in A$. Agora, supondo que $p \in A$, provemos que $s(p) \in A$. Se $p \in A$, então, pela hipótese de indução, temos $m \cdot_N (n \cdot_N p) = (m \cdot_N n) \cdot_N p$, de modo que
                \begin{align*}
                    m \cdot_N (n \cdot_N s(p)) &= m \cdot_N ((n \cdot_N p) +_N n) \\
                    &= (m \cdot_N (n \cdot_N p)) +_N (m \cdot_N n) \\
                    &= ((m \cdot_N n) \cdot_N p) +_N (m \cdot_N n) \\
                    &= (m \cdot_N n) \cdot_N s(p),
                \end{align*}
            isto é, $s(p) \in A$. Logo, pelo princípio de indução, temos $A = N$. Logo, a multiplicação $\cdot_N$ é associativa. \blackproof
        \end{enumerate}
\end{proof}

\begin{proposition}
    Seja $(N,s,a)$ um sistema de Peano. Para todo $n \in N$, temos
        \[
            n \cdot_N 1_N = n.
        \]
\end{proposition}

\begin{proof}
    Para todo $n \in N$, temos
        \[
            n \cdot_N 1_N = n \cdot_N s(a) = (n \cdot_N a) +_N n = a +_N n = n,
        \]
    isto é, $n \cdot_N 1_N = n$. \blackproof
\end{proof}

\subsubsection{Ordem}

\begin{proposition}
    Seja $(N,s,a)$ um sistema de Peano.
        \begin{enumerate}[leftmargin=*, align=left, label=\textbf{(\alph*)}]
            \item A relação
                \[
                    {\leq_N} := \{(m,n) \in N \times N : \exists p (p \in N \land m +_N p=n) \}
                \]
            é de ordem parcial em $N$.
            \item A relação
                \[
                    {<_N} := \{(m,n) \in N \times N : \exists p (p \in N_{\neq a} \land m +_N p=n) \}
                \]
            é de ordem estrita em $N$.
            \item 
                \begin{enumerate}[label=\roman*.]
                    \item ${<_N} = \{(m,n) \in N \times N : m \leq_N n \land m \neq n\}$.
                    \item ${\leq_N} = \{(m,n) \in N \times N : m <_N n \lor m = n\}$.
                \end{enumerate}
        \end{enumerate}
\end{proposition}

\begin{proof}
    \leavevmode
        \begin{enumerate}[leftmargin=*, align=left, label=\textbf{(\alph*)}]
            \item 
                \begin{claim}
                    ${\leq_N}$ é reflexiva.
                \end{claim}
                \begin{midproof}
                    Seja $m \in N$. Como $m +_N a = m$, temos $m \leq_N m$. Logo, ${\leq_N}$ é reflexiva. \whiteproof
                \end{midproof}
                \begin{claim}
                    ${\leq_N}$ é antissimétrica.
                \end{claim}
                \begin{midproof}
                    Sejam $m,n \in N$ tais que $m \leq_N n$ e $n \leq_N m$. Então, existem $p,q \in N$ tais que $m +_N p = n$ e $n +_N q = m$. Com isso,
                        \[
                            m +_N (p +_N q) = (m +_N p) +_N q = n +_N q = m = m +_N a,
                        \]
                    de modo que $p +_N q = a$, isto é, $p=q=a$. Logo, como $m +_N p = n$, temos $m=n$. Logo, ${\leq_N}$ é antissimétrica. \whiteproof
                \end{midproof}
                \begin{claim}
                    ${\leq_N}$ é transitiva.
                \end{claim}
                \begin{midproof}
                    Sejam $m,n,r \in N$ tais que $m \leq_N n$ e $n \leq_N r$. Então, existem $p,q \in N$ tais que $m +_N p = n$ e $n +_N q = r$. Com isso,
                        \[
                            m +_N (p +_N q) = (m +_N p) +_N q = n +_N q = r,
                        \]
                    de modo que $m \leq_N r$. Logo, ${\leq_N}$ é transitiva. \whiteproof
                \end{midproof}
            Logo, ${\leq_N}$ é uma relação de ordem parcial em $N$. \blackproof
            \item 
                \begin{claim}
                    ${<_N}$ é irreflexiva.
                \end{claim}
                \begin{midproof}
                    Seja $m \in N$. Se $m <_N m$, então existiria $p \in N_{\neq a}$ tal que $m +_N p = m$; como $m=m+_N a$, teríamos $m +_N p = m +_N a$, donde $p=a$, uma contradição. Logo, não ocorre $m <_N m$, para todo $m \in N$. Logo, ${<_N}$ é irreflexiva. \whiteproof
                \end{midproof}
                \begin{claim}
                    ${<_N}$ é transitiva.
                \end{claim}
                \begin{midproof}
                    Sejam $m,n,r \in N$ tais que $m <_N n$ e $n <_N r$. Então, existem $p,q \in N_{\neq a}$ tais que $m +_N p = n$ e $n +_N q = r$. Com isso,
                        \[
                            m +_N (p +_N q) = (m +_N p) +_N q = n +_N q = r,
                        \]
                    e como $p +_N q \neq a$ (pois, se fosse $p +_N q = a$, teríamos, $p=q=a$, uma contradição), existe $p +_N q \in N_{\neq a}$ tal que $m +_N (p +_N q) = r$, de modo que $m <_N r$. Logo, ${<_N}$ é transitiva. \whiteproof
                \end{midproof}
            Logo, ${<_N}$ é uma relação de ordem estrita em $N$. \blackproof
            \item 
                \begin{enumerate}[label=\roman*.]
                    \item Por um lado ($\subseteq$), se $(m,n) \in {<_N}$, então existe $p \in N_{\neq a}$ tal que $m +_N p = n$. Logo, existe $p \in N$ tal que $m +_N p=n$, de modo que $m \leq_N n$. Além disso, se fosse $m=n$, teríamos $m +_N p = m = m +_N a$, donde $p=a$, uma contradição. Logo, $m \neq n$, o que prova a inclusão $\subseteq$. Por outro lado ($\supseteq$), se $(m,n) \in N \times N$ é tal que $m \leq_N n$ e $m \neq n$, então existe $p \in N$ tal que $m +_N p=n$. Se fosse $p=a$, teríamos $n=m+_N a=m$, uma contradição. Logo, $p \in N_{\neq a}$, de modo que $m <_N n$, o que prova a inclusão $\supseteq$.
                    \item Por um lado ($\subseteq$), se $(m,n) \in {\leq_N}$, então existe $p \in N$ tal que $m +_N p=n$. Se $p=a$, então $n=m+_N a=m$. Se $p \neq a$, então $p \in N_{\neq a}$, de modo que $m <_N n$. Logo, $m <_N n$ ou $m=n$, o que prova a inclusão $\subseteq$. Por outro lado ($\supseteq$), se $(m,n) \in N \times N$ é tal que $m <_N n$ ou $m=n$, então, se $m <_N n$, existe $p \in N_{\neq a}$ tal que $m +_N p=n$; em particular, existe $p \in N$ tal que $m +_N p=n$, de modo que $m \leq_N n$. Se $m=n$, então $m+_N a=n$, de modo que $m \leq_N n$. Logo, $m \leq_N n$, o que prova a inclusão $\supseteq$.
                \end{enumerate}
            Logo, ${<_N}$ e ${\leq_N}$ se comportam como na \cref{prop.set:parcial_e_estrita}. \blackproof
        \end{enumerate}
\end{proof}

\begin{proposition}
    Sejam $(N,s,a)$ um sistema de Peano e $m,n,p \in N$.
        \begin{enumerate}[leftmargin=*, align=left, label=\textbf{(\alph*)}]
            \item Se $m \leq_N n$, então
                \begin{enumerate}[label=\roman*.]
                    \item $m +_N p \leq_N n +_N p$.
                    \item $m \cdot_N p \leq_N n \cdot_N p$.
                \end{enumerate}
            \item Se $m <_N n$, então
                \begin{enumerate}[label=\roman*.]
                    \item $m +_N p <_N n +_N p$.
                    \item $m \cdot_N p <_N n \cdot_N p$, desde que $p \neq a$.
                \end{enumerate}
        \end{enumerate}
\end{proposition}

\begin{proof}
    \leavevmode
        \begin{enumerate}[leftmargin=*, align=left, label=\textbf{(\alph*)}]
            \item Suponha que $m \leq_N n$. Então, existe $q \in N$ tal que $m +_N q=n$.
                \begin{enumerate}[label=\roman*.]
                    \item Temos
                        \[
                            (m +_N p) +_N q = (m +_N q) +_N p = n +_N p,
                        \]
                    de modo que existe $q \in N$ tal que $(m +_N p) +_N q = n +_N p$. Logo, $m +_N p \leq_N n +_N p$.

                    \item Temos
                        \begin{align*}
                            (m \cdot_N p) +_N (q \cdot_N p) &= (p \cdot_N m) +_N (p \cdot_N q) \\
                            &= p \cdot_N (m +_N q) \\
                            &= p \cdot_N n \\
                            &= n \cdot_N p,
                        \end{align*}
                    de modo que existe $q \cdot_N p \in N$ tal que $(m \cdot_N p) +_N (q \cdot_N p) = n \cdot_N p$. Logo, $m \cdot_N p \leq_N n \cdot_N p$.
                \end{enumerate}
            Logo, se $m \leq_N n$, então $m +_N p \leq_N n +_N p$ e $m \cdot_N p \leq_N n \cdot_N p$. \blackproof
            \item Suponha que $m <_N n$. Então, existe $q \in N_{\neq a}$ tal que $m +_N q=n$.
                \begin{enumerate}[label=\roman*.]
                    \item Temos
                        \[
                            (m +_N p) +_N q = (m +_N q) +_N p = n +_N p,
                        \]
                    de modo que existe $q \in N_{\neq a}$ tal que $(m +_N p) +_N q = n +_N p$. Logo, $m +_N p <_N n +_N p$.
                    \item Seja $p \neq a$. Como $q \neq a$ e $p \neq a$, existe $k \in N$ tal que $p=s(k)$. Com isso,
                        \[
                            q \cdot_N p = q \cdot_N s(k) = (q \cdot_N k) +_N q.
                        \]
                    Se fosse $q \cdot_N p=a$, teríamos $(q \cdot_N k) +_N q=a$, donde $q=a$, uma contradição. Logo, $q \cdot_N p \in N_{\neq a}$. Além disso,
                        \begin{align*}
                            (m \cdot_N p) +_N (q \cdot_N p) &= (p \cdot_N m) +_N (p \cdot_N q) \\
                            &= p \cdot_N (m +_N q) \\
                            &= p \cdot_N n \\
                            &= n \cdot_N p,
                        \end{align*}
                    de modo que existe $q \cdot_N p \in N_{\neq a}$ tal que $(m \cdot_N p) +_N (q \cdot_N p) = n \cdot_N p$. Logo, $m \cdot_N p <_N n \cdot_N p$.
                \end{enumerate}
            Logo, se $m <_N n$, então $m +_N p <_N n +_N p$, e se $p \neq a$, então $m \cdot_N p <_N n \cdot_N p$. \blackproof
        \end{enumerate}
\end{proof}

\begin{proposition}
    Seja $(N,s,a)$ um sistema de Peano.
        \begin{enumerate}[leftmargin=*, align=left, label=\textbf{(\alph*)}]
            \item Para todo $n \in N$, temos $n \geq_N a$.
            \item Para todo $n \in N$, temos $s(n) >_N n$.
        \end{enumerate}
\end{proposition}

\begin{proof}
    \leavevmode
        \begin{enumerate}[leftmargin=*, align=left, label=\textbf{(\alph*)}]
            \item Como $a +_N n = n$ para todo $n \in N$, temos $n \geq_N a$. \blackproof
            \item Como $s(n) = n +_N 1_N$ para todo $n \in N$, temos $s(n) >_N n$. \blackproof
        \end{enumerate}
\end{proof}

\begin{theorem} \label{teo.set:tricotomia_naturais}
    Sejam $(N,s,a)$ um sistema de Peano e $m,n \in N$.
        \begin{enumerate}[leftmargin=*, align=left, label=\textbf{(\alph*)}]
            \item Para quaisquer $m,n\in N$, $m \leq_N n$ ou $n \leq_N m$.
            \item Para quaisquer $m,n\in N$, ou $m <_N n$, ou $n <_N m$, ou $m=n$.
        \end{enumerate}
\end{theorem}

\begin{proof}
    \leavevmode
        \begin{enumerate}[leftmargin=*, align=left, label=\textbf{(\alph*)}]
            \item Fixe $m \in N$ e defina
                \[
                    A := \{n \in N : m \leq_N n \lor n \leq_N m\}.
                \]
            Como $a \leq_N m$, temos $a \in A$. Agora, supondo que $n \in A$, provemos que $s(n) \in A$. Se $n \in A$, então, pela hipótese de indução, $m \leq_N n$ ou $n \leq_N m$.
                \begin{itemize}
                    \item Se $m \leq_N n$, então existe $p \in N$ tal que $m +_N p = n$. Com isso, $s(n) = s(m +_N p) = m +_N s(p)$, de modo que $m \leq_N s(n)$, isto é, $s(n) \in A$.
                    \item Se $n \leq_N m$, existe $p \in N$ tal que $n +_N p = m$. Se $p=a$, então $n=m$, e como $n <_N s(n)$, temos $m <_N s(n)$, donde $m \leq_N s(n)$, de modo que $s(n) \in A$. Por outro lado, se $p \neq a$, então existe $k \in N$ tal que $p = s(k)$, de modo que
                        \[
                            m = n +_N s(k) = s(n +_N k) = s(k +_N n) = k +_N s(n) = s(n) +_N k,
                        \]
                    isto é, existe $k \in N$ tal que $s(n) +_N k = m$, donde $s(n) \leq_N m$, isto é, $s(n) \in A$. Em qualquer caso, $s(n) \in A$.
                \end{itemize}
            Logo, pelo princípio de indução, $A = N$. Logo, a ordem $\leq_N$ é total. \blackproof
            \item Para quaisquer $m,n \in N$, temos $m \leq_N n$ ou $n \leq_N m$. Se $m \leq_N n$, então, se $m=n$, vale uma das condições; se $m \neq n$, então $m <_N n$ e vale uma das condições. Analogamente, se $n \leq_N m$, então, se $m=n$, vale uma das condições; se $m \neq n$, então $n <_N m$ e vale uma das condições. Logo, ao menos uma das condições das três condições deve valer. Resta provar a exclusividade. Se $m <_N n$, então $m \neq n$. Analogamente, se $n <_N m$, então $n \neq m$. Por fim, se fosse $m <_N n$ e $n <_N m$, então, como $<_N$ é transitiva, teríamos $n < _N n$, o que contradiz a irreflexividade de $<_N$. Logo, vale uma, e só uma, dessas três condições. \blackproof
        \end{enumerate}
\end{proof}

\begin{corollary}
    Sejam $(N,s,a)$ um sistema de Peano e $m,n,p \in N$.
        \begin{enumerate}[leftmargin=*, align=left, label=\textbf{(\alph*)}]
            \item Se $m \cdot_N n = a$, então $m = a$ ou $n=a$.
            \item Se $m \cdot_N p = n \cdot_N p$ e $p \neq a$, então $m=n$.
        \end{enumerate}
\end{corollary}

\begin{proof}
    \leavevmode
        \begin{enumerate}[leftmargin=*, align=left, label=\textbf{(\alph*)}]
            \item Se fosse $m \neq a$ e $n \neq a$, então, como $m \geq_N a$, teríamos $a <_N m$, de modo que $a \cdot_N n <_N m \cdot_N n$; como $a \cdot_N n = a$ e $m \cdot_N n = a$, teríamos $a <_N a$, uma contradição. Logo, $m=a$ ou $n=a$. \blackproof
            \item Pela tricotomia de $<_N$, ou $m <_N n$, ou $n <_N m$, ou $m=n$. Se fosse $m <_N n$, então, como $p \neq a$, teríamos $m \cdot_N p <_N n \cdot_N p$, e como $m \cdot_N p = n \cdot_N p$, teríamos $m \cdot_N p <_N m \cdot_N p$, uma contradição. Analogamente, se fosse $n <_N m$, então, como $p \neq a$, teríamos $n \cdot_N p <_N m \cdot_N p$,
            donde $n \cdot_N p <_N n \cdot_N p$, uma contradição. Logo, $m=n$. \blackproof
        \end{enumerate}
\end{proof}

\begin{proposition}
    Sejam $(N,s,a)$ um sistema de Peano e $m,n \in N$.
        \begin{enumerate}[leftmargin=*, align=left, label=\textbf{(\alph*)}]
            \item $m \leq_N n$ se, e somente se, $m <_N s(n)$.
            \item $m <_N n$ se, e somente se, $s(m) \leq_N n$.
        \end{enumerate}
\end{proposition}

\begin{proof}
    \leavevmode
        \begin{enumerate}[leftmargin=*, align=left, label=\textbf{(\alph*)}]
            \item 
                \begin{itemize}
                    \item[($\Rightarrow$)] Se $m \leq_N n$, então existe $p \in N$ tal que $m +_N p = n$. Com isso, $m +_N s(p) = s(m +_N p) = s(n)$, e como $s(p) \in N_{\neq a}$, temos $m <_N s(n)$. \blackproof
                    \item[($\Leftarrow$)] Se $m <_N s(n)$, então existe $p \in N_{\neq a}$ tal que $m +_N p = s(n)$. Como $p \neq a$, existe $k \in N$ tal que $p = s(k)$. Com isso, $s(m +_N k) = m +_N s(k) = m +_N p = s(n)$, e como $s$ é injetiva, temos $m +_N k = n$, de modo que $m \leq_N n$. \blackproof
                \end{itemize}
            Logo, $m \leq_N n \Leftrightarrow m <_N s(n)$. \blackproof
            \item 
                \begin{itemize}
                    \item[($\Rightarrow$)] Se $m <_N n$, então existe $p \in N_{\neq a}$ tal que $m +_N p = n$. Como $p \neq a$, existe $k \in N$ tal que $p = s(k)$. Com isso, $s(m) +_N k = m +_N s(k) = m +_N p = n$, de modo que $s(m) \leq_N n$. \blackproof
                    \item[($\Leftarrow$)] Se $s(m) \leq_N n$, então existe $p \in N$ tal que $s(m) +_N p = n$. Com isso, $m +_N s(p) = s(m) +_N p = n$, e como $s(p) \in N_{\neq a}$, temos $m <_N n$. \blackproof
                \end{itemize}
            Logo, $m <_N n \Leftrightarrow s(m) \leq_N n$. \blackproof
        \end{enumerate}
\end{proof}

\begin{theorem} \label{teo.set:inducao_e_boaordem}
    Seja $(N,s,a)$ um sistema de Peano.
        \begin{enumerate}[leftmargin=*, align=left, label=\textbf{(\alph*)}]
            \item Seja $A \subseteq N$ tal que, para todo $n \in N$, se $m \in A$ para todo $m \in N$ tal que $m <_N n$, então $n \in A$. Então, $A = N$.
            \item Para todo $A \in \mathcal{P}(N)_{\neq \emptyset}$ existe $b \in A$ tal que $b \leq_N x$ para todo $x \in A$.
            \item As afirmações acima são equivalentes.
        \end{enumerate}
\end{theorem}

\begin{proof}
    \leavevmode
        \begin{enumerate}[leftmargin=*, align=left, label=\textbf{(\alph*)}]
            \item Defina
                \[
                    B := \{n \in N : \forall m (m \in N \land m \leq_N n \rightarrow m \in A)\}.
                \]
            Provemos por indução que $B = N$.
                \begin{claim}
                    $a \in B$.
                \end{claim}
                \begin{midproof}
                    Seja $m \in N$ tal que $m \leq_N a$. Então existe $p \in N$ tal que $m +_N p = a$, de modo que $m = p = a$. Como para todo $k \in N$ tal que $k <_N a$ temos $k \in A$ (por vacuidade, já que não existe um tal $k$), temos $a \in A$. Logo, temos $m \in A$, de modo que $a \in B$. \whiteproof
                \end{midproof}
                \begin{claim}
                    Para todo $n \in N$, se $n \in B$, então $s(n) \in B$.
                \end{claim}
                \begin{midproof}
                    Suponha, por hipótese de indução, que $n \in B$. Então, para todo $m \in N$, se $m \leq_N n$, então $m \in A$. Provemos que, para todo $m \in N$, se $m \leq_N s(n)$, então $m \in A$. Seja $m \in N$ tal que $m \leq_N s(n)$. Então $m <_N s(n)$ ou $m = s(n)$.
                        \begin{itemize}
                            \item Se $m <_N s(n)$, então $m \leq_N n$, de modo que, pela hipótese de indução, $m \in A$.
                            \item Se $m = s(n)$, então, se $k \in N$ é tal que $k <_N s(n)$, então $k \leq_N n$; com isso, pela hipótese de indução, temos $k \in A$, de modo que, pela hipótese sobre $A$, temos $s(n) \in A$, isto é, $m \in A$.
                        \end{itemize}
                    Logo, em qualquer caso, temos $m \in A$, de modo que $s(n) \in B$. \whiteproof
                \end{midproof}
            Logo, pelo princípio de indução, temos $B = N$. Por fim, seja $n \in N$. Como $n \in B$, para todo $m \in N$, se $m \leq_N n$, então $m \in A$, e como $n \leq_N n$, temos $n \in A$, de modo que $B = N \subseteq A$. Como $A \subseteq N$, temos $A = N$. \blackproof
            \item Seja $A \subseteq N$ tal que $A \neq \emptyset$. Suponha, por absurdo, que $A$ não tenha elemento mínimo. Então, para todo $b \in A$, existe $x \in A$ tal que $x < b$. Defina
                \[
                    B := \{n \in N : \forall m (m \in N \land m \leq_N n \rightarrow m \notin A)\}.
                \]
            Provemos por indução que $B = N$.
                \begin{claim}
                    $a \in B$.
                \end{claim}
                \begin{midproof}
                    Como $a$ é o elemento mínimo de $N$ e $A \subseteq N$ não tem elemento mínimo, temos $a \notin A$. Seja $m \in N$ tal que $m \leq_N a$. Como $a \leq_N m$, temos $m = a$, e como $a \notin A$, temos $m \notin A$. Logo, $a \in B$. \whiteproof
                \end{midproof}
                \begin{claim}
                    Para todo $n \in N$, se $n \in B$, então $s(n) \in B$.
                \end{claim}
                \begin{midproof}
                    Suponha, por hipótese de indução, que $n \in B$. Então, para todo $m \in N$, se $m \leq_N n$, então $m \notin A$. Provemos que, para todo $m \in N$, se $m \leq_N s(n)$, então $m \notin A$. Seja $m \in N$ tal que $m \leq_N s(n)$. Então $m <_N s(n)$ ou $m = s(n)$.
                        \begin{itemize}
                            \item Se $m <_N s(n)$, então $m \leq_N n$, de modo que, pela hipótese de indução, $m \notin A$.
                            \item Se $m = s(n)$, então, se $s(n) \in A$, então existiria $k \in A$ tal que $k <_N s(n)$, isto é, $k \leq_N n$; com isso, pela hipótese de indução, teríamos $k \notin A$, uma contradição. Logo, $s(n) \notin A$, isto é, $m \notin A$.
                        \end{itemize}
                    Logo, em qualquer caso, temos $m \notin A$, de modo que $s(n) \in B$. \whiteproof
                \end{midproof}
            Logo, pelo princípio de indução, temos $B = N$. Por fim, como $A \neq \emptyset$, existe $n \in A$, e como $n \in N$, temos $n \in B$, de modo que, para todo $m \in N$, se $m \leq_N n$, então $m \notin A$. Como $n \leq_N n$, temos $n \notin A$, uma contradição. Logo, $A$ tem elemento mínimo. \blackproof
            \item 
                \begin{itemize}
                    \item[($\Rightarrow$)] Seja $A \subseteq N$ tal que $A \neq \emptyset$. Suponha, por absurdo, que $A$ não tenha elemento mínimo. Defina $B:= N \setminus A$. Seja $n \in N$ tal que, para todo $m \in N$, se $m <_N n$, então $m \in B$, isto é, $m \notin A$. Se fosse $n \notin B$, teríamos $n \in A$, e como para todo $x \in A$ não se tem $x <_N n$, pois $m \notin A$ para todo $m \in N$ tal que $m <_N n$, teríamos $n \leq_N x$ para todo $x \in A$, de modo que $n = \min A$, uma contradição; logo, $n \in B$. Logo, pelo princípio de indução forte, temos $B = N$, e como $B = N \setminus A$, temos $A = \emptyset$, uma contradição. Logo, $A$ tem elemento mínimo.
                    \item[($\Leftarrow$)] Seja $A \subseteq N$ tal que, para todo $n \in N$, se $m \in A$ para todo $m \in N$ tal que $m <_N n$, então $n \in A$. Suponha, por absurdo, que $A \neq N$. Defina $B := N \setminus A$. Como $A \neq N$, temos $B \neq \emptyset$. Logo, pelo princípio da boa ordem, existe $b := \min B$. Como $b \in B$, temos $b \notin A$. Agora, para todo $m \in N$ tal que $m <_N b$, temos $m \notin B$ (pois $b \leq_N x$ para todo $x \in B$), isto é, $m \in A$, de modo que $b \in A$, uma contradição. Logo, $A = N$.
                \end{itemize}
            Logo, o princípio de indução forte é equivalente ao princípio da boa ordem. \blackproof
        \end{enumerate}
\end{proof}

\begin{theorem} \label{teo.set:isomorfismo_sistemas_peano_soma_produto_ordem}
    Sejam $(N,s,a)$ e $(M,t,b)$ sistemas de Peano. Existe uma única bijeção $\varphi : N \to M$ tal que, para quaisquer $m,n \in N$,
        \begin{enumerate}[label=\roman*.]
            \item $\varphi(m +_N n) = \varphi(m) +_M \varphi(n)$.
            \item $\varphi(m \cdot_N n) = \varphi(m) \cdot_M \varphi(n)$.
            \item $m \leq_N n \Leftrightarrow \varphi(m) \leq_M \varphi(n)$.
        \end{enumerate}
\end{theorem}

\begin{proof}
    Pelo \cref{teo.set:isomorfismo_sistemas_peano}, existe uma única bijeção $\varphi : N \to M$ tal que $\varphi(a) = b$ e $\varphi(s(n)) = t(\varphi(n))$ para todo $n \in N$.
        \begin{enumerate}[label=\roman*.]
            \item Fixe $m \in N$ e defina
                $
                    A := \{n \in N : \varphi(m +_N n) = \varphi(m) +_M \varphi(n)\}.
                $
            Como
                \[
                    \varphi(m +_N a) = \varphi(m) = \varphi(m) +_M b = \varphi(m) +_M \varphi(a),
                \]
            temos $a \in A$. Agora, supondo que $n \in A$, provemos que $s(n) \in A$. Se $n \in A$, então, pela hipótese de indução, temos $\varphi(m +_N n) = \varphi(m) +_M \varphi(n)$, de modo que
                \begin{align*}
                    \varphi(m +_N s(n)) &= \varphi(s(m +_N n)) \\
                    &= t(\varphi(m +_N n)) \\
                    &= t(\varphi(m) +_M \varphi(n)) \\
                    &= \varphi(m) +_M t(\varphi(n)) \\
                    &= \varphi(m) +_M \varphi(s(n)),
                \end{align*}
            isto é, $s(n) \in A$. Logo, $A = N$, isto é, $\varphi(m +_N n) = \varphi(m) +_M \varphi(n)$, para quaisquer $m, n \in N$.
            \item Fixe $m \in N$ e defina $A := \{n \in N : \varphi(m \cdot_N n) = \varphi(m) \cdot_M \varphi(n)\}$. Como
                \[
                    \varphi(m \cdot_N a) = \varphi(a) = b = \varphi(m) \cdot_M b = \varphi(m) \cdot_M \varphi(a),
                \]
            temos $a \in A$. Agora, supondo que $n \in A$, provemos que $s(n) \in A$. Se $n \in A$, então, pela hipótese de indução, temos $\varphi(m \cdot_N n) = \varphi(m) \cdot_M \varphi(n)$, de modo que
                \begin{align*}
                    \varphi(m \cdot_N s(n)) &= \varphi((m \cdot_N n) +_N m) \\
                    &= \varphi(m \cdot_N n) +_M \varphi(m) \\
                    &= (\varphi(m) \cdot_M \varphi(n)) +_M \varphi(m) \\
                    &= \varphi(m) \cdot_M t(\varphi(n)) \\
                    &= \varphi(m) \cdot_M \varphi(s(n)),
                \end{align*}
            isto é, $s(n) \in A$. Logo, $A = N$, isto é, $\varphi(m \cdot_N n) = \varphi(m) \cdot_M \varphi(n)$, para quaisquer $m, n \in N$.
            \item Se $m = n$, então $\varphi(m) = \varphi(n)$, e se $\varphi(m) = \varphi(n)$, então, pela injetividade de $\varphi$, temos $m=n$. 
                \begin{itemize}
                    \item[($\Rightarrow$)] Suponha que $m <_N n$. Então existe $k \in N_{\neq a}$ tal que $m +_N k = n$. Logo, $\varphi(m) +_M \varphi(k) = \varphi(n)$, e como $\varphi(k) \neq b$, temos $\varphi(m) <_M \varphi(n)$.
                    \item[($\Leftarrow$)] Suponha que $\varphi(m) <_M \varphi(n)$. Se $m \geq_N n$, então $\varphi(m) \geq_M \varphi(n)$, uma contradição. Logo, $m <_N n$. 
                \end{itemize}
            Logo, $m \leq_N n \Leftrightarrow \varphi(m) \leq_M \varphi(n)$.
        \end{enumerate}
    Logo, $\varphi$ também traduz a adição, a multiplicação e a ordem entre os sistemas de Peano $(N,s,a)$ e $(M,t,b)$. \blackproof
\end{proof}

\subsection{O Axioma do Infinito}

\begin{definition}
    O \emph{sucessor} de um conjunto $x$ é definido como $x^+ := x \cup \{x\}$.
\end{definition}

\begin{proposition}
    \leavevmode
        \begin{enumerate}[leftmargin=*, align=left, label=\textbf{(\alph*)}]
            \item $\forall x \forall y ((y \in x^+) \leftrightarrow ((y \in x) \lor (y = x)))$.
            \item $\forall x (x \in x^+)$.
            \item $\forall x (x \subseteq x^+)$.
        \end{enumerate}
\end{proposition}

\begin{proof}
    \leavevmode
        \begin{enumerate}[leftmargin=*, align=left, label=\textbf{(\alph*)}]
            \item Pela definição de sucessor, $x^+ = x \cup \{x\}$. Como $y \in x \cup \{x\}$ equivale a $y \in x$ ou $y \in \{x\}$, e como $y \in \{x\}$ equivale a $y = x$, temos $y \in x^+$ se, e somente se, $y \in x$ ou $y = x$. \blackproof
            \item Tomando $y=x$ no item anterior, obtemos $(x \in x^+) \leftrightarrow ((x \in x) \lor (x = x))$. Como $x = x$, temos $(x \in x) \lor (x = x)$. Logo, $x \in x^+$. \blackproof
            \item Se $z \in x$, então $(z \in x) \lor (z = x)$, de modo que $z \in x^+$. Logo, $x \subseteq x^+$. \blackproof
        \end{enumerate}
\end{proof}

\begin{definition}
    Um conjunto $x$ é \emph{indutivo} se
        \[
            (\emptyset \in x) \land \forall y ( (y \in x) \rightarrow (y^+ \in x)).
        \]
\end{definition}

\begin{axiom}[Infinito]
    Existe um conjunto indutivo.
        \[
            \boxed{
                \exists x ( (\emptyset \in x) \land \forall y ((y \in x) \rightarrow (y^+ \in x)) )
            }
        \]
\end{axiom}

\begin{proposition} \label{prop.set:omega}
    \leavevmode
        \begin{enumerate}[leftmargin=*, align=left, label=\textbf{(\alph*)}]
            \item Para todo $I$ indutivo,
                \[
                    \omega(I) := \bigcap \{A \in \mathcal{P}{(I)} : A \text{ é indutivo}\}.
                \]
            é indutivo.
            \item Para quaisquer $I$ e $J$ é indutivos, $\omega(J) = \omega(I)$.
        \end{enumerate}  
\end{proposition}

\begin{proof} 
    Pelo axioma das partes, existe $\mathcal{P}{(I)}$. Pelo axioma da separação em $\mathcal{P}{(I)}$ com a fórmula ``$A$ é indutivo'', existe $z(I) :=  \{A \in \mathcal{P}{(I)} : A \text{ é indutivo} \}$. Como $z(I) \neq \emptyset$ pois $I \in z(I)$, existe $\omega(I) := \bigcap z(I)$. Logo, $\omega(I)$ está bem definido.
        \begin{enumerate}[leftmargin=*, align=left, label=\textbf{(\alph*)}]
            \item Como $\emptyset \in A$ para todo $A \subseteq I$ indutivo, temos $\emptyset \in \omega(I)$. Se $x \in \omega(I)$, então $x \in A$ para todo $A \subseteq I$ indutivo, de modo que $x^+ \in A$ para todo $A \subseteq I$ indutivo, isto é, $x^+ \in \omega(I)$. Logo, $\omega(I)$ é indutivo. \blackproof
            \item 
                \begin{claim}
                    $I \cap J$ é indutivo.
                \end{claim}
                \begin{midproof}
                    Como $I$ e $J$ são indutivos, temos $\emptyset \in I$ e $\emptyset \in J$, de modo que $\emptyset \in I \cap J$. Se $x \in I \cap J$, então $x \in I$ e $x \in J$, e como $I$ e $J$ são indutivos, temos $x^+ \in I$ e $x^+ \in J$, de modo que $x^+ \in I \cap J$. Logo, a interseção de conjuntos indutivos é um conjunto indutivo.  
                    \whiteproof
                \end{midproof}
            Como $I \cap J \subseteq I$, temos $I \cap J \in z(I)$, e como $\omega(I) \subseteq I \cap J \subseteq J$, temos $\omega(I) \in z(J)$ de modo que $\omega(J) \subseteq \omega(I)$. Analogamente, $\omega(I) \subseteq \omega(J)$. Logo, $\omega(I) = \omega(J)$. \blackproof
        \end{enumerate}
\end{proof}

\begin{remark}
    A \cref{prop.set:omega} nos diz que $\omega(I)$ é a interseção da família de todos os conjuntos indutivos e que o parâmtro $I$ pode ser suprimido. A seguinte definição só é possível devido a esse teorema.
\end{remark}

\begin{definition}
    O conjunto dos números naturais é definido como a interseção de todos os conjuntos indutivos. Ele é denotado por $\omega$.
\end{definition}

\subsubsection{A Existência de um Sistema de Peano}

\begin{theorem} \label{teo.set:omega_indução}
    Seja $A \subseteq \omega$ indutivo. Então $A = \omega$.
\end{theorem}

\begin{proof}
    Como $A$ é indutivo, temos $\omega \subseteq A$. Como $A \subseteq \omega$, temos $A = \omega$. \blackproof
\end{proof}

\begin{definition}
    Dizemos que $A$ é \emph{transitivo} se
        \[
            \forall x \forall y (x \in y \land y \in A \rightarrow x \in A).
        \]
\end{definition}

\begin{proposition}
    Um conjunto $A$ é transitivo se, e somente se, $a \subseteq A$ para todo $a \in A$.    
\end{proposition}

\begin{proof}
    \leavevmode
        \begin{itemize}[align=left]
            \item[($\Rightarrow$)] Suponha, primeiro, que $A$ é transitivo. Seja $a \in A$. Se $x \in a$, então, como $a \in A$ e $A$ é transitivo, temos $x \in A$. Logo, $a \subseteq A$.
            \item[($\Leftarrow$)] Reciprocamente, suponha que $a \subseteq A$ para todo $a \in A$. Se $x \in y$ e $y \in A$, então, pela hipótese, temos $y \subseteq A$, de modo que $x \in A$. Logo, $A$ é transitivo.
        \end{itemize}
    \blackproof
\end{proof}

\begin{lemma}
    Todo $n \in \omega$ é transitivo.
\end{lemma}

\begin{proof}
    Por indução. Defina $A := \{n \in \omega : n \text{ é transitivo} \}$. Como $\emptyset$ é trivialmente transitivo (por vacuidade), temos $\emptyset \in A$. Agora, supondo que $n \in A$, provemos que $n^+ \in A$. Se $n \in A$, então, pela hipótese de indução, temos $x \subseteq n$ para todo $x \in n$. Seja $x \in n^+$. Então $x \in n$ ou $x=n$. Se $x \in n$, então, pela hipótese de indução, temos $x \subseteq n$, e como $n \subseteq n^+$, temos $x \subseteq n^+$. Se $x = n$, então, como $n \subseteq n^+$, temos $x \subseteq n^+$. Logo, $n^+ \in A$. Logo, pelo princípio de indução, $A = \omega$. \blackproof
\end{proof}

\begin{corollary} \label{lema.set:omega_irreflexivo}
    $n \notin n$ para todo $n \in \omega$.
\end{corollary}

\begin{proof}
    Por indução. Defina $A := \{n \in \omega : n \notin n\}$. Como $\emptyset \notin \emptyset$, temos $\emptyset \in A$. Agora, supondo que $n \in A$, provemos que $n^+ \in A$. Se $n \in A$, então, pela hipótese de indução, temos $n \notin n$. Se fosse $n^+ \in n^+$, teríamos $n^+ \in n \cup \{n\}$, de modo que $n^+ = n$ ou $n^+ \in n$. Se $n^+ = n$, então $n \cup \{n\} = n$, de modo que $n \in n$, o que contradiz a hipótese de indução. Se $n^+ \in n$, então, como $n$ é transitivo, teríamos $n^+ \subseteq n$, e como $n \in n^+$, teríamos $n \in n$, o que contradiz a hipótese de indução. Logo, $n^+ \notin n^+$, de modo que $n^+ \in A$. Pelo princípio de indução, $A = \omega$. \blackproof
\end{proof}

\begin{theorem} \label{teo.set:existe_peano}
    A tripla $(\omega, x \mapsto x^+, \emptyset)$ é um sistema de Peano.
\end{theorem}

\begin{proof}
    Defina $s : \omega \to \omega$ por $s(x) = x^+$, isto é,
        \[
            s := \{(x,y) \in \omega \times \omega : y = x^+ \}.
        \]
        \begin{claim}
            $s$ é injetiva.
        \end{claim}
        \begin{midproof}
            Sejam $x,y \in \omega$ tais que $s(x)=s(y)$. Então $x^+ = y^+$. Como $x \in x^+$, temos $x \in y^+$, de modo que $x \in y$ ou $x=y$. Analogamente, como $y \in y^+$, temos $y \in x^+$, de modo que $y \in x$ ou $y=x$. Se fosse $x \neq y$, teríamos $x \in y$ e $y \in x$, e como $y$ é transitivo, teríamos $x \subseteq y$, de modo que $y \in y$, o que contradiz o \cref{lema.set:omega_irreflexivo}. Logo, $x=y$. Logo, $s$ é injetiva. \whiteproof
        \end{midproof}
        \begin{claim}
            $\emptyset \in \omega \setminus s(\omega)$.
        \end{claim}
        \begin{midproof}
            Como $\omega$ é indutivo, temos $\emptyset \in \omega$. Suponha, por absurdo, que $\emptyset \in s(\omega)$. Então existe $x \in \omega$ tal que $s(x)=\emptyset$, isto é, $x^+=\emptyset$. Como $x \in x^+$, temos $x \in \emptyset$, uma contradição. Logo, $\emptyset \notin s(\omega)$, de modo que $\emptyset \in \omega \setminus s(\omega)$. \whiteproof
        \end{midproof}
        \begin{claim}
            Para todo $A \subseteq \omega$, se $\emptyset \in A$ e $x^+ \in A$ para todo $x \in A$, então $A = \omega$.
        \end{claim}
        \begin{midproof}
            Segue do \cref{teo.set:omega_indução}. \whiteproof
        \end{midproof}
    Logo, $(\omega, x \mapsto x^+, \emptyset)$ é um sistema de Peano. \blackproof
\end{proof}

\subsection*{Algumas Propriedades}

\begin{theorem} \label{cor.set:bigcupomega}
    \leavevmode
        \begin{enumerate}[leftmargin=*, align=left, label=\textbf{(\alph*)}]
            \item $\omega$ é transitivo.
            \item $\bigcup \omega = \omega$.
        \end{enumerate}
\end{theorem}

\begin{proof}
    \leavevmode
        \begin{enumerate}[leftmargin=*, align=left, label=\textbf{(\alph*)}]
            \item Por indução. Defina $A := \{n \in \omega : n \subseteq \omega\}$. Como $\emptyset \subseteq \omega$, temos $\emptyset \in A$. Agora, supondo que $n \in A$, provemos que $n^+ \in A$. Se $n \in A$, então, pela hipótese de indução, temos $n \subseteq \omega$. Como $n \in \omega$, temos $\{n\} \subseteq \omega$, de modo que $n \cup \{n\} \subseteq \omega$, isto é, $n^+ \subseteq \omega$. Isso completa o passo indutivo. Logo, pelo princípio de indução, temos $A = \omega$. Logo, $\omega$ é transitivo. \blackproof
            \item Por um lado ($\subseteq$), se $x \in \bigcup \omega$, então existe $y \in \omega$ tal que $x \in y$, e como $y \subseteq \omega$ pois $y$ é transitivo, temos $x \in \omega$, de modo que $\bigcup \omega \subseteq \omega$. Por outro lado ($\supseteq$), se $x \in \omega$, então existe $x^+ \in \omega$ tal que $x \in x^+$, de modo que $x \in \bigcup \omega$, isto é, $\bigcup \omega \supseteq \omega$. Logo, $\bigcup \omega = \omega$. \blackproof
        \end{enumerate}
\end{proof}

\begin{proposition}
    \leavevmode
        \begin{enumerate}[leftmargin=*, align=left, label=\textbf{(\alph*)}]
            \item Para quaisquer $m,n \in \omega$, temos $m <_{\omega} n$ se, e somente se, $m \in n$.
            \item Para quaisquer $m,n \in \omega$, temos $m \leq_{\omega} n$ se, e somente se, $m \subseteq n$.
        \end{enumerate}
\end{proposition}

\begin{proof}
    \leavevmode
        \begin{enumerate}[leftmargin=*, align=left, label=\textbf{(\alph*)}]
            \item 
        \end{enumerate}
\end{proof}

\begin{corollary}
    Para todo $n \in \omega$,
        \[
            n = \{k \in \omega : k < n\}.
        \]
\end{corollary}

\begin{proof}
    \blackproof
\end{proof}

\subsection{Os Teoremas da Recursão}

\begin{theorem}[da recursão finita] \label{teo.set:recfin}
    Sejam $A \neq \emptyset$, $a \in A$ e $f:A \to A$. Existe uma única função $\varphi : \omega \to A$ tal que
        \begin{itemize}
            \item $\varphi(0) = a$;
            \item $\varphi(n^+) = f(\varphi(n))$, para todo $n \in \omega$.
        \end{itemize}
\end{theorem}

\begin{proof}
    Segue como corolário dos \cref{teo.set:recursao_peano,teo.set:existe_peano}. \blackproof
\end{proof}

\begin{theorem}[da recursão com parâmetro] \label{teo.set:recpar}
    Sejam $a \in A$ e $f : A \times \omega \to A$. Existe uma única função $\varphi : \omega \to A$ tal que
        \begin{itemize}
            \item $\varphi(0) = a$;
            \item $\varphi(n^+) = f(\varphi(n), n)$, para todo $n \in \omega$.
        \end{itemize}
\end{theorem}

\begin{proof}
    Segue como corolário dos \cref{teo.set:recpar1,teo.set:existe_peano}. \blackproof
\end{proof}

\begin{definition}
    O conjunto de todas as funções de um certo $n \in \omega$ em $X$ é denotado por $X^{< \omega}$, isto é,
        \[
            X^{< \omega} := \{f \in \mathcal{P}(\omega \times X) : \exists n (n \in \omega \land f \in X^n) \}.
        \]
\end{definition}

\begin{theorem}[da recursão completa] \label{teo.set:reccom}
    Seja $f :  X^{<\omega} \to X$. Existe uma única  $\varphi : \omega \to X$ tal que $\varphi(n) = f(\varphi \restriction_n)$ para todo $n \in \omega$.
\end{theorem}

\begin{proof}
    Defina $h : X^{< \omega} \to X^{< \omega}$ por $h(g) := g \cup \{(\Dom{(g)}, f(g)) \}$. Pelo \cref{teo.set:recfin}, existe uma única função $\psi : \omega \to X^{< \omega}$ tal que $\psi(0) = \emptyset$ e $\psi(n^+) = h(\psi(n))$.
        \begin{claim} \label{set:reccom.claim.dompsin}
            $\psi(n) \in X^n$ para todo $n \in \omega$.
        \end{claim}
        \begin{midproof}
            Dado $n \in \omega$, como $\psi(n) \in X^{< \omega}$, existe $m \in \omega$ para o qual $\psi(n) : m \to X$. Com isso, $\Img{(\psi(n))} \subseteq X$ e a afirmação equivale a provar que $m = n$, isto é, basta provar que $\Dom{(\psi(n))} = n$ para todo $n \in \omega$. Façamos isso por indução em $n$. Para $n=0$, como $\psi(0) = \emptyset$, temos $\Dom{(\psi(0))} = \Dom{(\emptyset)} = \emptyset = 0$. Suponha, por hipótese de indução, que $\Dom{(\psi(n))} = n$. Como
                \[
                    \psi(n^+) = h(\psi(n)) = \psi(n) \cup \{ (\Dom{(\psi(n))}, f(\psi(n)))\},
                \]
            obtemos
                \[
                    \Dom{(\psi(n^+))} = \Dom{(\psi(n))} \cup \{ \Dom{(\psi(n))} \},
                \]
            e como $\Dom{(\psi(n))} = n$, vem $\Dom{(\psi(n^+))} = n \cup \{ n \} = n^+$. Isso completa o passo indutivo. Como $\Dom{(\psi(n))} = n$ e $\Img{\psi(n)} \subseteq X$, vem $\psi(n) : n \to X$, isto é, $\psi(n) \in X^n$, para todo $n \in \omega$. \whiteproof
        \end{midproof}
        \begin{claim} \label{set:reccom.claim.restricao}
            Para quaisquer $m,n \in \omega$, se $m \subseteq n$, então $\psi(n) \restriction_m = \psi(m)$.
        \end{claim}
        \begin{midproof}
            A prova se dará por indução em $n$. Para $n=0$, obtemos $\psi(0) \restriction_m = \emptyset \restriction_m = \emptyset$, e como $m \subseteq \emptyset$, vem $m = \emptyset$, de modo que $\psi(0) \restriction_{\emptyset} = \emptyset = \psi(\emptyset) = \psi(0)$. Supondo, por hipótese de indução, que $\psi(n) \restriction_m = \psi(m)$ para todo $m \subseteq n$, provemos que $\psi(n^+) \restriction_m = \psi(m)$ para todo $m \subseteq n^+$. Se $m = n^+$, então $\psi(n^+) \restriction_{n^+} = \psi(n^+)$ trivialmente. Se $m \neq n^+$, então $m \subsetneq n^+$, de modo que $m \subseteq n$ (\cref{cor.set:bigcupomega}). 
                \begin{itemize}
                    \item Por um lado ($\subseteq$), seja $(k,x) \in \psi(n^+) \restriction_m$. Então $(k,x) \in \psi(n^+)$, com $k \in m$. Como $\psi(n^+) = \psi(n) \cup \{(n, f(\psi(n)))\}$, vem $(k,x) \in \psi(n)$ ou $(k,x) = (n, f(\psi(n)))$. Se fosse $(k,x) = (n, f(\psi(n)))$, então seria $k=n$, de modo que $n \in n$ pois $k \in m \subseteq n$, uma contradição. Logo, só pode ser $(k,x) \in \psi(n)$, e como $k \in m$, vem $(k,x) \in \psi(m)$ pois $\psi(n) \restriction_m = \psi(m)$ pela hipótese de indução. Assim, fica provado que $\psi(n^+) \restriction_m \subseteq \psi(m)$.
                    \item Por outro lado ($\supseteq$), seja $(k,x) \in \psi(m)$. Então, como $\psi(m) = \psi(n) \restriction_m$ pela hipótese de indução, temos $(k,x) \in \psi(n) \restriction_m$. Como $\psi(n) \subseteq \psi(n^+)$, temos $(k,x) \in \psi(n^+)$, como $k \in m$, temos $(k,x) \in \psi(n^+) \restriction_m$. Assim, fica provado que $\psi(n^+) \restriction_m \supseteq \psi(m)$.
                \end{itemize}
            Com isso, vem $\psi(n^+) \restriction_m = \psi(m)$ para todo $m \subseteq n^+$. Isso completa o passo indutivo. \whiteproof
        \end{midproof}
    Agora, defina $\varphi := \bigcup \Img{(\psi)}$. Essa é a função procurada. Provemos isso.
        \begin{claim} \label{set:reccom.claim.varphifuncao}
            $\varphi$ é uma função de $\omega$ em $X$.
        \end{claim}
        \begin{midproof}
            Comecemos com $\varphi$ ser função. Como $\varphi \subseteq \omega \times X$, provemos que se $(k,x) \in \varphi$ e $(k,y) \in \varphi$, então $x=y$. Se $(k,x), (k,y) \in \varphi = \bigcup \Img{(\psi)}$, então existem $n,m \in \omega$ tais que $(k,x) \in \psi(n)$ e $(k,y) \in \psi(m)$. Como a relação $\subseteq$ é total em $\omega$, suponha, sem perda de generalidade, que $m \subseteq n$. Com isso, $\psi(n) \restriction_m = \psi(m)$, e como $k \in m$ pois $(k,y) \in \psi(m)$, vem $(k,x) \in \psi(m)$. Como $\psi(m)$ é uma função, vem $x=y$, de modo que $\varphi$ é uma função. Agora, vejamos que $\Dom{(\varphi)} = \omega$. Como
                \[
                    \varphi = \bigcup \Img{(\psi)} = \bigcup \{\psi(n) : n \in \omega \} = \bigcup_{n \in \omega} \psi(n),
                \]
            obtemos
                \[
                    \Dom{(\varphi)} = \Dom{\left( \bigcup_{n \in \omega} \psi(n) \right)} = \bigcup_{n \in \omega} \Dom{(\psi(n))} = \bigcup_{n \in \omega} n = \bigcup \omega = \omega,
                \]
            pois $\Dom{(\psi(n))} = n$ para todo $n \in \omega$ (\cref{set:reccom.claim.dompsin}) e $\bigcup \omega = \omega$ (\cref{cor.set:bigcupomega}). Por fim, verifiquemos que $\Img{(\varphi)} \subseteq X$. Para isso, basta ver que
                \[
                    \Img{(\varphi)} = \Img{\left(\bigcup_{n \in \omega} \psi(n) \right)} = \bigcup_{n \in \omega} \Img{(\psi(n))} \subseteq X,
                \]
            pois $\Img{(\psi(n))} \subseteq X$ para todo $n \in \omega$. Com isso, $\varphi : \omega \to X$. \whiteproof
        \end{midproof}
        \begin{claim} \label{set:reccom.claim.propriedaderecursiva}
            $\varphi(n) = f(\varphi \restriction_n)$ para todo $n \in \omega$.
        \end{claim}
        \begin{midproof}
            Seja $n \in \omega$. Temos $(n, f(\psi(n))) \in \psi(n^+)$ pois 
                \[
                    \psi(n^+) = \psi(n) \cup \{ (n, f(\psi(n)))\},
                \]
            e como $\varphi = \bigcup \Img{(\psi)}$, temos $\psi(n^+) \subseteq \varphi$, de modo que $(n, f(\psi(n))) \in \varphi$; como $\varphi : \omega \to X$ é uma função (\cref{set:reccom.claim.varphifuncao}), vem  $\varphi(n) = f(\psi(n))$. Para finalizarmos, vejamos que $\psi(n) = \varphi \restriction_n$. Por um lado ($\subseteq$), se $(k,x) \in \psi(n)$, então, como $\psi(n) \in \Img{(\psi)}$, temos $(k,x) \in \bigcup \Im{(\psi)} = \varphi$. Com isso, como $k \in n$ pois $\Dom{(\psi(n))} = n$ (\cref{set:reccom.claim.dompsin}), vem $(k,x) \in \varphi \restriction_n$. Por outro lado ($\supseteq$), se $(k,x) \in \varphi \restriction_n$, então $(k,x) \in \varphi$ com $k \in n$ e existe $m \in \omega$ tal que $(k,x) \in \psi(m)$. Se $n \subseteq m$, então $\psi(m) \restriction_n = \psi(n)$ (\cref{set:reccom.claim.restricao}), e como $k \in n$, temos $(k,x) \in \psi(m) \restriction_n$, de modo que $(k,x) \in \psi(n)$. Se $m \subseteq n$, então $\psi(m) = \psi(n) \restriction_m$, e como $k \in m$, temos $(k,x) \in \psi(n) \restriction_m$, isto é, $(k,x) \in \psi(n)$. Com isso, temos $\psi(n) = \varphi \restriction_n$ para todo $n \in \omega$, de modo que, por fim, como $\varphi(n) = f(\psi(n))$, temos $\varphi(n) = f(\varphi \restriction_n)$ para todo $n \in \omega$. \whiteproof
        \end{midproof}
    Com isso, fica provada a existência de uma função $\varphi: \omega \to X$ que tem as propriedades do enunciado. Provemos, por fim, sua unicidade.
        \begin{claim}
            A função $\varphi : \omega \to X$ é única.
        \end{claim}
        \begin{midproof}
            Seja $\Phi : \omega \to X$ tal que $\Phi(n) = f(\Phi \restriction_n)$ para todo $n \in \omega$. Provemos que $\Phi \restriction_n = \varphi \restriction_n$ para todo $n \in \omega$ por indução em $n$. Para $n=0$, trivialmente $\Phi \restriction_{0} = \emptyset = \varphi \restriction_{0}$. Suponha, por hipótese de indução, que $\Phi \restriction_n = \varphi \restriction_n$. Com isso,
                \begin{align*}
                    \Phi \restriction_{n^+} &= \Phi \restriction_n \cup \{ (n, \Phi(n)) \} \\
                    &= \varphi \restriction_n \cup \{ (n, f(\Phi \restriction_n)) \} \\
                    &= \varphi \restriction_n \cup \{ (n, f(\varphi \restriction_n)) \} \\
                    &= h(\varphi \restriction_n) = h(\psi(n)) \\
                    &= \psi(n^+) = \varphi \restriction_{n^+},
                \end{align*}
            o que completa o passo indutivo. Logo
                \[
                    \Phi(n) = \Phi \restriction_{n^+}(n) = \varphi \restriction_{n^+}(n) = \varphi(n)
                \]
            para todo $n \in \omega$, de modo que $\Phi = \varphi$. \whiteproof
        \end{midproof}
    Logo, existe uma única função $\varphi : \omega \to X$ tal que $\varphi(n) = f(\varphi \restriction_n)$ para todo $n \in \omega$. \blackproof
\end{proof}

\subsection{Produto Cartesiano Generalizado}

\begin{definition} \label{defi.set:produtocartesianogeneralizado}
    O produto cartesiano de uma função $f$ é definido como
        \[
            \prod f := \left\{ \varphi \in \left(\bigcup \Im f \right)^{\Dom f} : \forall x (x \in \Dom f \rightarrow \varphi(x) \in f(x)) \right\}.
        \]
\end{definition}

\begin{proposition}
    Seja $f$ uma função. Se $f(x) = A$ para todo $x \in \Dom f$, então
        \[
            \prod f = A^{\Dom f}.
        \]
\end{proposition}

\begin{proof}
    \blackproof
\end{proof}

\chapter{Equipotência}

\section{Os Teoremas de Cantor}

\begin{definition}
    Dizemos que $A$ e $B$ são \emph{equipotentes} se existe $f : A \to B$ bijetiva. Isso é denotado por $|A| = |B|$.
\end{definition}

\begin{proposition}
    \leavevmode
        \begin{enumerate}[leftmargin=*, align=left, label=\textbf{(\alph*)}]
            \item $|A| = |A|$.
            \item Se $|A| = |B|$, então $|B| = |A|$.
            \item Se $|A| = |B|$ e $|B| = |C|$, então $|A| = |C|$.
        \end{enumerate}
\end{proposition}

\begin{proof}
    \leavevmode
        \begin{enumerate}[leftmargin=*, align=left, label=\textbf{(\alph*)}]
            \item $\id_A$ é uma bijeção de $A$ em $A$. \blackproof
            \item Se $f$ é uma bijeção de $A$ em $B$, então $f^{-1}$ é uma bijeção de $B$ em $A$. \blackproof
            \item Se $f$ é uma bijeção de $A$ em $B$ e $g$ é uma bijeção de $B$ em $C$, então $g \circ f$ é uma bijeção de $A$ em $C$. \blackproof
        \end{enumerate}
\end{proof}

\begin{theorem}
    $|\mathcal{P}(A)| = |2^A|$.
\end{theorem}

\begin{proof}
    Seja $B \in \mathcal{P}(A)$. Defina $\mathcal{X}_B : A \to 2$ por
        \[
            \mathcal{X}_B(x) :=
                \begin{cases}
                    1, & \text{se } x \in B \\
                    0, & \text{se } x \in A \setminus B
                \end{cases}.
        \]
    Por fim, defina $\varphi : \mathcal{P}(A) \to 2^A$ por $\varphi(B) = \mathcal{X}_B$.
        \begin{claim}
            $\varphi$ é injetora.
        \end{claim}
        \begin{midproof}
            Sejam $B, C \in \mathcal{P}(A)$ tais que $\varphi(B) = \varphi(C)$, isto é, $\mathcal{X}_B = \mathcal{X}_C$. Com isso, para todo $x \in A$, temos $\mathcal{X}_B(x) = \mathcal{X}_C(x)$. Por um lado ($\subseteq$), se $x \in B$, então $\mathcal{X}_B(x) = 1$, donde $\mathcal{X}_C(x) = 1$, de modo que $x \in C$, isto é, $B \subseteq C$. Por outro lado ($\supseteq$), se $x \in C$, então $\mathcal{X}_C(x) = 1$, donde $\mathcal{X}_B(x) = 1$, de modo que $x \in B$, isto é, $C \subseteq B$. Logo, $B = C$. Logo, $\varphi$ é injetora. \whiteproof
        \end{midproof}
        \begin{claim}
            $\varphi$ é sobrejetora.
        \end{claim}
        \begin{midproof}
            Seja $f \in 2^A$, isto é, $f : A \to 2$. Defina $B := \{x \in A : f(x) = 1\}$. Como $B \subseteq A$, temos $B \in \mathcal{P}(A)$. Provemos que $\varphi(B) = f$, isto é, $\mathcal{X}_B = f$. Seja $x \in A$. Se $x \in B$, então $f(x) = 1$ e $\mathcal{X}_B(x) = 1$. Se $x \notin B$, então $f(x) \neq 1$, e como $f(x) \in 2 = \{0,1\}$, temos $f(x) = 0$; por outro lado, como $x \notin B$, temos $\mathcal{X}_B(x) = 0$. Em ambos os casos, $\mathcal{X}_B(x) = f(x)$. Logo, $\varphi(B) = f$. Logo, $\varphi$ é sobrejetora. \whiteproof
        \end{midproof}
    Logo, $\varphi$ é bijetora. Logo, $|\mathcal{P}(A)| = |2^A|$. \blackproof
\end{proof}

\begin{definition}
    Dizemos que \emph{$B$ domina $A$}, e que \emph{a cardinalidade de $A$ é menor ou igual que a de $B$}, se existe $f : A \to B$ injetiva. Isso é denotado por $|A| \leq |B|$.
\end{definition}

\begin{proposition}
    \leavevmode
        \begin{enumerate}[leftmargin=*, align=left, label=\textbf{(\alph*)}]
            \item $|A| \leq |A|$.
            \item Se $|A| \leq |B|$ e $|B| \leq |A|$, então $|A| = |B|$.
            \item Se $|A| \leq |B|$ e $|B| \leq |C|$, então $|A| \leq |C|$.
        \end{enumerate}
\end{proposition}

\begin{proof}
    \leavevmode
        \begin{enumerate}[leftmargin=*, align=left, label=\textbf{(\alph*)}]
            \item $\id_A$ é uma injeção de $A$ em $A$. \blackproof
            \item Segue do \cref{teo.set:CSB}. \blackproof
            \item Se $f$ é uma injeção de $A$ em $B$ e $g$ é uma injeção de $B$ em $C$, então $g \circ f$ é uma injeção de $A$ em $C$. \blackproof
        \end{enumerate}
\end{proof}

\begin{proposition}
    \leavevmode
        \begin{enumerate}[leftmargin=*, align=left, label=\textbf{(\alph*)}]
            \item Se $|A| \leq |B|$ e $|A| = |C|$, então $|C| \leq |B|$.
            \item Se $|A| \leq |B|$ e $|B| = |C|$, então $|A| \leq |C|$.
        \end{enumerate}
\end{proposition}

\begin{proof}
    \leavevmode
        \begin{enumerate}[leftmargin=*, align=left, label=\textbf{(\alph*)}]
            \item Seja $f : A \to B$ injetora e $g : C \to A$ bijetora. Então, $f \circ g : C \to B$ é injetora. Logo, $|C| \leq |B|$. \blackproof
            \item Seja $f : A \to B$ injetora e $g : B \to C$ bijetora. Então, $g \circ f : A \to C$ é injetora. Logo, $|A| \leq |C|$. \blackproof
        \end{enumerate}
\end{proof}

\begin{lemma} \label{lema.set:CSB}
    \leavevmode
        \begin{enumerate}[leftmargin=*, align=left, label=\textbf{(\alph*)}]
            \item \fonte[157]{hercules}
            Se $B \subseteq A$ e $|A| \leq |B|$, então $|B| = |A|$.
            \item \fonte[67]{jechhrbacek1999} 
            Se $C \subseteq B \subseteq A$ e $|C| = |A|$, então $|B| = |A|$.
            \item \fonte[138]{fajardo2024conjuntos}
            Seja $F : \mathcal{P}(A) \to \mathcal{P}(A)$ tal que, para quaisquer $X,Y \in \mathcal{P}(A)$, se $X \subseteq Y$, então $F(X) \subseteq F(Y)$. Existe $C \in \mathcal{P}(A)$ tal que $F(C) = C$.
        \end{enumerate}
\end{lemma}

\begin{proof}
    \leavevmode
        \begin{enumerate}[leftmargin=*, align=left, label=\textbf{(\alph*)}]
            \item Como $|A| \leq |B|$, existe $f : A \to B$ injetiva. Defina $F : \mathcal{P}(A) \to \mathcal{P}(A)$ por $F(S) := f[S]$ para todo $S \in \mathcal{P}(A)$. Pelo \cref{teo.set:recursao_peano}, defina recursivamente $X : \N \to \mathcal{P}(A)$ por $X_0 = A \setminus B$ e $X_{n+1} = F(X_n) = f[X_n]$ para todo $n \in \N$, e seja
                \[
                    Y := \bigcup_{n \geq 0} X_n.
                \]
                \begin{claim}
                    $f \restriction_Y \cup \id_{A \setminus Y}$ é uma bijeção de $A$ em $B$.
                \end{claim}
                \begin{midproof}
                    Como $Y \cap (A \setminus Y) = \emptyset$, temos 
                        \[
                            \Dom \, (f \restriction_Y \cup \id_{A \setminus Y}) = \Dom f \restriction_Y \cup \Dom \id_{A \setminus Y} = Y \cup (A \setminus Y) = A.
                        \]
                    Como
                        \[
                            \Img f \restriction_Y = f[Y] = f \left[ \bigcup_{n \geq 0} X_n \right] = \bigcup_{n \geq 0} f[X_n] = \bigcup_{n \geq 0} X_{n+1} = \bigcup_{n \geq 1} X_n,
                        \]
                    temos $Y = X_0 \cup f[Y]$. Temos
                        \begin{align*}
                            \Img f \restriction_Y \cap \Img \id_{A \setminus Y} &= f[Y] \cap (A \setminus Y) \\
                            &= f[Y] \cap (A \setminus (X_0 \cup f[Y])) \\
                            &= f[Y] \cap ( (A \setminus X_0) \cap (A \setminus f[Y]) ) \\
                            &= (A \setminus X_0) \cap (f[Y] \cap (A \setminus f[Y])) \\
                            &= (A \setminus X_0) \cap \emptyset = \emptyset,
                        \end{align*}
                    donde
                        \begin{align*}
                            \Img \, (f \restriction_Y \cup \id_{A \setminus Y}) &= \Img f \restriction_Y \cup \Img \id_{A \setminus Y} \\
                            &= f[Y] \cup (A \setminus Y) \\
                            &= f[Y] \cup (A \setminus (X_0 \cup f[Y])) \\
                            &= f[Y] \cup ( (A \setminus X_0) \cap (A \setminus f[Y]) ) \\
                            &= (f[Y] \cup (A \setminus X_0) ) \cap (f[Y] \cup (A \setminus f[Y]) ) \\
                            &= (f[Y] \cup (A \setminus (A \setminus B) ) ) \cap A \\
                            &= f[Y] \cup ( (A \setminus A) \cup (A \cap B) ) \\
                            &= f[Y] \cup B.
                        \end{align*}
                    Como $X_n \subseteq B$ para todo $n \geq 1$, temos $f[Y] \subseteq B$, donde $f[Y] \cup B = B$, isto é,
                        \[
                            \Img \, (f \restriction_Y \cup \id_{A \setminus Y}) = B.
                        \]
                    Como $f \restriction_Y$ é bijetiva em $f[Y]$ e $\id_{A \setminus Y}$ é bijetiva em $A \setminus Y$, das observações acima segue que $f \restriction_Y \cup \id_{A \setminus Y}$ é bijetiva em $B$.
                    \whiteproof
                \end{midproof}
            Logo, $|B| = |A|$. \blackproof
            \item Como $|C| = |A|$, existe $f : A \to C$ bijetiva. Defina $F : \mathcal{P}(A) \to \mathcal{P}(A)$ por $F(S) := f[S]$ para todo $S \in \mathcal{P}(A)$. Pelo \cref{teo.set:recursao_peano}, defina recursivamente $X, Y : \N \to \mathcal{P}(A)$ por $X_0 = A$, $Y_0 = B$, e para todo $n \in \N$, $X_{n+1} = F(X_n) = f[X_n]$ e $Y_{n+1} = F(Y_n) = f[Y_n]$. Por fim, defina $Z : \N \to \mathcal{P}(A)$ por $Z_n := X_n \setminus Y_n$, para todo $n \in \N$, e seja
                \[
                    D := \bigcup_{n \geq 0} Z_n.
                \]
                \begin{claim}
                    $f \restriction_D \cup \id_{A \setminus D}$ é uma bijeção de $A$ em $B$.
                \end{claim}
                \begin{midproof}
                    Como $D \cap (A \setminus D) = \emptyset$, temos
                        \[
                            \Dom \, (f \restriction_D \cup \id_{A \setminus D}) = \Dom f \restriction_D \cup \Dom \id_{A \setminus D} = D \cup (A \setminus D) = A.
                        \]
                    Como
                        \[
                            f[Z_n] = f[X_n \setminus Y_n] = f[X_n] \setminus f[Y_n] = X_{n+1} \setminus Y_{n+1} = Z_{n+1},
                        \]
                    temos
                        \[
                            \Img f \restriction_D = f[D] = f \left[ \bigcup_{n \geq 0} Z_n \right] = \bigcup_{n \geq 0} f[Z_n] = \bigcup_{n \geq 0} Z_{n+1} = \bigcup_{n \geq 1} Z_n,
                        \]
                    donde $D = Z_0 \cup f[D]$. Temos
                        \begin{align*}
                            \Img f \restriction_D \cap \Img \id_{A \setminus D} &= f[D] \cap (A \setminus D) \\
                            &= f[D] \cap (A \setminus (Z_0 \cup f[D])) \\
                            &= f[D] \cap ( (A \setminus Z_0) \cap (A \setminus f[D]) ) \\
                            &= (A \setminus Z_0) \cap (f[D] \cap (A \setminus f[D])) \\
                            &= (A \setminus Z_0) \cap \emptyset = \emptyset,
                        \end{align*}
                    donde
                        \begin{align*}
                            \Img \, (f \restriction_D \cup \id_{A \setminus D}) &= \Img f \restriction_D \cup \Img \id_{A \setminus D} \\
                            &= f[D] \cup (A \setminus D) \\
                            &= f[D] \cup (A \setminus (Z_0 \cup f[D])) \\
                            &= f[D] \cup ( (A \setminus Z_0) \cap (A \setminus f[D]) ) \\
                            &= (f[D] \cup (A \setminus Z_0) ) \cap (f[D] \cup (A \setminus f[D]) ) \\
                            &= (f[D] \cup (A \setminus (X_0 \setminus Y_0) ) ) \cap A \\
                            &= f[D] \cup ( (A \setminus X_0) \cup (A \cap Y_0) ) \\
                            &= f[D] \cup Y_0 = f[D] \cup B.
                        \end{align*}
                    Como $Z_n \subseteq B$ para todo $n \geq 1$, temos $f[D] \subseteq B$, donde $f[D] \cup B = B$, isto é,
                        \[
                            \Img \, (f \restriction_D \cup \id_{A \setminus D}) = B.
                        \]
                    Como $f \restriction_D$ é bijetiva em $f[D]$ e $\id_{A \setminus D}$ é bijetiva em $A \setminus D$, das observações acima segue que $f \restriction_D \cup \id_{A \setminus D}$ é bijetiva em $B$. \whiteproof
                \end{midproof}
            Logo, $|B| = |A|$. \blackproof
            \item \blackproof
        \end{enumerate}
\end{proof}

\begin{theorem}[Cantor-Schröder-Bernstein] \label{teo.set:CSB}
    Se $|A| \leq |B|$ e $|B| \leq |A|$, então $|A| = |B|$.
\end{theorem}

\begin{proof} \fonte[158]{hercules}
    Como $|A| \leq |B|$, existe $f : A \to B$ injetiva. Como $|B| \leq |A|$, existe $g : B \to A$ injetiva. Como $g$ é bijetiva em $g[B]$, temos $|g[B]| = |B|$. Como $g \circ f : A \to g[B]$ é injetiva, temos $|A| \leq |g[B]|$. Como $g[B] \subseteq A$, pelo \cref{lema.set:CSB}, temos $|g[B]| = |A|$, de modo que $|A|=|B|$. \blackproof    
\end{proof}

\begin{proof} \fonte[66]{jechhrbacek1999}
    Como $|A| \leq |B|$, existe $f : A \to B$ injetiva. Como $|B| \leq |A|$, existe $g : B \to A$ injetiva. Como $g$ é bijetiva em $g[B]$, temos $|g[B]| = |B|$. Como $g \circ f$ é bijetiva em $g[f[A]]$, temos $|g[f[A]]| = |A|$. Temos $g[B] \subseteq A$, e como $f[A] \subseteq B$, temos $g[f[A]] \subseteq g[B] \subseteq A$. Logo, pelo \cref{lema.set:CSB}, temos $|g[B]| = |A|$, de modo que $|A| = |B|$. \blackproof
\end{proof}

\begin{proof} \fonte[138]{fajardo2024conjuntos}
    
\end{proof}

\begin{definition}
    Dizemos que \emph{$B$ domina estritamente $A$}, e que \emph{a cardinalidade de $A$ é estritamente menor que a de $B$}, se $|A| \leq |B|$ e $|A| \neq |B|$. Isso é denotado por $|A| < |B|$.
\end{definition}

\begin{theorem}[Cantor]
    $|A| < |\mathcal{P}(A)|$.
\end{theorem}

\begin{proof}
    Defina $f : A \to \mathcal{P}(A)$ por $f(x) = \{x\}$ para todo $x \in A$. Como $f$ é injetiva, temos $|A| \leq |\mathcal{P}(A)|$. Seja $f : A \to \mathcal{P}(A)$ e defina
        \[
            B := \{x \in A : x \notin f(x)\}.
        \]
    Se $B \in \Img f$, então existe $b \in A$ tal que $f(b) = B$. Se $b \in B$, então $b \notin f(a) = B$, e se $b \notin B$, então $b \in f(a) = B$, uma contradição. Logo, $B \notin \Img f$.
\end{proof}

\begin{proof}
    Provemos que $|A| \leq |\mathcal{P}(A)|$ e $|A| \neq |\mathcal{P}(A)|$.
        \begin{claim}
            $|A| \leq |\mathcal{P}(A)|$.
        \end{claim}
        \begin{midproof}
            Defina $f : A \to \mathcal{P}(A)$ por $f(x) = \{x\}$ para todo $x \in A$. Como $f$ é injetiva, temos $|A| \leq |\mathcal{P}(A)|$. \whiteproof
        \end{midproof}
        \begin{claim}
            $|A| \neq |\mathcal{P}(A)|$.
        \end{claim}
        \begin{midproof}
            Suponha, por absurdo, que $|A| = |\mathcal{P}(A)|$. Então existe $f : A \to \mathcal{P}(A)$ bijetiva; em particular, sobrejetiva. Defina $B := \{x \in A : x \notin f(x)\}$. Como $f$ é sobrejetiva, existe $a \in A$ tal que $f(a) = B$. No entanto, temos $a \in B \Leftrightarrow a \notin B$, uma contradição. Logo, não existe $f : A \to \mathcal{P}(A)$ sobrejetiva. Logo, $|A| \neq |\mathcal{P}(A)|$. \whiteproof 
        \end{midproof}
    Logo, $|A| < |\mathcal{P}(A)|$. \blackproof
\end{proof}

\part{Cálculo I}
%!TEX root = main.tex

\chapter{Limites e Continuidade}

\section{Sequências}

\begin{definition}
    \leavevmode
        \begin{enumerate}[leftmargin=*, align=left, label=\textbf{(\alph*)}]
            \item Uma sequência $(x_n)$ é \emph{convergente} e \emph{converge} para $a \in \R$ se para todo $\epsilon \in \R_{>0}$ existe $n_0 \in \N$ tal que, para todo $n \in \N$, se $n > n_0$, então $|x_n - a| < \epsilon$. Isso é denotado por
                \[
                    \lim x_n =a.
                \]
            Dizemos, nesse caso, que $a$ é o \emph{limite} de $(x_n)$.
            \item Uma sequência $(x_n)$ é \emph{divergente} se não é convergente.
        \end{enumerate}
\end{definition}

\begin{remark}
    As notações ``$\ds \lim_{n \to +\infty} x_n = a$'',  ``$\ds \lim_{n \in \N} x_n = a$'', ``$x_n \to a$ quando $n \to +\infty$'' e ``$x_n \to a$'' também são frequentemente usadas para indicar que $\lim x_n =a$.
\end{remark}

\begin{proposition}[Unicidade]
    O limite de uma sequência convergente é único.
\end{proposition}

\begin{proof}
    Provemos que se $(x_n)$ converge para $a \in \R$ e para $b \in \R$, então $a=b$. Suponha, por absurdo, que $a \neq b$, e seja $\epsilon := \dfrac{|b-a|}{2} \in \R_{>0}$. Como $\lim x_n = a$, existe $n_1 \in \N$ tal que, para todo $n \in \N$, se $n > n_1$, então $|x_n-a| < \epsilon$. Analogamente, como $\lim x_n = b$, existe $n_2 \in \N$ tal que, para todo $n \in \N$, se $n > n_2$, então $|x_n-b| < \epsilon$. Pois tome $n_0 := \max\{n_1,n_2\}$. Com isso, para todo $n \in \N$, se $n > n_0$, então $|x_n-a| < \epsilon$ e $|x_n-b| < \epsilon$, donde
        \begin{align*}
            |a-b| &= |(a-x_n) + (x_n-b)| \\
            &\leq |a-x_n| + |x_n-b| \\
            &< 2\epsilon = |a-b|,
        \end{align*}
    uma contradição. Logo, $a=b$. \blackproof
\end{proof}

\begin{theorem} \label{teo.calc1:algebra_das_sequencias}
    Sejam $(x_n)$ e $(y_n)$ convergentes.
        \begin{enumerate}[leftmargin=*, align=left, label=\textbf{(\alph*)}]
            \item A sequência $(x_n + y_n)$ é convergente e
                \[
                    \lim \, (x_n + y_n) = \lim x_n + \lim y_n.
                \]
            \item A sequência $(x_n \cdot y_n)$ é convergente e
                \[
                    \lim \, (x_n \cdot y_n) = \lim x_n \cdot \lim y_n.
                \]
            \item Se $\lim y_n \neq 0$, então a sequência $(x_n / y_n)$ é convergente e
                \[
                    \lim \left( \dfrac{x_n}{y_n} \right) = \dfrac{\lim x_n}{\lim y_n}.
                \]
        \end{enumerate}
\end{theorem}

\begin{proof}
    Sejam $\lim x_n = a$ e $\lim y_n = b$.
        \begin{enumerate}[leftmargin=*, align=left, label=\textbf{(\alph*)}]
            \item Seja $\epsilon \in \R_{>0}$. Como $\lim x_n = a$, para $\ds \epsilon_1 = \dfrac{\epsilon}{2} \in \R_{>0}$ existe $n_1 \in \N$ tal que, para todo $n \in \N$, se $n > n_1$, então $|x_n - a| < \epsilon_1$. Como $\lim y_n = b$, para $\ds \epsilon_2 = \dfrac{\epsilon}{2} \in \R_{>0}$ existe $n_2 \in \N$ tal que, para todo $n \in \N$, se $n > n_2$, então $|y_n - b| < \epsilon_2$. Pois tome $n_0 := \max \{n_1, n_2 \}$. Com isso, para todo $n \in \N$ tal que $n > n_0$, temos
                \begin{align*}
                    |(x_n + y_n) - (a + b)|
                        &\leq |x_n - a| + |y_n - b| \\
                        &< \dfrac{\epsilon}{2} + \dfrac{\epsilon}{2} \\
                        &= \epsilon,
                \end{align*}
            de modo que $\lim \, (x_n + y_n) = a + b$. \blackproof
            \item Seja $\epsilon \in \R_{>0}$. Como $\lim x_n = a$, para $\epsilon_1 = 1$ existe $n_1 \in \N$ tal que, para todo $n \in \N$, se $n > n_1$, então $|x_n - a| < \epsilon_1 = 1$; em particular, $|x_n| \leq 1 + |a|$. De fato, 
                \[
                    |x_n| = |x_n - a + a| \leq |x_n - a| + |a| < 1 + |a|.
                \]
            Como $\lim x_n = a$, para $\ds \epsilon_2 = \dfrac{\epsilon}{2(1 + |b|)} \in \R_{>0}$ existe $n_2 \in \N$ tal que, para todo $n \in \N$, se $n > n_2$, então $|x_n - a| < \epsilon_2$. Como $\lim y_n = b$, para $\epsilon_3 = \dfrac{\epsilon}{2(1 + |a|)} \in \R_{>0}$ existe $n_3 \in \N$ tal que, para todo $n \in \N$, se $n > n_3$, então $|y_n - b| < \epsilon_3$. Pois tome $n_0 := \max \{n_1, n_2, n_3 \}$. Observando que
                \[
                    x_n y_n - a b = x_n (y_n - b) + b (x_n - a),
                \]
            para todo $n \in \N$ tal que $n > n_0$, temos
                \begin{align*}
                    |x_n y_n - a b|
                        &\leq |x_n| |y_n - b| + |b| |x_n - a| \\
                        &< (1+|a|) \dfrac{\epsilon}{2(1+|a|)} + |b| \dfrac{\epsilon}{2(1+|b|)} \\
                        &< \dfrac{\epsilon}{2} + 1 \cdot \dfrac{\epsilon}{2} \\
                        &< \epsilon,
                \end{align*}
            de modo que $\lim \, (x_n \cdot y_n) = a \cdot b$. \blackproof
            \item Suponha $b \neq 0$. Seja $\epsilon \in \R_{>0}$. Como $\lim y_n = b$, para $\ds \epsilon_1 = \dfrac{|b|}{2} \in \R_{>0}$ existe $n_1 \in \N$ tal que, para todo $n \in \N$, se $n > n_1$, então $|y_n - b| < \epsilon_1$; em particular, $\dfrac{1}{|y_n|} < \dfrac{2}{|b|}$. De fato,
                \[
                    |y_n| = |y_n - b + b| \geq |b| - |y_n - b| > |b| - \dfrac{|b|}{2} = \dfrac{|b|}{2},
                \]
            donde $\dfrac{1}{|y_n|} < \dfrac{2}{|b|}$. Como $\lim x_n = a$, para $\ds \epsilon_2 = \dfrac{|b| \epsilon}{4} \in \R_{>0}$ existe $n_2 \in \N$ tal que, para todo $n \in \N$, se $n > n_2$, então $|x_n - a| < \epsilon_2$. Como $\lim y_n = b$, para $\ds \epsilon_3 = \dfrac{|b|^2 \epsilon}{4(1 + |a|)} \in \R_{>0}$ existe $n_3 \in \N$ tal que, para todo $n \in \N$, se $n > n_3$, então $|y_n - b| < \epsilon_3$. Pois tome $n_0 := \max \{n_1, n_2, n_3 \}$. Observando que
                \[
                    \dfrac{x_n}{y_n} - \dfrac{a}{b} = \dfrac{b(x_n - a) - a(y_n - b)}{y_n b},
                \]
            para todo $n \in \N$ tal que $n > n_0$, temos
                \begin{align*}
                    \left| \dfrac{x_n}{y_n} - \dfrac{a}{b} \right|
                        &= \dfrac{|b(x_n - a) - a(y_n - b)|}{|y_n| |b|} \\
                        &\leq \dfrac{|b| |x_n - a| + |a| |y_n - b|}{|y_n| |b|} \\
                        &< \dfrac{2}{|b|^2} \left( |b| \dfrac{|b| \epsilon}{4} + |a| \dfrac{|b|^2 \epsilon}{4(1 + |a|)} \right) \\
                        &< \dfrac{2}{|b|^2} \left( \dfrac{|b|^2 \epsilon}{4} + 1 \cdot \dfrac{|b|^2 \epsilon}{4} \right) \\
                        &= \epsilon,
                \end{align*}
            de modo que $\lim \left( \dfrac{x_n}{y_n} \right) = \dfrac{a}{b}$. \blackproof
        \end{enumerate}
\end{proof}

\begin{proposition}
    Seja $(x_n)$ convergente.
        \begin{enumerate}[leftmargin=*, align=left, label=\textbf{(\alph*)}]
            \item Para todo $a \in \R$, se $a < \lim x_n$, então, para todo $n$ suficientemente grande, $a < x_n$.
            \item Para todo $b \in \R$, se $\lim x_n < b$, então, para todo $n$ suficientemente grande, $x_n < b$.
        \end{enumerate}
\end{proposition}

\begin{proof}
    \leavevmode
        \begin{enumerate}[leftmargin=*, align=left, label=\textbf{(\alph*)}]
            \item Seja $\epsilon := \lim x_n - a \in \R_{>0}$. Então existe $n_0 \in \N$ tal que, para todo $n \in \N$, se $n > n_0$, então $\lim x_n - \epsilon < x_n < \lim x_n + \epsilon$; como $a = \lim x_n - \epsilon$, temos $a < x_n$ para todo $n$ suficientemente grande. \blackproof
            \item Seja $\epsilon := b - \lim x_n \in \R_{>0}$. Então existe $n_0 \in \N$ tal que, para todo $n \in \N$, se $n > n_0$, então $\lim x_n - \epsilon < x_n < \lim x_n + \epsilon$; como $b = \lim x_n + \epsilon$, temos $x_n < b$ para todo $n$ suficientemente grande. \blackproof
        \end{enumerate}
\end{proof}

\begin{corollary}
    Sejam $(x_n)$ e $(y_n)$ convergentes. Se $x_n \leq y_n$ para todo $n$ suficientemente grande, então $\lim x_n \leq \lim y_n$.
\end{corollary}

\begin{proof}
    Sejam $a := \lim x_n$ e $b := \lim y_n$. Suponha, por absurdo, que $b < a$. Seja $c \in \R$ tal que $b < c < a$. Então, para todo $n$ suficientemente grande, teríamos $y_n < c < x_n$, o que contradiz a hipótese. \blackproof
\end{proof}

\begin{theorem}[Confronto]
    Sejam $(x_n)$, $(y_n)$ e $(z_n)$ tais que $x_n \leq y_n \leq z_n$ para todo $n$ suficientemente grande. Se $\lim x_n = a$ e $\lim z_n = a$, então $\lim y_n = a$.
\end{theorem}

\begin{proof}
    Seja $\epsilon \in \R_{>0}$. Como $\lim x_n = a$, existe $n_1 \in \N$ tal que, para todo $n \in \N$, se $n > n_1$, então $a-\epsilon < x_n < a+\epsilon$. Como $\lim z_n = a$, existe $n_2 \in \N$ tal que, para todo $n \in \N$, se $n > n_2$, então $a-\epsilon < z_n < a+\epsilon$. Como $x_n \leq y_n \leq z_n$ para todo $n$ suficientemente grande, existe $n_3 \in \N$ tal que, para todo $n \in \N$, se $n > n_3$, então $x_n \leq y_n \leq z_n$. Seja $n_0 := \max \{n_1, n_2, n_3 \}$. Com isso, para todo $n \in \N$, se $n > n_0$, então 
        \[
            a-\epsilon < x_n \leq y_n \leq z_n < a+\epsilon,
        \]
    de modo que $|y_n - a| < \epsilon$. Logo, $\lim y_n = a$. \blackproof
\end{proof}

\begin{proposition}
    Toda sequência convergente é limitada.
\end{proposition}

\begin{proof}
    Sejam $(x_n)$ e $a \in \R$ tais que $\lim x_n = a$. Para $\epsilon = 1$, existe $n_0 \in \N$ tal que, para todo $n \in \N$, se $n > n_0$, então $|x_n - a| < 1$, isto é, $a-1 < x_n < a+1$. Então, sendo $S := \{x_1, x_2, \ldots, x_{n_0}, a-1, a+1\}$, temos $|x_n| \leq \max\{|\max\{S\}|, |\min\{S\}|\}$ para todo $n \in \N$, de modo que $(x_n)$ é limitada. \blackproof
\end{proof}

\begin{theorem}[Convergência Monótona] \label{teo.calc1:convergencia_monotona}
    \leavevmode
        \begin{enumerate}[leftmargin=*, align=left, label=\textbf{(\alph*)}]
            \item Toda sequência crescente e limitada superiormente é convergente.
            \item Toda sequência decrescente e limitada inferiormente é convergente.
            \item Toda sequência monótona e limitada é convergente.
        \end{enumerate}
\end{theorem}

\begin{proof}
    \leavevmode
        \begin{enumerate}[leftmargin=*, align=left, label=\textbf{(\alph*)}]
            \item Seja $(x_n)$ crescente e limitada superiormente. Como $\Im{x}$ é não vazio e limitado superiormente, existe $\sup{\Im{x}}$. Agora, seja $\epsilon \in \R_{>0}$. Então existe $n_0 \in \N$ tal que $\sup{\Im{x}} - \epsilon < x_{n_0}$ pois $\sup{\Im{x}} - \epsilon$ não é uma cota superior de $\Im{x}$. Como $(x_n)$ é crescente, para todo $n \in \N$, se $n > n_0$, então
                \[
                    \sup{\Im{x}} - \epsilon < x_{n_0} \leq x_n \leq \sup{\Im{x}} < \sup{\Im{x}} + \epsilon,
                \]
            de modo que $\lim x_n = \sup{\Im{x}}$. \blackproof
            \item Seja $(x_n)$ decrescente e limitada inferiormente. Como $\Im{x}$ é não vazio e limitado inferiormente, existe $\inf{\Im{x}}$. Agora, seja $\epsilon \in \R_{>0}$. Então existe $n_0 \in \N$ tal que $x_{n_0} < \inf{\Im{x}} + \epsilon$ pois $\inf{\Im{x}} + \epsilon$ não é uma cota inferior de $\Im{x}$. Como $(x_n)$ é decrescente, para todo $n \in \N$, se $n > n_0$, então
                \[
                    \inf{\Im{x}} - \epsilon < \inf{\Im{x}} \leq x_n \leq x_{n_0} < \inf{\Im{x}} + \epsilon,
                \]
            de modo que $\lim x_n = \inf{\Im{x}}$. \blackproof
            \item Segue dos itens acima. \blackproof
        \end{enumerate}
\end{proof}

\subsection{Subsequências}

\begin{definition} \label{defi:subseq}
    Sejam $x : \N \to \R$ e $n : \N \to \N$ estritamente crescente. Dizemos que $x \circ n : \N \to \R$ é uma \emph{subsequência} de $x$. Isso é denotado por $(x_{n_{k}})$.
\end{definition}

\begin{proposition} \fonte[26]{analisereal1}
    Seja $(x_n)$ monótona. Se existe $(x_{n_{k}})$ limitada, então $(x_n)$ é limitada.
\end{proposition}

\begin{proof}
    \blackproof
\end{proof}

\begin{proposition} \label{prop.calc1:subseqconv}
    Se $(x_n)$ converge para $a \in \R$, então toda subsequência $(x_{n_{k}})$ converge para $a$.
\end{proposition}

\begin{proof}
    Seja $\epsilon \in \R_{>0}$. Como $\lim x_n = a$, existe $n_0 \in \N$ tal que, para todo $n \in \N$, se $n > n_0$, então $|x_n - a| < \epsilon$. Como $(n_k)$ é estritamente crescente, existe $k_0 \in \N$ tal que $n_{k_0} > n_0$; com isso, para todo $k \in \N$, se $k > k_0$, então $n_k > n_{k_0} > n_0$, de modo que $|x_{n_{k}} - a| < \epsilon$. Logo, $\lim x_{n_{k}} = a$. \blackproof
\end{proof}

\begin{theorem}[Bolzano-Weierstrass] \label{teo.calc1:bolzanoweierstrass}
    Toda sequência limitada possui uma subsequência convergente.
\end{theorem}

\begin{proof}
    \blackproof
\end{proof}

\section{Séries}

\section{Topologia da Reta}

\begin{definition}
    Sejam $a \in \R$ e $r \in \R_{>0}$.
        \begin{enumerate}[leftmargin=*, align=left, label=\textbf{(\alph*)}]
            \item A \emph{$r$-vizinhança de $a$} é definida como
                \[
                    V_r(a) := \{x \in \R : |x-a|<r \}.
                \]
            \item A \emph{$r$-vizinhança furada de $a$} é definida como
                \[
                    V^*_r(a) := \{x \in \R : 0<|x-a|<r \}.
                \]
        \end{enumerate}
\end{definition}

\begin{corollary}
    Sejam $a \in \R$ e $r \in \R_{>0}$.
        \begin{enumerate}[leftmargin=*, align=left, label=\textbf{(\alph*)}]
            \item $V_r(a) = \left]a-r, a+r \right[$.
            \item $V^*_r(a) = V_r(a) \setminus \{a\}$.
        \end{enumerate}
\end{corollary}

\begin{proof}
    \leavevmode
        \begin{enumerate}[leftmargin=*, align=left, label=\textbf{(\alph*)}]
            \item \blackproof
            \item \blackproof
        \end{enumerate}
\end{proof}

\begin{definition}
    Dizemos que $a \in \R$ é um \emph{ponto de acumulação} de $A \subseteq \R$ se 
        \[
            V^*_{r}(a) \cap A \neq \emptyset
        \]
    para todo $r \in \R_{>0}$. O conjunto de todos os pontos de acumulação de $A$ é denotado por $A'$.
\end{definition}

\begin{theorem} \label{teo.calc1:a_linha_sse_sequencia}
    Sejam $A \subseteq \R$ e $a \in \R$. Então, $a \in A'$ se, e somente se, existe $(x_n)$ em $A_{\neq a}$ tal que $\lim x_n = a$. 
\end{theorem}

\begin{proof}
    Ver \cite{analisereal1}, página 49, teorema 6. \blackproof
\end{proof}

\section{Limites}

\begin{definition}
    Dizemos que $L \in \R$ é o \emph{limite} de uma função $f : A \subseteq \R \to \R$ em $a \in A'$ se para todo $\epsilon \in \R_{>0}$ existe $\delta = \delta(\epsilon, a) \in \R_{>0}$ tal que, para todo $x \in A$, se $0<|x-a|<\delta$, então $|f(x)-L| < \epsilon$. Isso é denotado por 
        \[
            \lim_{x \to a} f(x) = L.
        \]
\end{definition}

\begin{lemma} \label{calc1.lemma:existe_x0}
    Sejam $a \in \R$ e $A \subseteq \R$. Temos $a \in A'$ se, e somente se, para todo $\delta \in \R_{>0}$ existe $x_0 \in V^*_{\delta}(a) \cap A$.
\end{lemma}

\begin{proof}
    Por um lado ($\Rightarrow$), suponha que $a \in A'$. Seja $\delta \in \R_{>0}$. Como $a \in A'$, temos $V^*_{\delta}(a) \cap A \neq \emptyset$. Logo, existe $x_0 \in V^*_{\delta}(a) \cap A$. Por outro lado ($\Leftarrow$), suponha que, para todo $\delta \in \R_{>0}$, existe $x_0 \in V^*_{\delta}(a) \cap A$. Então $V^*_{\delta}(a) \cap A \neq \emptyset$ para todo $\delta \in \R_{>0}$, de modo que $a \in A'$. \blackproof
\end{proof}

\begin{proposition}[Unicidade] \label{teo.calc1:unicidade_limite_bilateral}
    O limite de $f : A \subseteq \R \to \R$ em $a \in A'$, quando existe, é único.
\end{proposition}

\begin{proof} \fonte[135]{cursodeanalise1}
    Provemos que se $\ds \lim_{x \to a} f(x) = L_1$ e $\ds \lim_{x \to a} f(x) = L_2$, então $L_1 = L_2$. Suponha, por absurdo, que $L_1 \neq L_2$, e seja $\epsilon := \dfrac{|L_1 - L_2|}{2} \in \R_{>0}$. Como $\ds \lim_{x \to a} f(x) = L_1$, existe $\delta_1 \in \R_{>0}$ tal que, para todo $x \in A$, se $0 < |x-a| < \delta_1$, então $|f(x) - L_1| < \epsilon$. Analogamente, como $\ds \lim_{x \to a} f(x) = L_2$, existe $\delta_2 \in \R_{>0}$ tal que, para todo $x \in A$, se $0 < |x-a| < \delta_2$, então $|f(x) - L_2| < \epsilon$. Pois tome $\delta := \min\{\delta_1, \delta_2 \}$. como $a \in A'$, pelo \cref{calc1.lemma:existe_x0} existe $x_0 \in A$ tal que $0<|x_0-a|<\delta$; com isso, obtemos $|f(x_0) - L_1| < \epsilon$ e $|f(x_0) - L_2| < \epsilon$, donde
        \begin{align*}
            |L_1 - L_2| &= |(f(x_0) - L_2) - (f(x_0)-L_1)| \\
            &\leq |f(x_0) - L_2| + |f(x_0)-L_1| \\
            &< 2 \epsilon = |L_1 - L_2|,
        \end{align*}
    uma contradição. Logo, $L_1 = L_2$. \blackproof
\end{proof}

\begin{theorem} \label{teo.calc1:limite_sequencial} \fonte[136]{cursodeanalise1}
    Sejam $f : A \subseteq \R \to \R$, $a \in A'$ e $L \in \R$. Temos $\ds \lim_{x \to a} f(x) = L$ se, e somente se, para toda sequência $(x_n)$ em $A \setminus \{a\}$, se $\lim x_n = a$, então $\lim f(x_n) = L$.
\end{theorem}

\begin{proof}
    \leavevmode
        \begin{itemize}[align=left]
            \item[($\Rightarrow$)] Suponha que $\ds \lim_{x \to a} f(x) = L$, e seja $(x_n)$ em $A \setminus \{a\}$ tal que $\lim x_n = a$. Seja $\epsilon \in \R_{>0}$. Então existe $\delta \in \R_{>0}$ tal que, para todo $x \in A$, se $0<|x-a|<\delta$, então $|f(x) - L|<\epsilon$. Como $\lim x_n = a$, existe $N \in \N$ tal que, para todo $n \in \N$, se $n > N$, então $|x_n - a| < \delta$. Como $x_n \neq a$ para todo $n \in \N$, temos $0<|x_n - a| < \delta$ para todo $n > N$. Logo, para todo $n \in \N$, se $n > N$, então $|f(x_n) - L| < \epsilon$, isto é, $\lim f(x_n) = L$.
            \item[($\Leftarrow$)] Suponha que $\lim f(x_n) = L$ para toda sequência $(x_n)$ em $A \setminus \{a\}$ tal que $\lim x_n = a$. Suponha, por absurdo, que $\ds \lim_{x \to a} f(x) \neq L$. Então existe $\xi \in \R_{>0}$ tal que, para todo $\delta \in \R_{>0}$, existe $x = x(\delta) \in A$ tal que $0<|x(\delta)-a|<\delta$ e $|f(x(\delta))-L| \geq \xi$. Em particular, para todo $n \in \N$, tomando $\delta_n := 1/n$, existe $x_n := x(\delta_n) \in A$ tal que $0<|x_n - a|< 1/n$ e $|f(x_n)-L| \geq \xi$.
                \begin{claim}
                    A sequência $(x_n)$ converge para $a$, isto é, $\lim x_n = a$.
                \end{claim}
                \begin{midproof}
                    Seja $\epsilon \in \R_{>0}$. Pela propriedade arquimediana, existe $N \in \N$ tal que $1/N < \epsilon$. Então, para todo $n \in \N$, se $n > N$, temos $1/n < 1/N$, de modo que
                        \[
                            |x_n - a| < 1/n < 1/N < \epsilon,
                        \]
                    isto é, $\lim x_n = a$ e $(x_n)$ converge para $a$. \whiteproof
                \end{midproof}
            Como $0<|x_n - a|$ para todo $n \in \N$, então $(x_n)$ é uma sequência em $A \setminus \{a\}$. Como $\lim x_n = a$, deveríamos ter $\lim f(x_n) = L$, mas existe $\xi \in \R_{>0}$ tal que $|f(x_n)-L| \geq \xi$ para todo $n \in \N$, uma contradição. Logo, $\ds \lim_{x \to a} f(x) = L$.
        \end{itemize}
    Isso completa a prova. \blackproof
\end{proof}

\begin{theorem} \label{teo.calc1:algebra_dos_limites}
    Sejam $f,g : A \subseteq \R \to \R$ e $a \in A'$ tais que existem $\ds \lim_{x \to a} f(x)$ e $\ds \lim_{x \to a} g(x)$.
        \begin{enumerate}[leftmargin=*, align=left, label=\textbf{(\alph*)}]
            \item O limite da soma é a soma dos limites:
                \[
                    \lim_{x \to a} \, (f+g)(x) = \lim_{x \to a} f(x) + \lim_{x \to a} g(x).
                \]
            \item O limite do produto é o produto dos limites:
                \[
                    \lim_{x \to a} \, (f \cdot g)(x) = \lim_{x \to a} f(x) \cdot \lim_{x \to a} g(x).
                \]
            \item O limite do quociente é igual ao quociente dos limites, desde que o denominador seja diferente de $0$: se $\ds \lim_{x \to a} g(x) \neq 0$, então
                \[
                    \lim_{x \to a} \left(\dfrac{f}{g} \right)(x) = \dfrac{\ds \lim_{x \to a} f(x)}{\ds \lim_{x \to a} g(x)}.
                \]
        \end{enumerate}
\end{theorem}

\begin{proof}
    Sejam $\ds \lim_{x \to a} f(x) = L_1$ e $\ds \lim_{x \to a} g(x) = L_2$.
        \begin{enumerate}[leftmargin=*, align=left, label=\textbf{(\alph*)}]
            \item Seja $\epsilon \in \R_{>0}$. Como $\ds \lim_{x \to a} f(x) = L_1$, para $\ds \epsilon_1 = \dfrac{\epsilon}{2} \in \mathbb{R}_{>0}$ existe $\delta_1 \in \mathbb{R}_{>0}$ tal que, para todo $x \in A$, se $0<|x-a|<\delta_1$, então $|f(x) - L_1| < \epsilon_1$. Como $\ds \lim_{x \to a} g(x) = L_2$, para $\ds \epsilon_2 = \dfrac{\epsilon}{2} \in \mathbb{R}_{>0}$ existe $\delta_2 \in \mathbb{R}_{>0}$ tal que, para todo $x \in A$, se $0<|x-a|<\delta_2$, então $|g(x) - L_2| < \epsilon_2$. Pois tome $\delta := \min \{\delta_1, \delta_2 \}$. Com isso, para todo $x \in A$ tal que $0 < |x-a| < \delta$, temos
                \begin{align*}
                    |(f(x) + g(x)) - (L_1 + L_2)|
                        &\leq |f(x) - L_1| + |g(x) - L_2| \\
                        &< \dfrac{\epsilon}{2} + \dfrac{\epsilon}{2} \\
                        &= \epsilon,
                \end{align*}
            de modo que $\ds \lim_{x \to a} (f+g)(x) = L_1 + L_2$. \blackproof
            \item \fonte[30]{aurichi}
            Seja $\epsilon \in \R_{>0}$. Como $\ds \lim_{x \to a} f(x) = L_1$, para $\epsilon_1 = 1$ existe $\delta_1 \in \R_{>0}$ tal que, para todo $x \in A$, se $0<|x-a|<\delta_1$, então $|f(x) - L_1| < \epsilon_1 = 1$; em particular, $|f(x)| \leq 1 + |L_1|$. De fato, 
                \[
                    |f(x)| = |f(x) - L_1 + L_1| \leq |f(x) - L_1| + |L_1| < 1 + |L_1|.
                \]
            Como $\ds \lim_{x \to a} f(x) = L_1$, para $\ds \epsilon_2 = \dfrac{\epsilon}{2(1 + |L_2|)} \in \R_{>0}$ existe $\delta_2 \in \R_{>0}$ tal que, para todo $x \in A$, se $0<|x-a|<\delta_2$, então $|f(x) - L_1| < \epsilon_2$. Como $\ds \lim_{x \to a} g(x) = L_2$, para $\epsilon_3 = \dfrac{\epsilon}{2(1 + |L_1|)} \in \R_{>0}$ existe $\delta_3 \in \R_{>0}$ tal que, para todo $x \in A$, se $0<|x-a|<\delta_3$, então $|g(x) - L_2| < \epsilon_3$. Pois tome $\delta := \min \{\delta_1, \delta_2, \delta_3 \}$. Observando que
                \[
                    f(x) g(x) - L_1 L_2 = f(x) (g(x) - L_2) + L_2 (f(x) - L_1),
                \]
            para todo $x \in A$ tal que $0 < |x-a| < \delta$, temos
                \begin{align*}
                    |f(x) g(x) - L_1 L_2|
                        &\leq |f(x)| |g(x) - L_2| + |L_2| |f(x) - L_1| \\
                        &< (1+|L_1|) \dfrac{\epsilon}{2(1+|L_1|)} + |L_2| \dfrac{\epsilon}{2(1+|L_2|)} \\
                        &< \dfrac{\epsilon}{2} + 1 \cdot \dfrac{\epsilon}{2} \\
                        &= \epsilon,
                \end{align*}
            de modo que $\ds \lim_{x \to a} f(x)g(x) = L_1 L_2$. \blackproof
            \item Suponha $L_2 \neq 0$. Seja $\epsilon \in \R_{>0}$. Como $\ds \lim_{x \to a} g(x) = L_2$, para $\ds \epsilon_1 = \dfrac{|L_2|}{2} \in \R_{>0}$ existe $\delta_1 \in \R_{>0}$ tal que, para todo $x \in A$, se $0<|x-a|<\delta_1$, então $|g(x) - L_2| < \epsilon_1$; em particular, $\dfrac{1}{|g(x)|} < \dfrac{2}{|L_2|}$. De fato,
                \[
                    |g(x)| = |g(x) - L_2 + L_2| \geq |L_2| - |g(x) - L_2| > |L_2| - \dfrac{|L_2|}{2} = \dfrac{|L_2|}{2},
                \]
            donde $\dfrac{1}{|g(x)|} < \dfrac{2}{|L_2|}$. Como $\ds \lim_{x \to a} f(x) = L_1$, para $\ds \epsilon_2 = \dfrac{|L_2| \epsilon}{4} \in \R_{>0}$ existe $\delta_2 \in \R_{>0}$ tal que, para todo $x \in A$, se $0<|x-a|<\delta_2$, então $|f(x) - L_1| < \epsilon_2$. Como $\ds \lim_{x \to a} g(x) = L_2$, para $\ds \epsilon_3 = \dfrac{|L_2|^2 \epsilon}{4(1 + |L_1|)} \in \R_{>0}$ existe $\delta_3 \in \R_{>0}$ tal que, para todo $x \in A$, se $0<|x-a|<\delta_3$, então $|g(x) - L_2| < \epsilon_3$. Pois tome $\delta := \min \{\delta_1, \delta_2, \delta_3 \}$. Observando que
                \[
                    \dfrac{f(x)}{g(x)} - \dfrac{L_1}{L_2} = \dfrac{L_2(f(x) - L_1) - L_1(g(x) - L_2)}{g(x)L_2},
                \]
            para todo $x \in A$ tal que $0 < |x-a| < \delta$, temos
                \begin{align*}
                    \left| \dfrac{f(x)}{g(x)} - \dfrac{L_1}{L_2} \right|
                        &= \dfrac{|L_2(f(x) - L_1) - L_1(g(x) - L_2)|}{|g(x)| |L_2|} \\
                        &\leq \dfrac{|L_2| |f(x) - L_1| + |L_1| |g(x) - L_2|}{|g(x)| |L_2|} \\
                        &< \dfrac{2}{|L_2|^2} \left( |L_2| \dfrac{|L_2| \epsilon}{4} + |L_1| \dfrac{|L_2|^2 \epsilon}{4(1 + |L_1|)} \right) \\
                        &< \dfrac{2}{|L_2|^2} \left( \dfrac{|L_2|^2 \epsilon}{4} + 1 \cdot \dfrac{|L_2|^2 \epsilon}{4} \right) \\
                        &= \epsilon,
                \end{align*}
            de modo que $\ds \lim_{x \to a} \dfrac{f(x)}{g(x)} = \dfrac{L_1}{L_2}$. \blackproof
        \end{enumerate}
\end{proof}

\begin{corollary}
    Sejam $n \in \N$, $A \subseteq \R$, $a \in A'$. Para cada $i \in [n]$, seja $f_i : A \to \R$ tal que existe $\ds \lim_{x \to a} f_i(x)$.
        \begin{enumerate}[leftmargin=*, align=left, label=\textbf{(\alph*)}]
            \item O limite da soma é igual à soma dos limites:
                \[
                    \lim_{x \to a} \left( \sum_{i=1}^{n} f_i(x) \right) = \sum_{i=1}^{n} \left( \lim_{x \to a} f_i(x) \right).
                \]
            \item O limite do produto é igual ao produto dos limites:
                \[
                    \lim_{x \to a} \left( \prod_{i=1}^{n} f_i(x) \right) = \prod_{i=1}^{n} \left( \lim_{x \to a} f_i(x) \right).
                \]
        \end{enumerate}
\end{corollary}

\begin{proof}
    Por indução. \blackproof
        \begin{enumerate}[leftmargin=*, align=left, label=\textbf{(\alph*)}]
            \item \blackproof
            \item \blackproof
        \end{enumerate}
\end{proof}

\begin{theorem} \fonte[135]{cursodeanalise1}
    Sejam $f : A \subseteq \R \to \R$, $B \subseteq A$, $a \in B'$ e $L \in \R$.
        \begin{enumerate}[leftmargin=*, align=left, label=\textbf{(\alph*)}]
            \item Se $\ds \lim_{x \to a} f(x) = L$, então $\ds \lim_{x \to a} f \restriction_B(x) = L$.
            \item Se $B = I \cap A$, onde $I \subseteq \R$ é um intervalo aberto tal que $a \in I$, então
                \[
                    \lim_{x \to a} f \restriction_{B}(x) = L \Rightarrow \lim_{x \to a} f(x) = L
                \]
        \end{enumerate}
\end{theorem}

\begin{proof}
    Seja $\epsilon \in \R_{>0}$.
        \begin{enumerate}[leftmargin=*, align=left, label=\textbf{(\alph*)}]
            \item Como $\ds \lim_{x \to a} f(x) = L$, existe $\delta \in \R_{>0}$ tal que, para todo $x \in A$, se $0<|x-a|<\delta$, então $|f(x)-L|<\epsilon$. Como $B \subseteq A$, para todo $x \in B$, se $0<|x-a|<\delta$, então $x \in A$ e, portanto,
                \[
                    |f \restriction_B(x)-L| = |f(x)-L| < \epsilon.
                \]
            Logo, $\ds \lim_{x \to a} f \restriction_B(x) = L$. \blackproof
            \item Como $\ds \lim_{x \to a} f \restriction_B(x) = L$, existe $\delta_1 \in \R_{>0}$ tal que, para todo $x \in B$, se $0<|x-a|<\delta_1$, então $|f \restriction_B(x)-L|<\epsilon$. Como $I$ é aberto e $a \in I$, existe $\delta_2 \in \R_{>0}$ tal que $V_{\delta_2}(a) \subseteq I$. Pois tome $\delta := \min \{\delta_1, \delta_2\}$. Logo, para todo $x \in A$, se $0<|x-a|<\delta$, então $x \in V_{\delta_2}(a) \subseteq I$, donde $x \in I \cap A = B$; com isso,
                \[
                    |f(x)-L| = |f \restriction_B(x)-L| < \epsilon.
                \]
            Logo, $\ds \lim_{x \to a} f(x) = L$. \blackproof
        \end{enumerate}
\end{proof}

\begin{theorem}[Conservação da Ordem] \label{teo.calc1:conservacao_da_ordem}
    Sejam $f,g : A \subseteq \R \to \R$ e $a \in A'$.
        \begin{enumerate}[leftmargin=*, align=left, label=\textbf{(\alph*)}]
            \item Se $\ds \lim_{x \to a} f(x) < \lim_{x \to a} g(x)$, então existe $r \in \R_{>0}$ tal que $f < g$ em $V^*_r(a) \cap A$.
            \item Se existe $r \in \R_{>0}$ tal que $f \leq g$ em $V^*_r(a) \cap A$, então $\ds \lim_{x \to a} f(x) \leq \lim_{x \to a} g(x)$.
        \end{enumerate}
\end{theorem}

\begin{proof}
    Sejam $\ds \lim_{x \to a} f(x) = L_1$ e $\ds \lim_{x \to a} g(x) = L_2$.
        \begin{enumerate}[leftmargin=*, align=left, label=\textbf{(\alph*)}]
            \item Seja $\ds \epsilon = \dfrac{L_2 - L_1}{2} \in \R_{>0}$. Como $\ds \lim_{x \to a} f(x) = L_1$, existe $\delta_1 \in \R_{>0}$ tal que, para todo $x \in A$, se $0<|x-a|<\delta_1$, então $|f(x) - L_1| < \epsilon$; em particular, $f(x) < L_1 + \epsilon$. Como $\ds \lim_{x \to a} g(x) = L_2$, existe $\delta_2 \in \R_{>0}$ tal que, para todo $x \in A$, se $0<|x-a|<\delta_2$, então $|g(x) - L_2| < \epsilon$; em particular, $L_2 - \epsilon < g(x)$. Pois tome $r := \min \{\delta_1, \delta_2 \}$. Logo, para todo $x \in V^*_r(a) \cap A$, isto é, $x \in A$ tal que $0 < |x-a| < r$, temos
                \[
                    f(x) < L_1 + \epsilon = L_2 - \epsilon < g(x),
                \]
            isto é, $f(x) < g(x)$ para todo $x \in V^*_r(a) \cap A$. \blackproof
            \item Suponha, por absurdo, que $L_2 < L_1$. Pelo item anterior, existe $r' \in \R_{>0}$ tal que $f > g$ em $V^*_{r'}(a) \cap A$. Pois tome $\delta := \min \{r, r'\}$. Então $f \leq g$ e $f > g$ em $V^*_{\delta}(a) \cap A$, uma contradição. Logo, $L_1 \leq L_2$. \blackproof
            \repeatitem Suponha, por absurdo, que $L_2 < L_1$. Como $\ds \lim_{x \to a} (f(x) - g(x)) = L_1 - L_2$, tomando $L:= L_1 - L_2$, para $\epsilon = L/2 \in \R_{>0}$ existe $\delta_1 \in \R_{>0}$ tal que, para todo $x \in A$, se $0<|x-a|<\delta_1$, então $|(f(x) - g(x)) - L| < L/2$. Em particular, $g < f$ em $V^*_{\delta_1}(a) \cap A$. Agora, seja $\delta := \min \{r, \delta_1\}$. Como $a \in A'$, existe $x_0 \in V^*_{\delta}(a) \cap A$. Como $x_0 \in V^*_{\delta_1}(a) \cap A$, temos $g(x_0) < f(x_0)$, e como $x_0 \in V^*_{r}(a) \cap A$, temos $f(x_0) \leq g(x_0)$, uma contradição. Logo, $L_1 \leq L_2$. \blackproof
        \end{enumerate}
\end{proof}

\begin{comment}
\begin{theorem}[Comparação]
    Sejam $f,g : A \subseteq \R \to \R$ e $a \in A'$. Se existe $r \in \R_{>0}$ tal que $f \leq g$ em $V^*_r(a) \cap A$, então
        \[
            \lim_{x \to a} f(x) \leq \lim_{x \to a} g(x),
        \]
    desde que os limites existam.
    %tais que $f(x) \leq g(x)$ para todo $x \in A$. Se existem $\ds \lim_{x \to a} f(x)$ e $\ds \lim_{x \to a} g(x)$, então $\ds \lim_{x \to a} f(x) \leq \lim_{x \to a} g(x)$.
    %Se $\ds \lim_{x \to a} f(x) < \lim_{x \to a} g(x)$, então existe $r \in \R_{>0}$ tal que $f \leq g$ em $V^*_{r}(a) \cap A$.
\end{theorem}

\begin{proof}
    Sejam $\ds \lim_{x \to a} f(x) = L_1$ e $\ds \lim_{x \to a} g(x) = L_2$. Suponha, por absurdo, que $L_2 < L_1$. Como
        $ \ds
            \lim_{x \to a} (f(x) - g(x)) = L_1 - L_2,
        $
    tomando $L:= L_1 - L_2$, para $\epsilon = \frac{L}{2} \in \R_{>0}$ existe $\delta_1 \in \R_{>0}$ tal que, para todo $x \in A$, se $0<|x-a|<\delta_1$, então
        $
            |(f(x) - g(x)) - L| < \frac{L}{2}.
        $
    Em particular,
    %$\ds 0 < \dfrac{L_1 - L_2}{2} < (f(x) - g(x)),$ de modo que
    $g(x) < f(x)$ para todo $x \in V^*_{\delta_1}(a) \cap A$. Agora, seja $\delta := \min \{r, \delta_1\}$. Como $a \in A'$, existe $x_0 \in V^*_{\delta}(a) \cap A$. Como $x_0 \in V^*_{\delta_1}(a) \cap A$, temos $g(x_0) < f(x_0)$, e como $x_0 \in V^*_{r}(a) \cap A$, temos $f(x_0) \leq g(x_0)$, uma contradição. Logo $L_1 \leq L_2$. \blackproof
\end{proof}

\begin{proof}
    Sejam $\ds \lim_{x \to a} f(x) = L_1$ e $\ds \lim_{x \to a} g(x) = L_2$. Suponha, por absurdo, que $L_2 < L_1$. Como
        \[
            \lim_{x \to a} (f(x) - g(x)) = \lim_{x \to a} f(x) - \lim_{x \to a} g(x) = L_1 - L_2,
        \]
    para $\epsilon = \dfrac{L_1 - L_2}{2} \in \R_{>0}$ existe $\delta_1 \in \R_{>0}$ tal que, para todo $x \in A$, se $0<|x-a|<\delta_1$, então
        \[
            |(f(x) - g(x)) - (L_1 - L_2)| < \dfrac{L_1 - L_2}{2},
        \]
    isto é,
        \[
            -\dfrac{L_1 - L_2}{2} < (f(x) - g(x)) - (L_1 - L_2) < \dfrac{L_1 - L_2}{2}.
        \]
    Em particular,
    %$\ds 0 < \dfrac{L_1 - L_2}{2} < (f(x) - g(x)),$ de modo que
    $g(x) < f(x)$ para todo $x \in V^*_{\delta_1}(a) \cap A$. Agora, seja $\delta := \min \{r, \delta_1\}$. Como $a \in A'$, existe $x_0 \in V^*_{\delta}(a) \cap A$. Como $x_0 \in V^*_{\delta_1}(a) \cap A$, temos $g(x_0) < f(x_0)$, e como $x_0 \in V^*_{r}(a) \cap A$, temos $f(x_0) \leq g(x_0)$, uma contradição. Logo $L_1 \leq L_2$. \blackproof
\end{proof}
\end{comment}

\begin{theorem}[Confronto]
    Sejam $f,g,h : A \subseteq \R \to \R$, $a \in A'$ e $L \in \R$. Se existe $r \in \R_{>0}$ tal que $f \leq g \leq h$ em $V^*_r(a) \cap A$, e se $\ds \lim_{x \to a} f(x) = L$ e $\ds \lim_{x \to a} h(x) = L$, então $\ds \lim_{x \to a} g(x) = L$.
\end{theorem}

\begin{proof}
    Seja $\epsilon \in \R_{>0}$. Como $\ds \lim_{x \to a} f(x) = L$, existe $\delta_1 \in \R_{>0}$ tal que, para todo $x \in A$, se $0<|x-a|<\delta_1$, então $L-\epsilon < f(x) < L+\epsilon$. Como $\ds \lim_{x \to a} h(x) = L$, existe $\delta_2 \in \R_{>0}$ tal que, para todo $x \in A$, se $0<|x-a|<\delta_2$, então $L-\epsilon < h(x) < L+\epsilon$. Seja $\delta := \min \{r, \delta_1, \delta_2 \}$. Com isso, para todo $x \in A$, se $0<|x-a|<\delta$, então 
        \[
            L-\epsilon < f(x) \leq g(x) \leq h(x) < L+\epsilon,
        \]
    isto é, $|g(x) - L| < \epsilon$. Logo, $\ds \lim_{x \to a} g(x) = L$. \blackproof
\end{proof}

\begin{definition}
    Dizemos que $f : A \subseteq \R \to \R$ é \emph{localmente limitada} em $a \in A \cup A'$ se existe $r \in \R_{>0}$ tal que $f \restriction_{V_r(a) \cap A}$ é limitada.
\end{definition}

\begin{proposition}
    Sejam $f,g : A \subseteq \R \to \R$ e $a \in A'$.
        \begin{enumerate}[leftmargin=*, align=left, label=\textbf{(\alph*)}]
            \item Se $\ds \lim_{x \to a} f(x) \in \R$, então $f$ é localmente limitada em $a$.
            \item Se $\ds \lim_{x \to a} f(x) = 0$ e $g$ é localmente limitada em $a$, então $\ds \lim_{x \to a} f(x)g(x) = 0$.
        \end{enumerate}
\end{proposition}

\begin{proof}
    \leavevmode
        \begin{enumerate}[leftmargin=*, align=left, label=\textbf{(\alph*)}]
            \item Seja $\ds L := \lim_{x \to a} f(x)$. Para $\epsilon = 1$, existe $\delta \in \R_{>0}$ tal que, para todo $x \in A$, se $0 < |x-a| < \delta$, então $|f(x) - L| < 1$. Com isso, para todo $x \in V^*_{\delta}(a) \cap A$, temos
                \[
                    |f(x)| = |f(x) - L + L| \leq |f(x)-L| + |L| < 1 + |L|.
                \]
            Se $a \notin A$, então $V^*_{\delta}(a) \cap A = V_{\delta}(a) \cap A$ e $f$ é localmente limitada em $a$. Se $a \in A$, então, para todo $x \in V_{\delta}(a) \cap A$, temos
                \[
                    |f(x)| \leq \max \{ 1+|L|, |f(a)| \},
                \]
            de modo que $f$ é localmente limitada em $a$. \blackproof
            \item Seja $\epsilon \in \R_{>0}$. Como $g$ é localmente limitada em $a$, existem $\delta_1, M \in \R_{>0}$ tal que, para todo $x \in A$, se $x \in V_{\delta_1}(a)$, então $|g(x)| \leq M$. Como $\ds \lim_{x \to a} f(x) = 0$, existe $\delta_2 \in \R_{>0}$ tal que, para todo $x \in A$, se $0<|x-a|<\delta_2$, então $|f(x)| < \frac{\epsilon}{M}$. Agora, tome $\delta := \min \{\delta_1, \delta_2\}$. Com isso, para todo $x \in A$, se $0<|x-a|<\delta$, então
                \[
                    |f(x) g(x)| = |f(x)| |g(x)| < M \cdot \dfrac{\epsilon}{M} = \epsilon,
                \]
            de modo que $\ds \lim_{x \to a} f(x) g(x) = 0$. \blackproof
        \end{enumerate}
\end{proof}

\begin{proposition}[Conservação do sinal]
    Sejam $f : A \subseteq \R \to \R$ e $a \in A'$.
        \begin{enumerate}[leftmargin=*, align=left, label=\textbf{(\alph*)}]
            \item Se $\ds \lim_{x \to a} f(x) \in \R_{>0}$, então existe $r \in \R_{>0}$ tal que $f(x) > 0$ para todo $x \in V^*_{r}(a) \cap A$.
            \item Se $\ds \lim_{x \to a} f(x) \in \R_{<0}$, então existe $r \in \R_{>0}$ tal que $f(x) < 0$ para todo $x \in V^*_{r}(a) \cap A$.
        \end{enumerate}
\end{proposition}

\begin{proof}
    Seja $\ds \lim_{x \to a} f(x) = L$.
        \begin{enumerate}[leftmargin=*, align=left, label=\textbf{(\alph*)}]
            \item Como $L \in \R_{>0}$, para $\epsilon = L$ existe $\delta \in \R_{>0}$ tal que, para todo $x \in A$, se $0<|x-a|<\delta$, então $|f(x) - L| < L$; em particular,
                \[
                    0 = L - L < f(x) < 2L.
                \]
            Logo, $f(x) > 0$ para todo $x \in V^*_{\delta}(a) \cap A$. \blackproof
            \item Como $L \in \R_{<0}$, para $\epsilon = -L$ existe $\delta \in \R_{>0}$ tal que, para todo $x \in A$, se $0<|x-a|<\delta$, então $|f(x) - L| < -L$; em particular,
                \[
                    2L < f(x) < 0 = L - L.
                \]
            Logo, $f(x) < 0$ para todo $x \in V^*_{\delta}(a) \cap A$. \blackproof
        \end{enumerate}
\end{proof}

\subsection{Limites Laterais}

\begin{definition}
    Sejam $a \in \R$ e $A \subseteq \R$.
        \begin{enumerate}[leftmargin=*, align=left, label=\textbf{(\alph*)}]
            \item Dizemos que $a$ é um \emph{ponto de acumulação à direita} de $A$ se 
                $
                    \left]a, a+ \delta \right[ \cap A \neq \emptyset
                $
            para todo $\delta \in \R_{>0}$. O conjunto de todos os pontos de acumulação à direita de $A$ é denotado por $A'_+$.
            \item Dizemos que $a$ é um \emph{ponto de acumulação à esquerda} de $A$ se
                $
                    \left]a - \delta, a\right[ \cap A \neq \emptyset
                $
            para todo $\delta \in \R_{>0}$. O conjunto de todos os pontos de acumulação à esquerda de $A$ é denotado por $A'_-$.
        \end{enumerate}
\end{definition}

\begin{definition}[Limites Laterais]
    Seja $f : A \subseteq \R \to \R$.
        \begin{enumerate}[leftmargin=*, align=left, label=\textbf{(\alph*)}]
            \item Dizemos que $L \in \R$ é o \emph{limite lateral à direita} de $f$ em $a \in A'_+$ se para todo $\epsilon \in \R_{>0}$ existe $\delta = \delta(\epsilon, a) \in \R_{>0}$ tal que, para todo $x \in A$, se $a < x < a + \delta$, então $|f(x)-L| < \epsilon$. Isso é denotado por
                \[
                    \lim_{x \to a^+} f(x) = L.
                \]
            \item Dizemos que $L \in \R$ é o \emph{limite lateral à esquerda} de $f$ em $a \in A'_-$ se para todo $\epsilon \in \R_{>0}$ existe $\delta = \delta(\epsilon, a) \in \R_{>0}$ tal que, para todo $x \in A$, se $a - \delta < x < a$, então $|f(x)-L| < \epsilon$. Isso é denotado por
                \[
                    \lim_{x \to a^-} f(x) = L.
                \]
        \end{enumerate}
\end{definition}

\begin{proposition}[Unicidade] \label{prop.calc1:unicidade_limites_laterais}
    Seja $f : A \subseteq \R \to \R$
        \begin{enumerate}[leftmargin=*, align=left, label=\textbf{(\alph*)}]
            \item O limite lateral à direita de $f$ em $a \in A'_+$, quando existe, é único.
            \item O limite lateral à esquerda de $f$ em $a \in A'_-$, quando existe, é único.
        \end{enumerate}
\end{proposition}

\begin{proof}
    \leavevmode
        \begin{enumerate}[leftmargin=*, align=left, label=\textbf{(\alph*)}]
            \item Provemos que se $\ds \lim_{x \to a^+} f(x) = L_1$ e $\ds \lim_{x \to a^+} f(x) = L_2$, então $L_1 = L_2$. Suponha, por absurdo, que $L_1 \neq L_2$, e seja $\epsilon := \dfrac{|L_1 - L_2|}{2} \in \R_{>0}$. Como $\ds \lim_{x \to a^+} f(x) = L_1$, existe $\delta_1 \in \R_{>0}$ tal que, para todo $x \in A$, se $a < x < a+\delta_1$, então $|f(x) - L_1| < \epsilon$. Analogamente, como $\ds \lim_{x \to a^+} f(x) = L_2$, existe $\delta_2 \in \R_{>0}$ tal que, para todo $x \in A$, se $a < x < a+\delta_2$, então $|f(x) - L_2| < \epsilon$. Pois tome $\delta := \min \{\delta_1, \delta_2\}$. Como $a \in A'_+$, existe $x_0 \in A$ tal que $a < x_0 < a+\delta$; com isso, obtemos $|f(x_0) - L_1| < \epsilon$ e $|f(x_0) - L_2| < \epsilon$, donde
                \begin{align*}
                    |L_1 - L_2| &= |(f(x_0) - L_2) - (f(x_0)-L_1)| \\
                    &\leq |f(x_0) - L_2| + |f(x_0)-L_1| \\
                    &< 2\epsilon \\
                    &= |L_1 - L_2|,
                \end{align*}
            uma contradição. Logo, $L_1 = L_2$. \blackproof
            \item Provemos que se $\ds \lim_{x \to a^-} f(x) = L_1$ e $\ds \lim_{x \to a^-} f(x) = L_2$, então $L_1 = L_2$. Suponha, por absurdo, que $L_1 \neq L_2$, e seja $\epsilon := \dfrac{|L_1 - L_2|}{2} \in \R_{>0}$. Como $\ds \lim_{x \to a^-} f(x) = L_1$, existe $\delta_1 \in \R_{>0}$ tal que, para todo $x \in A$, se $a-\delta_1 < x < a$, então $|f(x) - L_1| < \epsilon$. Analogamente, como $\ds \lim_{x \to a^-} f(x) = L_2$, existe $\delta_2 \in \R_{>0}$ tal que, para todo $x \in A$, se $a-\delta_2 < x < a$, então $|f(x) - L_2| < \epsilon$. Pois tome $\delta := \min \{\delta_1, \delta_2\}$. Como $a \in A'_-$, existe $x_0 \in A$ tal que $a-\delta < x_0 < a$; com isso, obtemos $|f(x_0) - L_1| < \epsilon$ e $|f(x_0) - L_2| < \epsilon$, donde
                \begin{align*}
                    |L_1 - L_2| &= |(f(x_0) - L_2) - (f(x_0)-L_1)| \\
                    &\leq |f(x_0) - L_2| + |f(x_0)-L_1| \\
                    &< 2\epsilon \\
                    &= |L_1 - L_2|,
                \end{align*}
            uma contradição. Logo, $L_1 = L_2$. \blackproof
        \end{enumerate}
\end{proof}

\begin{theorem}[Bilateral $\Leftrightarrow$ Laterais] \label{teo.calc1:bilatsselat}
    Sejam $f : A \subseteq \R \to \R$, $a \in A'_+ \cap A'_-$ e $L \in \R$. Então, $\ds \lim_{x \to a} f(x) = L$ se, e somente se, $\ds \lim_{x \to a^+} f(x) = L$ e $\ds \lim_{x \to a^-} f(x) = L$.
\end{theorem}

\begin{proof}
    \leavevmode
        \begin{itemize}[align=left]
            \item[($\Rightarrow$)] Suponha que $\ds \lim_{x \to a} f(x) = L$. Seja $\epsilon \in \R_{>0}$. Como $\ds \lim_{x \to a} f(x) = L$, existe $\delta \in \R_{>0}$ tal que, para todo $x \in A$, se $0<|x-a|<\delta$, então $|f(x)-L| < \epsilon$. Com isso, para todo $x \in A$, se $a < x < a+\delta$, então $0 < x-a < \delta$, isto é, $0 < |x-a| < \delta$, donde $|f(x)-L| < \epsilon$. Logo, $\ds \lim_{x \to a^+} f(x) = L$. Analogamente, para todo $x \in A$, se $a-\delta < x < a$, temos $0 < a-x < \delta$, isto é, $0 < |x-a| < \delta$, donde $|f(x)-L| < \epsilon$. Logo, $\ds \lim_{x \to a^-} f(x) = L$.
            \item[($\Leftarrow$)] Suponha que $\ds \lim_{x \to a^+} f(x) = L$ e $\ds \lim_{x \to a^-} f(x) = L$. Seja $\epsilon \in \R_{>0}$. Como $\ds \lim_{x \to a^+} f(x) = L$, existe $\delta_1 \in \R_{>0}$ tal que, para todo $x \in A$, se $a < x < a+\delta_1$, então $|f(x)-L| < \epsilon$. Como $\ds \lim_{x \to a^-} f(x) = L$, existe $\delta_2 \in \R_{>0}$ tal que, para todo $x \in A$, se $a-\delta_2 < x < a$, então $|f(x)-L| < \epsilon$. Pois tome $\delta := \min\{\delta_1, \delta_2\}$. Com isso, para todo $x \in A$ tal que $0 < |x-a| < \delta$, temos duas possibilidades: ou $x > a$ ou $x < a$. Se $x > a$, então $a < x < a+\delta \leq a+\delta_1$, donde $|f(x)-L| < \epsilon$. Se $x < a$, então $a-\delta_2 \leq a-\delta < x < a$, donde $|f(x)-L| < \epsilon$. Em ambos os casos, $|f(x) - L| < \epsilon$. Logo, $\ds \lim_{x \to a} f(x) = L$.
        \end{itemize}
    Logo, $\ds \lim_{x \to a} f(x) = L \Leftrightarrow \lim_{x \to a^+} f(x) = \lim_{x \to a^-} f(x) = L$. \blackproof
\end{proof}

\begin{proposition}
    Sejam $a \in \R$ e $A \subseteq \R$.
        \begin{enumerate}[leftmargin=*, align=left, label=\textbf{(\alph*)}]
            \item Temos $a \in A'_+$ se, e somente se, $a \in (A \cap \left]a,+\infty \right[)'$.
            \item Temos $a \in A'_-$ se, e somente se, $a \in (A \cap \left]-\infty, a \right[)'$.
        \end{enumerate}
\end{proposition}

\begin{proof}
    \leavevmode
        \begin{enumerate}[leftmargin=*, align=left, label=\textbf{(\alph*)}]
            \item Seja $\delta \in \R_{>0}$.
                \begin{itemize}
                    \item[($\Rightarrow$)] Como $a \in A'_+$, temos $\left]a, a+\delta\right[ \cap A \neq \emptyset$, isto é, existe $x \in \left]a, a+\delta\right[ \cap A$. Como $x \in \left]a, a+\delta\right[$, temos $a < x < a+\delta$, donde $x > a$ e $0 < |x-a| < \delta$. Com isso, $x \in \left]a,+\infty\right[$ e $x \in V^*_{\delta}(a)$. Como $x \in A$, temos $x \in V^*_{\delta}(a) \cap (A \cap \left]a,+\infty\right[)$, de modo que $V^*_{\delta}(a) \cap (A \cap \left]a,+\infty\right[) \neq \emptyset$. Logo, $a \in (A \cap \left]a,+\infty\right[)'$.
                    \item[($\Leftarrow$)] Como $a \in (A \cap \left]a,+\infty\right[)'$, temos $V^*_{\delta}(a) \cap (A \cap \left]a,+\infty\right[) \neq \emptyset$, isto é, existe $x \in V^*_{\delta}(a) \cap (A \cap \left]a,+\infty\right[)$. Como $x \in \left]a,+\infty\right[$, temos $x > a$. Como $x \in V^*_{\delta}(a)$, temos $0 < |x-a| < \delta$, isto é, $a - \delta < x < a + \delta$. Logo, $a < x < a + \delta$, donde $x \in \left]a, a+\delta\right[$. Como $x \in A$, temos $x \in \left]a, a+\delta\right[ \cap A$, de modo que $\left]a, a+\delta\right[ \cap A \neq \emptyset$. Logo, $a \in A'_+$.
                \end{itemize}
            Logo, $a \in A'_+ \Leftrightarrow a \in (A \cap \left]a,+\infty \right[)'$. \blackproof
            \item Seja $\delta \in \R_{>0}$.
                \begin{itemize}
                    \item[($\Rightarrow$)] Como $a \in A'_-$, temos $\left]a-\delta, a\right[ \cap A \neq \emptyset$, isto é, existe $x \in \left]a-\delta, a\right[ \cap A$. Como $x \in \left]a-\delta, a\right[$, temos $a-\delta < x < a$, donde $x < a$ e $0 < |x-a| < \delta$. Com isso, $x \in \left]-\infty,a\right[$ e $x \in V^*_{\delta}(a)$. Como $x \in A$, temos $x \in V^*_{\delta}(a) \cap (A \cap \left]-\infty,a\right[)$, de modo que $V^*_{\delta}(a) \cap (A \cap \left]-\infty,a\right[) \neq \emptyset$. Logo, $a \in (A \cap \left]-\infty,a\right[)'$.
                    \item[($\Leftarrow$)] Como $a \in (A \cap \left]-\infty,a\right[)'$, temos $V^*_{\delta}(a) \cap (A \cap \left]-\infty,a\right[) \neq \emptyset$, isto é, existe $x \in V^*_{\delta}(a) \cap (A \cap \left]-\infty,a\right[)$. Como $x \in \left]-\infty,a\right[$, temos $x < a$. Como $x \in V^*_{\delta}(a)$, temos $0 < |x-a| < \delta$, isto é, $a - \delta < x < a + \delta$. Logo, $a-\delta < x < a$, donde $x \in \left]a-\delta, a\right[$. Como $x \in A$, temos $x \in \left]a-\delta, a\right[ \cap A$, de modo que $\left]a-\delta, a\right[ \cap A \neq \emptyset$. Logo, $a \in A'_-$.
                \end{itemize}
            Logo, $a \in A'_- \Leftrightarrow a \in (A \cap \left]-\infty, a \right[)'$. \blackproof
        \end{enumerate}
\end{proof}

\begin{proposition}
    Seja $f : A \subseteq \R \to \R$ e $L \in \R$.
        \begin{enumerate}[leftmargin=*, align=left, label=\textbf{(\alph*)}]
            \item Se $a \in A'_+$, então
                \[
                    \lim_{x \to a^+} f(x) = L \Leftrightarrow \lim_{x \to a} f \restriction_{A \cap \left]a,+\infty \right[ } (x) = L.
                \]
            \item Se $a \in A'_-$, então
                \[
                    \lim_{x \to a^-} f(x) = L \Leftrightarrow \lim_{x \to a} f \restriction_{A \cap \left]-\infty, a \right[ } (x) = L.
                \]
        \end{enumerate}
\end{proposition}

\begin{proof}
    \leavevmode
        \begin{enumerate}[leftmargin=*, align=left, label=\textbf{(\alph*)}]
            \item Seja $\epsilon \in \R_{>0}$. Como $a \in A'_+$, temos $a \in (A \cap \left]a,+\infty\right[)'$.
                \begin{itemize}
                    \item[($\Rightarrow$)] Como $\ds \lim_{x \to a^+} f(x) = L$, existe $\delta \in \R_{>0}$ tal que, para todo $x \in A$, se $a < x < a+\delta$, então $|f(x) - L| < \epsilon$. Seja $x \in A \cap \left]a,+\infty\right[$ tal que $0 < |x-a| < \delta$. Como $x \in \left]a,+\infty\right[$, temos $x > a$, donde $a < x < a+\delta$. Como $x \in A$, segue que $|f(x) - L| < \epsilon$, isto é, $|f \restriction_{A \cap \left]a,+\infty\right[}(x) - L| < \epsilon$. Logo, $\ds \lim_{x \to a} f \restriction_{A \cap \left]a,+\infty\right[}(x) = L$.
                    \item[($\Leftarrow$)] Como $\ds \lim_{x \to a} f \restriction_{A \cap \left]a,+\infty\right[}(x) = L$, existe $\delta \in \R_{>0}$ tal que, para todo $x \in A \cap \left]a,+\infty\right[$, se $0 < |x-a| < \delta$, então $|f \restriction_{A \cap \left]a,+\infty\right[}(x) - L| < \epsilon$. Seja $x \in A$ tal que $a < x < a+\delta$. Então $x \in \left]a,+\infty\right[$ e $x \in V^*_{\delta}(a) \cap \left]a,+\infty\right[ \cap A$, donde $|f \restriction_{A \cap \left]a,+\infty\right[}(x) - L| < \epsilon$, isto é, $|f(x) - L| < \epsilon$. Logo, $\ds \lim_{x \to a^+} f(x) = L$.
                \end{itemize}
            Logo, $\ds \lim_{x \to a^+} f(x) = L \Leftrightarrow \lim_{x \to a} f \restriction_{A \cap \left]a,+\infty \right[ } (x) = L$. \blackproof
            \item Seja $\epsilon \in \R_{>0}$. Como $a \in A'_-$, temos $a \in (A \cap \left]-\infty,a\right[)'$.
                \begin{itemize}
                    \item[($\Rightarrow$)] Como $\ds \lim_{x \to a^-} f(x) = L$, existe $\delta \in \R_{>0}$ tal que, para todo $x \in A$, se $a-\delta < x < a$, então $|f(x) - L| < \epsilon$. Seja $x \in A \cap \left]-\infty,a\right[$ tal que $0 < |x-a| < \delta$. Como $x \in \left]-\infty,a\right[$, temos $x < a$, donde $a-\delta < x < a$. Como $x \in A$, segue que $|f(x) - L| < \epsilon$, isto é, $|f \restriction_{A \cap \left]-\infty,a\right[}(x) - L| < \epsilon$. Logo, $\ds \lim_{x \to a} f \restriction_{A \cap \left]-\infty,a\right[}(x) = L$.
                    \item[($\Leftarrow$)] Como $\ds \lim_{x \to a} f \restriction_{A \cap \left]-\infty,a\right[}(x) = L$, existe $\delta \in \R_{>0}$ tal que, para todo $x \in A \cap \left]-\infty,a\right[$, se $0 < |x-a| < \delta$, então $|f \restriction_{A \cap \left]-\infty,a\right[}(x) - L| < \epsilon$. Seja $x \in A$ tal que $a-\delta < x < a$. Então $x \in \left]-\infty,a\right[$ e $x \in V^*_{\delta}(a) \cap \left]-\infty,a\right[ \cap A$, donde $|f \restriction_{A \cap \left]-\infty,a\right[}(x) - L| < \epsilon$, isto é, $|f(x) - L| < \epsilon$. Logo, $\ds \lim_{x \to a^-} f(x) = L$.
                \end{itemize}
            Logo, $\ds \lim_{x \to a^-} f(x) = L \Leftrightarrow \lim_{x \to a} f \restriction_{A \cap \left]-\infty, a \right[ } (x) = L$. \blackproof
        \end{enumerate}
\end{proof}

\part{Álgebra Linear}
%!TEX root = main.tex

\chapter{Matrizes e Sistemas Lineares}

\begin{definition}
    Sejam $m,n \in \N$.
        \leavevmode
            \begin{enumerate}[leftmargin=*, align=left, label=\textbf{(\alph*)}]
                \item Uma \emph{matriz de tamanho $m \times n$} é uma função $A : [m] \times [n] \to \mathbb{R}$.
                \item Uma matriz $A : [m] \times [n] \to \mathbb{R}$ é denotada por 
                    \begin{enumerate}[label=\roman*.]
                        \item 
                            \begin{itemize}
                                \item $A = (a_{ij})_{m \times n}$, onde $a_{ij} := A(i,j)$, para quaisquer $i \in [m]$ e $j \in [n]$.
                                \item $A = ([A]_{ij})_{m \times n}$, onde $[A]_{ij} := A(i,j)$, para quaisquer $i \in [m]$ e $j \in [n]$.
                            \end{itemize}
                         Dizemos que $a_{ij}$, ou $[A]_{ij}$, é a \emph{entrada} de posição $i, j$.
                        \item uma tabela com $m$ linhas e $n$ colunas
                            \[
                                A =
                                    \begin{bmatrix}
                                        a_{11} & a_{12} & \cdots & a_{1n} \\
                                        a_{21} & a_{22} & \cdots & a_{2n} \\
                                        \vdots & \vdots & \ddots &  \vdots \\
                                        a_{m1} & a_{m2} & \cdots & a_{mn}
                                    \end{bmatrix}_{m \times n}.
                            \]
                    \end{enumerate}
            \end{enumerate}
\end{definition}

\begin{definition}
    A \emph{adição de matrizes} é a operação
        \[
            + : \R^{[m] \times [n]} \times \R^{[m] \times [n]} \to \R^{[m] \times [n]}
        \]
    tal que, para quaisquer $A, B \in \R^{[m] \times [n]}$,
        \[
            [A+B]_{ij} := [A]_{ij} + [B]_{ij},
        \]
    para quaisquer $i \in [m]$ e $j \in [n]$.
\end{definition}

\begin{proposition}
    Sejam $A,B,C \in \R^{[m] \times [n]}$.
        \begin{enumerate}[leftmargin=*, align=left, label=\textbf{(\alph*)}]
            \item $A + (B + C) = (A + B) + C$.
            \item $A+B=B+A$.
            \item $A + 0 = A$.
            \item $A + (-A) = 0$.
        \end{enumerate}
\end{proposition}

\begin{proof}
    \leavevmode
        \begin{enumerate}[leftmargin=*, align=left, label=\textbf{(\alph*)}]
            \item
            \item
            \item
            \item 
        \end{enumerate}
\end{proof}

\begin{definition}
    A \emph{multiplicação de uma matriz por um escalar} é a operação
        \[
            \cdot : \R \times \R^{[m] \times [n]} \to \R^{[m] \times [n]}
        \]
    tal que, para quaisquer $\lambda \in \R$ e $A \in \R^{[m] \times [n]}$,
        \[
            [\lambda \cdot A]_{ij} := \lambda \cdot [A]_{ij},
        \]
    para quaisquer $i \in [m]$ e $j \in [n]$.
\end{definition}

\begin{proposition}
    Sejam $A,B \in \R^{[m] \times [n]}$ e $\lambda_1, \lambda_2 \in \R$.
        \begin{enumerate}[leftmargin=*, align=left, label=\textbf{(\alph*)}]
            \item $\lambda_1 \cdot (\lambda_2 \cdot A) = (\lambda_1 \cdot \lambda_2) \cdot A$.
            \item $(\lambda_1 + \lambda_2) \cdot A = \lambda_1 \cdot A + \lambda_2 \cdot A$.
            \item $\lambda_1 \cdot (A + B) = \lambda_1 \cdot A + \lambda_1 \cdot B$.
            \item $1 \cdot A = A$.
        \end{enumerate}
\end{proposition}

\begin{proof}
    \leavevmode
        \begin{enumerate}[leftmargin=*, align=left, label=\textbf{(\alph*)}]
            \item
            \item
            \item
            \item 
        \end{enumerate}
\end{proof}

\begin{definition}
    A \emph{multiplicação de matrizes} é a operação
        \[
            \cdot : \R^{[m] \times [n]} \times \R^{[n] \times [p]} \to \R^{[m] \times [p]}
        \]
    tal que, para quaisquer $A \in \R^{[m] \times [n]}$ e $B \in \R^{[n] \times [p]}$,
        \[
            [A \cdot B]_{ij} := \sum_{k=1}^{n} [A]_{ik} \cdot [B]_{kj},
        \]
    para quaisquer $i \in [m]$ e $j \in [p]$.
\end{definition}

\begin{proposition}
    Sejam $A = (a_{ij})_{m \times n}$, $B = (b_{jk})_{n \times p}$ e $C = (c_{ik})_{m \times p}$. Se $C = A \cdot B$, então
        \[
            [C]_{ik} = \sum_{j=1}^{n} [A]_{ij} [B]_{jk},
        \]
    para quaisquer $i \in [m]$ e $k \in [p]$.
\end{proposition}

\begin{proof}
    \blackproof
\end{proof}

\begin{proposition}
    Sejam $A$, $B$ e $C$ matrizes de tamanhos compatíveis e $\lambda \in \R$.
        \begin{enumerate}[leftmargin=*, align=left, label=\textbf{(\alph*)}]
            \item $A \cdot (B \cdot C) = (A \cdot B) \cdot C$.
            \item $A \cdot \mathrm{I} = \mathrm{I} \cdot A = A$.
            \item 
                \begin{enumerate}[label=\roman*.]
                    \item $A \cdot (B+C) = A \cdot B + A \cdot C$
                    \item $(A+B) \cdot C = A \cdot C + B \cdot C$
                \end{enumerate}
            \item $(\lambda \cdot A) \cdot B = \lambda \cdot (A \cdot B) = A \cdot (\lambda \cdot B)$.
        \end{enumerate}
\end{proposition}

\begin{proof}
    \leavevmode
        \begin{enumerate}[leftmargin=*, align=left, label=\textbf{(\alph*)}]
            \item
            \item
            \item
            \item
        \end{enumerate}
\end{proof}

\begin{definition}
    Seja $A \in \R^{[m] \times [n]}$.
        \begin{enumerate}[leftmargin=*, align=left, label=\textbf{(\alph*)}]
            \item Dizemos que $A$ é \emph{invertível à esquerda} se existe $B \in \R^{[n] \times [m]}$ tal que $BA=\mathrm{I}_n$. Dizemos, nesse caso, que $B$ é uma \emph{inversa à esquerda} de $A$.
            \item Dizemos que $A$ é \emph{invertível à direita} se existe $B \in \R^{[n] \times [m]}$ tal que $AB = \mathrm{I}_m$. Dizemos, nesse caso, que $B$ é uma \emph{inversa à direita} de $A$.
            \item Dizemos que $A$ é \emph{invertível}, ou \emph{não singular}, se é invertível à esquerda e à direita.
            \item Dizemos que $A$ é \emph{singular} se não é invertível. 
        \end{enumerate}
\end{definition}

\begin{proposition}
    A matriz inversa de uma matriz invertível é única.
\end{proposition}

\begin{proof} 
    Seja $A \in \R^{[m] \times [n]}$ invertível. Então, por definição, existem $B_1, B_2 \in \R^{[n] \times [m]}$ tais que $B_1 A = \mathrm{I}_n$ e $A B_2 = \mathrm{I}_{m}$. Com isso,
        \[
            B_1 = B_1 \mathrm{I}_{m} = B_1 (A B_2) = (B_1 A) B_2 = \mathrm{I}_{n} B_2 = B_2,
        \]
    isto é, $B_1 = B_2$. \blackproof
\end{proof}

\section{Operações e Matrizes Elementares}

Dada uma matriz, podemos 
    \begin{itemize}
        \item trocar a posição de duas de suas linhas;
        \item multiplicar uma de suas linhas por um escalar; e
        \item somar a uma de suas linhas uma outra linha que foi multiplicada por um escalar.
    \end{itemize}
Estas são as chamadas \emph{operações elementares}. Elas serão formalizadas nas definições a seguir e serão úteis no nosso estudo dos sistemas de equações lineares.

\begin{definition}
    Seja $A \in \R^{[m] \times [n]}$.
        \begin{enumerate}[leftmargin=*, align=left, label=\textbf{(\alph*)}]
            \item Seja $i \in [m]$. A \emph{$i$-ésima linha} de $A$ é a matriz $L_i(A) \in \R^{[1] \times [n]}$ definida por $[L_i(A)]_{1j} := a_{ij}$ para todo $j \in [n]$. Isso é denotado por
                \[
                    L_i (A):=
                        \begin{bmatrix}
                            a_{i1} & a_{i2} & \cdots & a_{in}
                        \end{bmatrix}.
                \]
            \item Seja $j \in [n]$. A \emph{$j$-ésima coluna} de $A$ é a matriz $C_j(A) \in \R^{[m] \times [1]}$ definida por $[C_j(A)]_{i1} := a_{ij}$ para todo $i \in [m]$. Isso é denotado por
                \[
                    C_j(A) :=
                        \begin{bmatrix}
                            a_{1j} \\
                            a_{2j} \\
                            \vdots \\
                            a_{mj} \\
                        \end{bmatrix}.
                \]
        \end{enumerate}
\end{definition}

\begin{definition}
    Sejam $A \in \R^{[m] \times [n]}$ e $p, q \in [m]$, com $p \neq q$.
        \begin{enumerate}[leftmargin=*, align=left, label=\textbf{(\alph*)}]
            \item Definimos $A_{L_p \leftrightarrow L_q} \in \R^{[m] \times [n]}$ por
                \[
                    L_i \left( A_{L_p \leftrightarrow L_q} \right) :=
                        \begin{cases}
                            L_q (A) & \text{se } i = p \\
                            L_p (A) & \text{se } i = q \\
                            L_i (A) & \text{se } i \neq p, q
                        \end{cases},
                \]
            isto é,
                \[
                    [A_{L_p \leftrightarrow L_q}]_{ij} :=
                        \begin{cases}
                            [A]_{qj} & \text{se } i=p \\
                            [A]_{pj} & \text{se } i=q \\
                            [A]_{ij} & \text{se } i \neq p, q
                        \end{cases}.
                \]
            \item Seja $\lambda \in \R_{\neq 0}$. Definimos $A_{L_p \leftarrow \lambda L_p} \in \R^{[m] \times [n]}$ por
                \[
                    L_i \left( A_{L_p \leftarrow \lambda L_p} \right) := 
                        \begin{cases}
                            \lambda \cdot L_p (A), & \text{se } i = p \\
                            L_i (A), & \text{se } i \neq p
                        \end{cases}.
                \]
            isto é,
                \[
                    [A_{L_p \leftarrow \lambda L_p}]_{ij} :=
                        \begin{cases}
                            \lambda \cdot [A]_{pj} & \text{se } i=p \\
                            [A]_{ij} & \text{se } i \neq p
                        \end{cases}.
                \]
            \item Seja $\lambda \in \R_{\neq 0}$. Definimos $A_{L_p \leftarrow L_p + \lambda L_q} \in \R^{[m] \times [n]}$ por
                \[
                    L_i \left( A_{L_p \leftarrow L_p + \lambda L_q} \right) := 
                        \begin{cases}
                            L_p (A) + \lambda \cdot L_q (A), & \text{se } i = p \\
                            L_i (A), & \text{se } i \neq p
                        \end{cases}.
                \]
            isto é,
                \[
                    [A_{L_p \leftarrow L_p + \lambda L_q}]_{ij} :=
                        \begin{cases}
                            [A]_{pj} + \lambda \cdot [A]_{qj} & \text{se } i=p \\
                            [A]_{ij} & \text{se } i \neq p
                        \end{cases}.
                \]
        \end{enumerate}
\end{definition}

Intuitivamente, uma matriz elementar é uma matriz que foi obtida da matriz identidade pela aplicação de exatamente uma operação elementar.

\begin{definition}
    Uma matriz $E \in \R^{[n] \times [n]}$ é dita \emph{elementar} se existem $p, q \in [n]$, com $p \neq q$, tais que, exclusivamente,
        \begin{enumerate}[label=\roman*.]
            \item ou $E = \mathrm{I}_{L_p \leftrightarrow L_q}$,
            \item ou $E =  \mathrm{I}_{L_p \leftarrow \lambda L_p}$, para algum $\lambda \in \R_{\neq 0}$,
            \item ou $E =  \mathrm{I}_{L_p \leftarrow L_p + \lambda L_q}$, para algum $\lambda \in \R_{\neq 0}$.
        \end{enumerate}
\end{definition}

\begin{proposition}
    As matrizes elementares são invertíveis.
\end{proposition}

\begin{proof}
    Exibimos, explicitamente, as inversas.
        \begin{itemize}
            \item A inversa de $\mathrm{I}_{L_p \leftrightarrow L_q}$ é ela mesma, isto é, $\mathrm{I}_{L_p \leftrightarrow L_q} \cdot \mathrm{I}_{L_p \leftrightarrow L_q} = \mathrm{I}$.
            \item A inversa de $\mathrm{I}_{L_p \leftarrow \lambda L_p}$ é $\mathrm{I}_{L_p \leftarrow \frac{1}{\lambda} L_p}$, isto é,
                \[
                    \mathrm{I}_{L_p \leftarrow \lambda L_p} \cdot \mathrm{I}_{L_p \leftarrow \frac{1}{\lambda} L_p} = \mathrm{I}_{L_p \leftarrow \frac{1}{\lambda} L_p} \cdot \mathrm{I}_{L_p \leftarrow \lambda L_p} = \mathrm{I}.
                \]
            \item A inversa de $\mathrm{I}_{L_p \leftarrow L_p + \lambda L_q}$ é $\mathrm{I}_{L_p \leftarrow L_p -\lambda L_q}$, isto é,
                \[
                    \mathrm{I}_{L_p \leftarrow L_p + \lambda L_q} \cdot \mathrm{I}_{L_p \leftarrow L_p - \lambda L_q} = \mathrm{I}_{L_p \leftarrow L_p - \lambda L_q} \cdot \mathrm{I}_{L_p \leftarrow L_p + \lambda L_q} = \mathrm{I}.
                \]
        \end{itemize}
    Logo, as matrizes elementares são invertíveis. \blackproof
\end{proof}

Na proposição a seguir, mostramos que fazer uma operação elementar numa matriz $A$ equivale à multiplicar a matriz $A$ por uma matriz elementar pela esquerda.

\begin{proposition}
    Sejam $A \in \R^{[m] \times [n]}$ e $p, q \in [m]$, com $p \neq q$, e $\lambda \in \R$.
        \begin{enumerate}[leftmargin=*, align=left, label=\textbf{(\alph*)}]
            \item Temos $\ds A_{L_p \leftrightarrow L_q} = \mathrm{I}_{L_p \leftrightarrow L_q} \cdot A$.
            \item Temos $\ds A_{L_p \leftarrow \lambda L_p} = \mathrm{I}_{L_p \leftarrow \lambda L_p} \cdot A$.
            \item Temos $\ds A_{L_p \leftarrow L_p + \lambda L_q} = \mathrm{I}_{L_p \leftarrow L_p + \lambda L_q} \cdot A$.
        \end{enumerate}
\end{proposition}

\begin{proof}
    \leavevmode
        \begin{enumerate}[leftmargin=*, align=left, label=\textbf{(\alph*)}]
            \item
            \item
            \item
        \end{enumerate}
\end{proof}

\begin{definition}
    Sejam $A, B \in \R^{[m] \times [n]}$. Dizemos que $A$ e $B$ são \emph{equivalentes por linhas} se existe uma sequência finita de matrizes elementares $E_1, \ldots, E_k \in \R^{[m] \times [m]}$ tais que 
        \[
            E_k \cdot E_{k-1} \cdot \ldots \cdot E_2 \cdot E_1 \cdot A = B.
        \]
    Isso é denotado por $A \sim B$.
\end{definition}

\begin{proposition}
    A equivalência por linhas é uma relação de equivalência.
\end{proposition}

\begin{proof}
    \blackproof
\end{proof}

\section{Eliminação Gaussiana e Decomposição LU}

Intuitivamente, uma matriz é dita \emph{escalonada} se o primeiro elemento não nulo de cada linha não nula está à esquerda do primeiro elemento não nulo de cada uma das linhas não nulas seguintes e se as linhas nulas estão abaixo das demais. Em uma matriz escalonada, o primeiro elemento não nulo de cada linha não nula é chamado de \emph{pivô}, enquanto o número de linhas não nulas é chamado de \emph{posto}.

\begin{definition}
    Uma matriz $A \in \R^{[m] \times [n]}$ é dita \emph{escalonada} se existem $r \in [m]_0$ e $b : [r] \to [n]$ estritamente crescente tais que:
        \begin{enumerate}[label=\roman*.]
            \item $a_{ib_i} \neq 0$ para todo $i \in [r]$;
            \item $a_{ij} = 0$ para quaisquer $i \in [r]$ e $j \in [n]$ tais que $j < b_i$; e
            \item $a_{ij} = 0$ para quaisquer $i \in [m] \setminus [r]$ e $j \in [n]$.
        \end{enumerate}
    %Dizemos que $A$ é uma matriz \emph{escalonada} se existe uma sequência de índices $1 \leq b_1 < b_2 < \ldots < b_r \leq n$ tal que $a_{ib_i} \neq 0$ para todo $i \in [r]$ e $a_{ij} = 0$ para todo $1 \leq j < b_i$. Os termos $a_{ib_i}$ são chamados de \emph{pivôs}, enquanto o número de pivôs, $r$, é chamado de \emph{posto}.
\end{definition}

\begin{theorem}
    Toda matriz é equivalente por linhas a uma matriz escalonada.
\end{theorem}

\begin{proof}
    \blackproof
\end{proof}

Intuitivamente, uma matriz escalonada é dita \emph{reduzida} se todo pivô é unitário e se todos os outros elementos da coluna de um pivô são iguais a 0.

\begin{definition}
    Uma matriz $A \in \R^{[m] \times [n]}$ é dita \emph{escalonada reduzida} se existem $r \in [m]_0 $ e $b : [r] \to [n]$ estritamente crescente tais que:
        \begin{enumerate}[label=\roman*.]
            \item $a_{ib_i} = 1$ para todo $i \in [r]$;
            \item $a_{k b_i} = 0$ para quaisquer $i \in [r]$ e $k \in [m] \setminus \{i\}$;
            \item $a_{ij} = 0$ para quaisquer $i \in [r]$ e $j \in [n]$ tais que $j < b_i$; e
            \item $a_{ij} = 0$ para quaisquer $i \in [m] \setminus [r]$ e $j \in [n]$.
        \end{enumerate}
\end{definition}

\begin{theorem}
    Toda matriz é equivalente por linhas a uma única matriz escalonada reduzida.
\end{theorem}

\begin{proof}
    \blackproof
\end{proof}

\begin{definition}
    Uma matriz $P \in \R^{[m] \times [m]}$ é dita uma \emph{matriz de permutação} se $P = \mathrm{I}_m$ ou se existem $k \in \N$ e matrizes elementares $E_1, \ldots, E_k \in \R^{[m] \times [m]}$ da forma $\mathrm{I}_{L_p \leftrightarrow L_q}$ tais que
        \[
            P = E_k \cdot E_{k-1} \cdot \ldots \cdot E_2 \cdot E_1.
        \]
\end{definition}

\begin{definition}
    Uma matriz $A \in \R^{[m] \times [n]}$ é dita
        \begin{enumerate}[label=\roman*.]
            \item \emph{triangular superior} se $a_{ij} = 0$ para quaisquer $i \in [m]$ e $j \in [n]$ com $i > j$.
            \item \emph{triangular inferior} se $a_{ij} = 0$ para quaisquer $i \in [m]$ e $j \in [n]$ com $i < j$.
            \item \emph{triangular inferior unitária} se é triangular inferior e $a_{ii} = 1$ para todo $i \in [\min\{m, n\}]$.
        \end{enumerate}
\end{definition}

\begin{theorem}
    Seja $A \in \R^{[m] \times [n]}$. Existem uma matriz de permutação $P \in \R^{[m] \times [m]}$, uma matriz triangular inferior unitária $L \in \R^{[m] \times [m]}$ e uma matriz triangular superior $U \in \R^{[m] \times [n]}$ tais que
        \[
            P \cdot A = L \cdot U.
        \]
\end{theorem}

\begin{proof}
    \blackproof
\end{proof}

\section{Sistemas Lineares}

Sejam $a_1, a_2, \ldots , a_n \in \R$ e $b \in \R$. Existem $x_1, x_2, \ldots, x_n \in \R$ tais que
    \[
        a_1 x_1 + a_2 x_2 + \cdots + a_n x_n = b \text{?}
    \]
Toda equação desse tipo é chamada de \emph{equação linear} nas \emph{incógnitas} $x_1, x_2, \ldots, x_n$ com \emph{coeficientes} $a_1, a_2, \ldots, a_n$ e \emph{termo independente} $b$. \emph{Resolver} uma equação linear consiste em caracterizar explicitamente o conjunto
    \[
        \{ (x_1, \ldots, x_n) \in \R^n : a_1 x_1 + a_2 x_2 + \cdots + a_n x_n = b \}.
    \]
Como formalizar essas ideias? Fazemos isso usando matrizes.


\begin{definition}
    Seja $n \in \N$.
        \begin{enumerate}[leftmargin=*, align=left, label=\textbf{(\alph*)}]
            \item Uma \emph{equação linear} é um par $(a,b)$, onde $a \in \R^{[1] \times [n]}$ e $b \in \R^{[1] \times [1]}$.
            \item Dizemos que $x \in \R^{[n] \times [1]}$ é uma \emph{solução} de uma equação linear $(a,b) \in \R^{[1] \times [n]} \times \R^{[1] \times [1]}$ se $ax=b$.
            \item O conjunto de todas as soluções de uma equação linear $(a,b) \in \R^{[1] \times [n]} \times \R^{[1] \times [1]}$ é denotado por $S(a,b)$, isto é,
                \[
                    S(a,b) := \{x \in \R^{[n] \times [1]} : ax=b\}.
                \]
        \end{enumerate}
\end{definition}

Intuitivamente, um \emph{sistema de equações lineares} é um conjunto finito de equações lineares. Com $m$ equações e $n$ incógnitas, a representação canônica é
    \[
        \left\{
            \begin{matrix}
                a_{11} x_1 + a_{12} x_2 + \cdots + a_{1n} x_n = b_1 \\
                a_{21} x_1 + a_{22} x_2 + \cdots + a_{2n} x_n = b_2 \\
                \vdots \\
                a_{m1} x_1 + a_{m2} x_2 + \cdots + a_{mn} x_n = b_m
            \end{matrix}
        \right. .
    \]

\emph{Resolver} um sistema de equações lineares consiste em caracterizar explicitamente o conjunto
    \[
        \bigcap_{i=1}^{m} \{ (x_1, \ldots, x_n) \in \R^n : a_{i1} x_1 + \cdots + a_{in} x_n = b_i \}.
    \]
Como formalizar essas ideias? Fazemos isso usando matrizes.

\begin{definition}
    Sejam $m,n \in \N$.
        \begin{enumerate}[leftmargin=*, align=left, label=\textbf{(\alph*)}]
            \item Um \emph{sistema de equações lineares} é um par $(A,b)$, onde $A \in \R^{[m] \times [n]}$ e $b \in \R^{[m] \times [1]}$.
            \item Dizemos que $x \in \R^{[n] \times [1]}$ é uma \emph{solução} de um sistema de equações lineares $(A,b) \in \R^{[m] \times [n]} \times \R^{[m] \times [1]}$ se $Ax=b$.
            \item O conjunto de todas as soluções de um sistema de equações lineares $(A,b) \in \R^{[m] \times [n]} \times \R^{[m] \times [1]}$ é denotado por $S(A,b)$, isto é,
                \[
                    S(A,b) := \{x \in \R^{[n] \times [1]} : Ax=b\}.
                \]
        \end{enumerate}
\end{definition}

Por simplicidade, vamos nos referir a um sistema de equações lineares como um \emph{sistema linear}, ou ainda, como um \emph{sistema}.

\begin{proposition}
    Seja $(A,b) \in \R^{[m] \times [n]} \times \R^{[m] \times [1]}$ um sistema linear. Então
        \[
            S(A,b) = \bigcap_{i \in [m]} S(L_i(A),L_i(b)).
        \]
\end{proposition}

\begin{proof}
    \blackproof
\end{proof}

\begin{definition}
    Seja $(A,b) \in \R^{[m] \times [n]} \times \R^{[m] \times [1]}$ um sistema linear.
        \begin{enumerate}[leftmargin=*, align=left, label=\textbf{(\alph*)}]
            \item Dizemos que $A$ é a \emph{matriz de coeficientes} do sistema $(A,b)$.
            \item Dizemos que $[A|b] \in \R^{[m] \times [n+1]}$, definida por
                \[
                    [A|b]_{ij} :=
                        \begin{cases}
                            [A]_{ij} & \text{se } i \in [m] \text{ e } j \in [n] \\
                            [b]_{i1} & \text{se } i \in [m] \text{ e } j = n+1
                        \end{cases},
                \]
            isto é,
                \[
                    [A | b] :=
                        \begin{bmatrix}
                            a_{11} & a_{12} & \cdots & a_{1n} & b_{11}\\
                            a_{21} & a_{22} & \cdots & a_{2n} & b_{21} \\
                            \vdots & \vdots & \ddots &  \vdots & \vdots \\
                            a_{m1} & a_{m2} & \cdots & a_{mn} & b_{m1}
                        \end{bmatrix}_{m \times (n+1)},
                \]
            é a \emph{matriz aumentada} do sistema $(A,b)$.
        \end{enumerate}
\end{definition}

\begin{definition}
    Seja $(A,b)$ um sistema linear.
        \begin{enumerate}[leftmargin=*, align=left, label=\textbf{(\alph*)}]
            \item Dizemos que $(A,b)$ é \emph{possível} se $S(A,b) \neq \emptyset$.
            \item Dizemos que $(A,b)$ é \emph{impossível} se $S(A,b) = \emptyset$.
        \end{enumerate}
\end{definition}

\begin{theorem}
    Seja $(A,b)$ um sistema linear. Se $|S(A,b)| \geq 2$, então $|S(A,b)| = |\R|$.
\end{theorem}

\begin{proof}
    Como $|S(A,b)| \geq 2$, sejam $x_1, x_2 \in S(A,b)$, com $x_1 \neq x_2$. Defina $f : \R \to \R^{[n]\times[1]}$ por $f(\lambda) := \lambda x_1 + (1-\lambda)x_2$ para todo $\lambda \in \R$. Então, para todo $\lambda \in \R$, temos
        \[
            A f(\lambda) = A [\lambda x_1 + (1 - \lambda)x_2] = b,
        \]
    de modo que $f(\lambda) \in S(A,b)$. Com isso, $\Im f \subseteq S(A,b)$. Agora, para quaisquer $\lambda_1, \lambda_2 \in \R$, se $f(\lambda_1) = f(\lambda_2)$, então $(\lambda_1 - \lambda_2)(x_1 - x_2) = 0$, e como $x_1 \neq x_2$, temos $\lambda_1 = \lambda_2$. Logo, $f$ é injetiva, de modo que $|\R| \leq |S(A,b)|$. Agora, como $S(A,b) \subseteq \R^{[n] \times [1]}$ e $|\R^{[n]\times[1]}| = |\R|$, temos $|S(A,b)| \leq |\R|$. Logo, pelo teorema de Cantor-Schroeder-Berstein, temos $|S(A,b)| = |\R|$.
\end{proof}

\begin{definition}
    Seja $(A,b)$ um sistema linear possível.
        \begin{enumerate}[leftmargin=*, align=left, label=\textbf{(\alph*)}]
            \item Dizemos que $(A,b)$ é \emph{determinado} se $|S(A,b)| = 1$.
            \item Dizemos que $(A,b)$ é \emph{indeterminado} se $|S(A,b)| \geq 2$.
        \end{enumerate}
\end{definition}

\begin{definition}
    Dois sistemas lineares $(A,b) \in \R^{[m] \times [n]} \times \R^{[m] \times [1]}$ e $(C,d) \in \R^{[p] \times [n]} \times \R^{[p] \times [1]}$ são ditos \emph{equivalentes} se $S(A,b) = S(C,d)$.
\end{definition}

\begin{proposition}
    Sejam $(A,b),(C,d) \in \R^{[m]\times[n]}\times\R^{[m]\times[1]}$ sistemas lineares.
        \begin{enumerate}[leftmargin=*, align=left, label=\textbf{(\alph*)}]
            \item Se as matrizes aumentadas $[A|b]$ e $[C|d]$ são equivalentes por linhas, então os sistemas $(A,b)$ e $(C,d)$ são equivalentes.
            \item Se $(A,b)$ e $(C,d)$ são possíveis, então $(A,b)$ e $(C,d)$ são equivalentes se, e somente se, as matrizes aumentadas $[A|b]$ e $[C|d]$ são equivalentes por linhas.
        \end{enumerate}
\end{proposition}

\begin{proof}
    Reginaldo, 32.
\end{proof}

\begin{definition}
    Um sistema linear $(A,b)$ é \emph{homogêneo} se $b=0$.
\end{definition}

\begin{proposition}
    Todo sistema linear homogêneo $(A,0)$ é possível.
\end{proposition}

\begin{proof}
    Basta ver que $x = 0$ é uma solução, dita \emph{trivial}. \blackproof
\end{proof}

\begin{theorem} \label{teo.alglin:sislinhom}
    Seja $(A,0) \in \R^{[m] \times [n]} \times \R^{[m] \times [1]}$ um sistema linear homogêneo.
        \begin{enumerate}[leftmargin=*, align=left, label=\textbf{(\alph*)}]
            \item Se $m < n$, então $(A,0)$ tem uma solução não trivial, isto é, existe $x \in \R^{[n] \times [1]} \setminus \{ 0\}$ tal que $Ax=0$.
            \item Se $m < n$, então $(A,0)$ tem infinitas soluções.
        \end{enumerate}
    %Todo sistema linear homogêneo em que o número de incógnitas é maior que o número de equações possui uma solução não trivial (e, portanto, possui infinitas soluções).
\end{theorem}

\begin{proof}
    Elon, 27. \blackproof
\end{proof}

\chapter{Espaços Vetoriais}

\section{Espaços e Subespaços Vetoriais}

\begin{definition}
    Um \emph{espaço vetorial real} é uma tripla $(V,+,\cdot)$, onde $+ : V \times V \to V$ e $\cdot : \R \times V \to V$ são operações tais que
        \begin{itemize}
            \item A1: $u + (v + w) = (u + v) + w$ para quaisquer $u,v,w \in V$;
            \item A2: $u + v = v + u$ para quaisquer $u,v \in V$;
            \item A3: existe $0_V \in V$ tal que $u + 0_V = u$ para todo $u \in V$;
            \item A4: para cada $u \in V$ existe $v \in V$ tal que $u+v=0_V$;
            \item M1: $\lambda_1 \cdot (\lambda_2 \cdot u) = (\lambda_1 \cdot \lambda_2) \cdot u$ para quaisquer $u \in V$ e $\lambda_1,\lambda_2 \in \R$;
            \item M2: $(\lambda_1 + \lambda_2) \cdot u = \lambda_1 \cdot u + \lambda_2 \cdot u$ para quaisquer $u \in V$ e $\lambda_1,\lambda_2 \in \R$;
            \item M3: $\lambda_1 \cdot (u + v) = \lambda_1 \cdot u + \lambda_1 \cdot v$ para quaisquer $u,v \in V$ e $\lambda_1 \in \R$;
            \item M4: $1 \cdot u = u$ para todo $u \in V$.
        \end{itemize}
    Os elementos de $V$ são chamados de \emph{vetores}.
\end{definition}

\begin{remark}
    Quando não houver perigo de confunsão, podemos nos referir ao espaço vetorial $(V, +, \cdot)$ simplesmente como o conjunto $V$.
\end{remark}

\begin{example}
    \leavevmode
        \begin{enumerate}[leftmargin=*, align=left, label=\textbf{(\alph*)}]
            \item O espaço $\R^n$, com soma e produto por escalar usuais, é um espaço vetorial real.
            \item Seja $X \subseteq \R$. O conjunto $\R^X$, com soma e produto por escalar usuais, é um espaço vetorial real.
            \item O conjunto $\R_{>0}$, com as operações $x\oplus y:=x \cdot y$ e $\alpha \odot x := x^{\alpha}$, é um espaço vetorial real.
        \end{enumerate}
\end{example}

\begin{proposition} \label{prop.alglin:manobras_algebricas}
    Seja $(V, +, \cdot)$ um espaço vetorial.
        \begin{enumerate}[leftmargin=*, align=left, label=\textbf{(\alph*)}]
            \item (Unicidade)
                \begin{enumerate}[label=\roman*.]
                    \item O elemento neutro $0_V$ de $+$ é único.
                    \item O oposto aditivo de cada elemento de $V$ é único.
                \end{enumerate}
            \item (Integridade) Para quaisquer $u \in V$ e $\lambda \in \R$, temos
                \begin{enumerate}[label=\roman*.]
                    \item $\lambda \cdot 0_V = 0_V$;
                    \item $0 \cdot u = 0_V$;
                    \item se $\lambda \cdot u = 0_V$, então $\lambda = 0$ ou $u = 0_V$.
                \end{enumerate}
            \item (Sinais) Para quaisquer $u \in V$ e $\lambda \in \R$, temos
                \begin{enumerate}[label=\roman*.]
                    \item $(-1) \cdot u = -u$;
                    \item $-(-u) = u$;
                    \item $(- \lambda) \cdot u = \lambda (-u) = -(\lambda \cdot u)$.
                \end{enumerate}
            \item (Lei do Corte) Para quaisquer $u,v,w \in V$ e $\lambda, \lambda_1, \lambda_2 \in \R$, temos
                \begin{enumerate}[label=\roman*.]
                    \item $u + w = v + w \Rightarrow u = v$;
                    \item $\lambda \neq 0 \land \lambda \cdot u = \lambda \cdot v \Rightarrow u = v$; 
                    \item $u \neq 0_V \land \lambda_1 u = \lambda_2 u \Rightarrow \lambda_1 = \lambda_2$.
                \end{enumerate}
            \item Se $V \neq \{0 \}$, então $V$ é infinito.
        \end{enumerate}
\end{proposition}

\begin{proof}
\end{proof}

\begin{definition}
    Sejam $(V, +_V, \cdot_V)$ um espaço vetorial real e $W \subseteq V$. Dizemos que $W$ é um \emph{subespaço vetorial} de $V$ se $(W, +_V \restriction_{W \times W}, \cdot_V \restriction_{\R \times W})$ é um espaço vetorial.
\end{definition}

\begin{proposition} 
    Sejam $(V, +_V, \cdot_V)$ um espaço vetorial real e $W \subseteq V$. As seguintes afirmações são equivalentes. 
        \begin{enumerate}[label=\roman*.]
            \item $(W, +_V \restriction_{W \times W}, \cdot_V \restriction_{\R \times W})$ é um espaço vetorial.
            \item $0_V \in W$ e, para quaisquer $u, v \in W$ e $\lambda \in \R$, $u +_V v \in W$ e $\lambda \cdot_V u \in W$.
        \end{enumerate}
\end{proposition}

\begin{proof}
    \blackproof
\end{proof}

\begin{example}
    Todo espaço vetorial $V$ tem pelo menos dois subespaços vetoriais, dito \emph{triviais}: $\{0_V\}$ e o próprio $V$.
\end{example}

\section{Combinações Lineares e Geradores}

\begin{definition}
    Seja $(V, +, \cdot)$ um espaço vetorial. Dizemos que $u \in V$ é uma \emph{combinação linear} de $u_1, u_2 , \ldots, u_n \in V$ se existem $\lambda_1, \lambda_2, \ldots, \lambda_n \in \R$ tais que $\ds u = \sum_{i=1}^{n} \lambda_i u_i$.
\end{definition}

\section{Dependência e Independência Linear}

\begin{definition}
    Seja $(V, +, \cdot)$ um espaço vetorial.
        \begin{enumerate}[leftmargin=*, align=left, label=\textbf{(\alph*)}]
            \item Dizemos que $u_1, u_2, \ldots, u_n \in V$ são \emph{linearmente independentes} (L.I.) se a única solução da equação $\lambda_1 u_1 + \cdots + \lambda_n u_n = 0$ é $\lambda_i = 0$ para todo $i \in [n]$. Caso contrário, isto é, se existe uma solução tal que $\lambda_i \neq 0$ para algum $i \in [n]$, dizemos que $u_1, u_2, \ldots, u_n \in V$ são \emph{linearmente dependentes} (L.D.).
            \item Seja $S \subseteq V$ finito.\footnote{Se $S$ é finito, então $S = V$ somente se $V = \{ 0\}$.}
                \begin{enumerate}[label=\roman*.]
                    \item Dizemos que $S$ é L.I. se os vetores de $S$ são L.I..
                    \item Dizemos que $S$ é L.D. se os vetores de $S$ são L.D..
                \end{enumerate}
            \item Seja $S \subseteq V$ infinito.
                \begin{enumerate}[label=\roman*.]
                    \item Dizemos que $S$ é L.I. se todo subconjunto finito de $S$ é L.I..
                    \item Dizemos que $S$ é L.D. se pelo menos um subconjunto finito de $S$ é L.D..
                \end{enumerate}
        \end{enumerate}
\end{definition}

\section{Base e Dimensão}

\begin{lemma} \label{lema.alglin:steinitz}
    Seja $(V,+,\cdot)$ um espaço vetorial. Seja $S \subseteq V$ finito e L.I.. Para todo $T \subseteq [S]$, se $|T| = |S|+1$, então $T$ é L.D..
\end{lemma}

\begin{proof}
    \blackproof
\end{proof}

\begin{definition}
    Seja $(V, +, \cdot)$ um espaço vetorial. Dizemos que $\mathcal{B} \subsetneq V$ é uma \emph{base} de $V$ se $\mathcal{B}$ é L.I. e $[\mathcal{B}] = V$.
\end{definition}

\begin{theorem}[Completamento]
     Todo espaço vetorial finitamente gerado possui uma base. 
\end{theorem}

\begin{theorem}[Invariância] \label{teo.alglin:invariancia}
    Seja $(V, +, \cdot)$ um espaço vetorial finitamente gerado. Se $\mathcal{B}_1$ e $\mathcal{B}_2$ são bases de $V$, então $|\mathcal{B}_1| = |\mathcal{B}_2|$.
\end{theorem}

\begin{definition}
    A \emph{dimensão} de um espaço vetorial finitamente gerado $(V, +, \cdot)$ é o número de vetores de qualquer uma de suas bases. Mais especificamente, se $\mathcal{B}$ é uma base de $V$, a dimensão de $V$ é definida como $\dim V := |\mathcal{B}|$.
\end{definition}

\section{Produto Interno}

\chapter{Transformações Lineares}

\begin{definition}
    Sejam $(V, +_V, \cdot_V)$ e $(W, +_W, \cdot_W)$ espaços vetoriais.
        \begin{enumerate}[leftmargin=*, align=left, label=\textbf{(\alph*)}]
            \item Uma \emph{transformação linear} é uma função $T : V \to W$ tal que $T(u +_V v) = T(u) +_W T(v)$ e $T(\lambda \cdot_V u) = \lambda \cdot_W T(u)$ para quaisquer $u, v \in V$ e $\lambda \in \R$.
            \item Um \emph{operador linear} é uma transformação linear $T : V \to V$.
        \end{enumerate}
\end{definition}

\begin{example}
    Sejam $(V, +, \cdot)$ e $(W, +, \cdot)$ espaços vetoriais.
        \begin{enumerate}[leftmargin=*, align=left, label=\textbf{(\alph*)}]
            \item A função $0 : V \to W$ definida por $0(v) = 0_W$ para todo $v \in V$ é uma transformação linear, chamada de \emph{transformação nula}.
            \item A função $I_V : V \to V$ definida por $I_V(v) = v$ para todo $v \in V$ é um operador linear, chamada de \emph{transformação identidade}. 
        \end{enumerate}
\end{example}

\begin{definition}
    Sejam $(V, +_V, \cdot_V)$ e $(W, +_W, \cdot_W)$ espaços vetoriais. O conjunto de todas as transformações lineares de $V$ em $W$ é denotado por $\mathcal{L}(V,W)$, isto é,
        \[
            \mathcal{L}(V,W) := \{T \in W^V : T \text{ é transformação linear} \}.
        \]
\end{definition}

\begin{proposition}
    Sejam $(V, +_V, \cdot_V)$ e $(W, +_W, \cdot_W)$ espaços vetoriais e $T : V \to W$ uma transformação linear.
        \begin{enumerate}[leftmargin=*, align=left, label=\textbf{(\alph*)}]
            \item $T(0_V) = 0_W$.
        \end{enumerate}
\end{proposition}

\begin{definition}
    Seja $T : V \to W$ uma transformação linear.
        \begin{enumerate}[leftmargin=*, align=left, label=\textbf{(\alph*)}]
            \item O \emph{núcleo} de $T$ é definido como
                \[
                    \ker{T} := \{ v \in V : T(v) = 0_W \}.
                \]
            \item A \emph{imagem} de $T$ é definida como
                \[
                    \Im{T} := \{ w \in W : \exists v(v \in V \land T(v) = w) \}.
                \]
        \end{enumerate}
\end{definition}

\begin{theorem}
    Seja $T : V \to W$ uma transformação linear.
        \begin{enumerate}[leftmargin=*, align=left, label=\textbf{(\alph*)}]
            \item $\ker{T}$ é um subespaço vetorial de $V$.
            \item $\Im{T}$ é um subespaço vetorial de $W$.
        \end{enumerate}
\end{theorem}

\part{Cálculo II}
%!TEX root = main.tex

\chapter{Limites e Continuidade}

\section{Topologia do Espaço Euclidiano}

\begin{definition}
    Sejam $a \in \R^n$ e $r \in \R_{>0}$.
        \begin{enumerate}[leftmargin=*, align=left, label=\textbf{(\alph*)}]
            \item A \emph{bola aberta de centro $a$ e raio $r$} é definida como
                \[
                    B_r(a) := \{x \in \R^n : \|x-a\| < r\}.
                \]
            \item A \emph{bola aberta furada de centro $a$ e raio $r$} é definida como
                \[
                    B^*_r(a) := \{x \in \R^n : 0 < \|x-a\| < r\}.
                \]
        \end{enumerate}
\end{definition}

\begin{corollary}
    Sejam $a \in \R^n$ e $r \in \R_{>0}$. Então $B^*_r(a) = B_r(a) \setminus \{a\}$.
\end{corollary}

\begin{proof}
\end{proof}

\begin{definition}
    Seja $A \subseteq \R^n$.
        \begin{enumerate}[leftmargin=*, align=left, label=\textbf{(\alph*)}]
            \item Dizemos que $a \in A$ é um \emph{ponto interior} de $A$ se existe $r \in \R_{>0}$ tal que $B_r(a) \subseteq A$.
            \item O \emph{interior} de $A$ é definido como
                \[
                    \op{int}{A} := \{ x \in A : \exists r (r \in \R_{>0} \land B_r(x) \subseteq A ) \}.
                \]
        \end{enumerate}
\end{definition}

\begin{definition}
    Dizemos que $A \subseteq \R^n$ é \emph{aberto} se $\op{int}{A} = A$.
\end{definition}

\begin{corollary}
    Toda bola aberta é aberta.
\end{corollary}

\begin{proof}
    \blackproof
\end{proof}

\begin{theorem}
    \leavevmode
        \begin{enumerate}[leftmargin=*, align=left, label=\textbf{(\alph*)}]
            \item $\emptyset$ e $\R^n$ são abertos.
            \item Toda união de abertos do $\R^n$ é aberta.
            \item Toda interseção finita de abertos do $\R^n$ é aberta.
        \end{enumerate}
\end{theorem}

\begin{proof}
\end{proof}

\begin{definition}
    Sejam $a \in \R^n$ e $r \in \R_{>0}$.
        \begin{enumerate}[leftmargin=*, align=left, label=\textbf{(\alph*)}]
            \item A \emph{bola fechada de centro $a$ e raio $r$} é definida como
                \[
                    B_r[a] := \{x \in \R^n : \|x-a\| \leq r\}.
                \]
            \item A \emph{bola fechada furada de centro $a$ e raio $r$} é definida como
                \[
                    B^*_r[a] := \{x \in \R^n : 0 < \|x-a\| \leq r\}.
                \]
        \end{enumerate}
\end{definition}

\begin{definition}
    Dizemos que $A \subseteq \R^n$ é \emph{fechado} se $A^C$ é aberto.
\end{definition}

\begin{corollary}
    Toda bola fechada é fechada.
\end{corollary}

\begin{proof}
    \blackproof
\end{proof}

\begin{theorem}
    \leavevmode
        \begin{enumerate}[leftmargin=*, align=left, label=\textbf{(\alph*)}]
            \item $\emptyset$ e $\R^n$ são fechados.
            \item Toda união finita de fechados do $\R^n$ é fechada.
            \item Toda interseção de fechados do $\R^n$ é fechada.
        \end{enumerate}
\end{theorem}

\begin{proof}
    \blackproof
\end{proof}

\begin{definition}
    Seja $A \subseteq \R^n$.
        \begin{enumerate}[leftmargin=*, align=left, label=\textbf{(\alph*)}]
            \item Dizemos que $a \in \R^n$ é um \emph{ponto exterior} de $A$ se $a \in \op{int}{A^C}$.
            \item O \emph{exterior} de $A$ é definido como $\op{ext}{A} := \op{int}{A^C}$.
        \end{enumerate}
\end{definition}

\begin{definition}
    Seja $A \subseteq \R^n$.
        \begin{enumerate}[leftmargin=*, align=left, label=\textbf{(\alph*)}]
            \item Dizemos que $a \in \R^n$ é um \emph{ponto da fronteira} de $A$ se $a \notin \op{int}{A}$ e $a \notin \op{ext}{A}$.
            \item A \emph{fronteira} de $A$ é definida como
                \[
                    \partial A := \{ x \in \R^n : x \notin \op{int}{A} \land x \notin \op{ext}{A} \}.
                \]
        \end{enumerate}
\end{definition}

\begin{definition}
    Sejam $a \in \R^n$ e $r \in \R_{>0}$. A \emph{esfera de centro $a$ e raio $r$} é definida como
        \[
            S_r(a) := \{x \in \R^n : \|x-a\| = r\}.
        \]
\end{definition}

\begin{proposition}
    Sejam $a \in \R^n$ e $r \in \R_{>0}$.
        \begin{enumerate}[leftmargin=*, align=left, label=\textbf{(\alph*)}]
            \item $B_r[a] = B_r(a) \cup S_r(a)$.
            \item $\partial (B_r(a)) = \partial (B_r[a]) = S_r(a)$.
        \end{enumerate}
\end{proposition}

\begin{proof}
    \blackproof
\end{proof}

\begin{definition}
    Seja $A \subseteq \R^n$.
        \begin{enumerate}[leftmargin=*, align=left, label=\textbf{(\alph*)}]
            \item Dizemos que $a \in \R^n$ é um \emph{ponto de acumulação} de $A$ se
                \[
                    B^*_r(a) \cap A \neq \emptyset
                \]
            para todo $r \in \R_{>0}$.
            \item O conjunto de todos os pontos de acumulação de $A$ é denotado por $A'$, isto é,
                \[
                    A' := \{x \in \R^n : \forall r (r \in \R_{>0} \rightarrow B^*_r(x) \cap A \neq \emptyset) \}.
                \]
        \end{enumerate}
\end{definition}

\begin{definition}
    Dizemos que $a \in A$ é um \emph{ponto isolado} de $A \subseteq \R^n$ se $a \notin A'$.
\end{definition}

\begin{definition}
    Seja $A \subseteq \R^n$.
        \begin{enumerate}[leftmargin=*, align=left, label=\textbf{(\alph*)}]
            \item Dizemos que $a \in \R^n$ é um \emph{ponto de aderência} de $A$ se
                \[
                    B_r(a) \cap A \neq \emptyset
                \]
            para todo $r \in \R_{>0}$.
            \item O \emph{fecho} de $A$ é definido como
                \[
                    \bar{A} := \{x \in \R^n : \forall r (r \in \R_{>0} \rightarrow B_r(x) \cap A \neq \emptyset) \}.
                \]
        \end{enumerate}
\end{definition}

\begin{proposition}
    Seja $A \subseteq \R^n$. 
        \begin{enumerate}[leftmargin=*, align=left, label=\textbf{(\alph*)}]
            \item 
                \begin{enumerate}[label=\roman*.]
                    \item $\op{int}{A}$ é aberto.
                    \item $\bar{A}$ é fechado.
                    \item $\partial A$ é fechado.
                \end{enumerate}
            \item 
                \begin{enumerate}[label=\roman*.]
                    \item $\op{int}{A} \cap \partial A = \emptyset$.
                    \item $\bar{A} = A \cup A'$.
                    \item $\bar{A} = A \cup \partial A$.
                \end{enumerate}
            \item Se $a \in A'$, então $a \in \op{int}{A}$ ou $a \in \partial A$.
            \item $A$ é fechado se, e somente se, $\bar{A} = A$.
        \end{enumerate}
\end{proposition}

\begin{proof}
    \blackproof
\end{proof}

\begin{definition}
    Dizemos que $A \subseteq \R^n$ é \emph{limitado} se existe $r \in \R_{>0}$ tal que $A \subseteq B_r(0)$.
\end{definition}

\begin{definition}
    Dizemos que $A \subseteq \R^n$ é \emph{compacto} se é fechado e limitado.
\end{definition}

\section{Limites}

\begin{definition}
    Dizemos que $L \in \R^m$ é o \emph{limite} de $f : A \subseteq \R^n \to \R^m$ em $a \in A'$ se para todo $\epsilon \in \R_{>0}$ existe $\delta = \delta(\epsilon, a) \in \R_{>0}$ tal que, para todo $x \in A$, se $0 < \| x-a \| < \delta$, então $\| f(x)-L \| < \epsilon$. Isso é denotado por 
        \[
            \lim_{x \to a} f(x) = L.
        \]
\end{definition}

\begin{proposition}[Unicidade] \label{teo.calc2:unicidade_limite_bilateral}
    O limite de $f : A \subseteq \R^n \to \R^m$ em $a \in A'$, quando existe, é único.
\end{proposition}

\begin{proposition}
    Sejam $f : A \subseteq \R^n \to \R^m$, $a \in A'$ e $L \in \R^m$. Então
        \[
            \lim_{x \to a} f(x) = L \Leftrightarrow \lim_{\| x-a \| \to 0} \|f(x) - L \| = 0.
        \]
\end{proposition}

\begin{proof}
    \blackproof
\end{proof}

\section{Continuidade}

\begin{definition}[Continuidade] \label{def.calc2:continuidade}
    Dizemos que $f : A \subseteq \R^n \to \R^m$ é \emph{contínua em $a \in A$} se para todo $\epsilon \in \R_{>0}$ existe $\delta = \delta(\epsilon, a) \in \R_{>0}$ tal que, para todo $x \in A$, se $\| x-a \| < \delta$, então $\| f(x) - f(a) \| < \epsilon$.
\end{definition}

\begin{definition}
    Sejam $f : A \subseteq \R^n \to \R^m$ e $X \subseteq A$.
        \begin{enumerate}[leftmargin=*, align=left, label=\textbf{(\alph*)}]
            \item Dizemos que $f$ é \emph{contínua em $X$} se $f$ é contínua em todos os pontos de $X$.
            \item Dizemos que $f$ é \emph{contínua} se $f$ é contínua em $A$.
        \end{enumerate}
\end{definition}

\begin{theorem} \label{teo.calc2:continuidade_sse_limite}
    Sejam $f : A \subseteq \R^n \to \R^m$ e $a \in A \cap A'$. Então, $f$ é contínua em $a$ se, e somente se, $\ds \lim_{x \to a} f(x) = f(a)$.
\end{theorem}

\begin{proof}
    \blackproof
\end{proof}

\begin{proposition}
    Sejam $f : A \subseteq \R^n \to \R^m$ e $g : B \subseteq \R^m \to \R^k$ tais que $f(A) \subseteq B$. Se $f$ é contínua em $a \in A$ e $g$ é contínua em $f(a) \in B$, então $g \circ f : X \subseteq \R^n \to \R^k$ é contínua em $a$.
\end{proposition}

\begin{proof}
    \blackproof
\end{proof}
\part{Espaços Métricos}
%!TEX root = main.tex

\chapter{Espaços Métricos}

\begin{definition}
    \leavevmode
        \begin{enumerate}[leftmargin=*, align=left, label=\textbf{(\alph*)}]
            \item Uma \emph{métrica} em $X \neq \emptyset$ é uma função $d : X \times X \to \R$ tal que, para quaisquer $x,y,z \in X$,
                \begin{enumerate}
                    \item $d(x,y) \geq 0$;
                    \item $d(x,y) = 0$ se, e somente se, $x=y$.
                    \item $d(x,y) = d(y,x)$; e
                    \item $d(x,z) \leq d(x,y) + d(y,z)$.
                \end{enumerate}
            \item Um \emph{espaço métrico} é um par $(X,d)$, onde $d : X \times X \to \R$ é uma métrica em $X \neq \emptyset$.
        \end{enumerate}
\end{definition}

%%%%%%%%%%%%%%%%%%%%%%%%%%%%%%

\nocite{*}
\printbibliography[heading=bibintoc]

\end{document}